% !TEX program = xelatex
\documentclass[11pt]{article}
\usepackage[a4paper,top=24mm,bottom=24mm,left=20mm,right=20mm]{geometry}
\usepackage{amsmath,amssymb,amsthm,mathtools}
\usepackage{booktabs,longtable,array}
\usepackage{xcolor}
\usepackage[obeyspaces]{url}
\usepackage{xurl}
\usepackage[colorlinks=true,linkcolor=blue!50!black,urlcolor=blue!50!black]{hyperref}
\usepackage{microtype}
\usepackage{fontspec}
\IfFontExistsTF{Consolas}{\setmonofont{Consolas}[Scale=MatchLowercase]}{\setmonofont{Latin Modern Mono}}
\setlength{\emergencystretch}{2em}
\setlength{\parskip}{0.2em}
\Urlmuskip=0mu plus 1mu
\newcommand{\R}{\mathbb R}
\newcommand{\C}{\mathbb C}
\newcommand{\tr}{\operatorname{tr}}
\newcommand{\rank}{\operatorname{rank}}
\newcommand{\diag}{\operatorname{diag}}
\newcommand{\norm}[1]{\left\lVert #1\right\rVert}
\newcommand{\code}[1]{\ifmmode\text{\small\nolinkurl{#1}}\else{\small\nolinkurl{#1}}\fi}
\newtheorem{theorem}{Theorem}[section]
\newtheorem{lemma}[theorem]{Lemma}
\newtheorem{proposition}[theorem]{Proposition}
\newtheorem{corollary}[theorem]{Corollary}
\theoremstyle{remark}
\newtheorem{remark}[theorem]{Remark}
\numberwithin{equation}{section}
\setcounter{tocdepth}{2}
\hypersetup{pdftitle={Polynomial approximation: operator routes and finite-dimensional proofs},pdfauthor={Mathematical derivations and audit}}
\title{Polynomial approximation\\Operator routes and finite-dimensional proofs}
\author{Mathematical handoff and statement audit}
\date{12 October 2026}
\begin{document}
\maketitle
\begin{abstract}
This note treats all 35 entries currently marked open under Polynomial
approximation in the pinned NLAlib catalogue. It provides complete
mathematical derivations for 33 entries, with one copied quadrature
statement explicitly corrected. A second pass proves the universal sharp
Lebesgue lower bound, Bernstein's polynomial limit, and the reciprocal-
polynomial root rate $1/3$. It also proves the sharp rational exponent
$\pi$ for absolute value. Two entries remain partial only at their
strong rational asymptotics, including the prefactor $8$ and the Halphen
equivalent. A third pass raises the elementary absolute-value lower
prefactor to $4$, gives exact finite dual and symmetric-operator
certificates, and reduces the half-line exponential target to
polynomial targets with negligible relative error. The main tools are
finite Gram matrices, interpolation
coordinates, polynomial division, finite orthogonal projections,
positive operators, and scalar recurrences. Existing scalar bounds
transfer to Hermitian operators by one diagonalization. These are
mathematical proofs and inspected reuse notes, not Lean implementations.
\end{abstract}
\clearpage
{\small\setlength{\parskip}{0pt}\tableofcontents}
\clearpage
\section{Scope, proof status, and corrections}
\label{poly:scope}


The public catalogue snapshot is \code{95699ace3c2cddbdc2a50760d97b51719eda9cf0}.
Its Polynomial approximation group, including its five child areas,
contains 53 entries: 18 marked proved and 35 marked candidate.
The downloaded declaration index and the live site's embedded catalogue
were checked against this target list. We reuse the 18 proved entries
when applicable; they are not counted as new proof obligations.
Sources: \href{https://research.chen.pw/NLAlib/}{live catalogue},
\href{https://github.com/tchen-research/NLAlib/blob/95699ace3c2cddbdc2a50760d97b51719eda9cf0/atlas/atlas.json}{pinned atlas},
\href{https://github.com/tchen-research/NLAlib/blob/95699ace3c2cddbdc2a50760d97b51719eda9cf0/atlas/declarations.json}{pinned declaration index}.

\paragraph{What is complete and what remains.}
The 33 complete derivations preserve the requested constants and full
conclusions, except where the source statement itself needs the explicit
corrections below. Existing standard results such as polynomial density,
the spectral theorem, Cauchy's formula, and maximum modulus are reused;
the relevant interfaces are identified in the text. The following two
entries are \emph{not} claimed fully re-derived:
\begin{itemize}
\item \code{abs-best-approx-asymptotics}: Bernstein's polynomial limit
$nE_n\to\beta>0$ and the sharp rational root exponent
$-\log E_{n,n}/\sqrt n\to\pi$ are proved here. The stronger rational
limit $e^{\pi\sqrt n}E_{n,n}\to8$ remains unproved in this handoff.
The third pass proves the stronger lower bound
$\liminf e^{\pi\sqrt n}E_{n,n}\ge4$ and isolates an exact scalar
slack limit equivalent to the missing factor $8$
(Section~\ref{thirdstahl:main}).
The numerical decimal digits of $\beta$ are not certified here.
\item \code{exp-rational-halfline}: the exact type-$(0,n)$ root limit
$1/3$ is proved, with explicit finite-degree upper and lower bounds.
The diagonal Halphen equivalent $E_{n,n}\sim2H^{n+1/2}$ remains
unproved here.
\end{itemize}
These two remaining gaps are limitations of this handoff,
not requests to weaken the original catalogue statements or to mark them
proved. A consumer needing the exact sharp statement must retain that
precise theorem as an outstanding dependency.


\paragraph{What the second attempt completed.}
The universal lower bound for arbitrary interpolation nodes is now
fully proved: explicit cotangent sums give the constant, a finite
sum-of-squares Wronskian gives the required matrix signs, and a
positive normal plus a proper-map argument gives the global peak
comparison. Bernstein's polynomial limit follows from coefficient
compactness and a contracted entire comparison function. The sharp
reciprocal rate $1/3$ follows from a finite Laguerre Gram identity and
the nilpotent differentiation resolvent. For rational approximation of
absolute value, a tapered finite product gives the exact exponent
$\pi$ but not the prefactor $8$.

All of these are mathematical derivations. Their selection favors
reusable finite-dimensional steps; it does not certify that the
current Lean files already contain every required interface, or that
formalization will be easy. The two remaining results are known
theorems in mathematics, not open research conjectures.

\paragraph{What the retrieval-assisted third attempt completed.}
Section~\ref{thirdstahl:main} proves an explicit finite-degree lower
bound with asymptotic prefactor $4$ for $|x|$.  It also proves that the
optimal even rational function is a product built from finitely many
positive nodes, and identifies precisely the logarithmic slack whose
limit would recover $8$.  Sections~\ref{polythird:dual}
and~\ref{polythird:spectral} give exact finite certificates, including
a symmetric operator and certified approximate eigenpairs.
Section~\ref{polythird:finite} replaces the exponential on the half-line
by monomials on $[0,1]$ with a quantified error negligible relative to
$H^n$, then gives complete finite linear certificates for strict upper
bounds.  The sharp limits are still not proved: the missing work is a
uniform asymptotic construction and evaluation of the certificates.

The user-specified retrieval engine at
\code{E:/Database-RA01-reviewed-2026-10-06} was used in read-only mode.
Exact-name source searches found no Stahl, Halphen, Aptekarev, or
Zolotarev text in its index.  Broader searches found related-looking
finite matrix results, but their contracts did not provide the needed
bridge.  For example, \code{MS-01/T-0021}, inspected through its proof,
computes Cauchy leading coefficients for a particular triangular
RPCholesky example with dimension fixed; it supplies no uniform
large-degree rational approximation estimate.  No database card is
used as an unsupported premise.  The useful rational Remez operator
connection was verified separately in the primary paper cited below.

\paragraph{Audited corrections and endpoint conventions.}
The Gauss and eventual Clenshaw--Curtis variation denominator must be
$(2n+1-\nu)^\nu$, rather than $(n-2\nu-1)^{2\nu+1}$; this is confirmed
by the author's \href{https://people.maths.ox.ac.uk/trefethen/ATAP/}{ATAP errata}.
The sharper quadrature constants for Lebesgue measure do not extend to
arbitrary measures merely by multiplying by the total mass; an explicit
counterexample appears in Section~\ref{polyquad:errors}.
The degree-one colleague matrix is $[-a_0/a_1]$, requiring a separate
branch from the usual half-scaled last row. The zero rational function
has alternation parameter $N=m$ in type $(m,n)$; treating its numerator
as an ordinary degree-zero polynomial gives an incorrect defect.
Chebyshev coefficients use the undoubled constant coefficient, degree
means polynomial degree with $\deg0=-\infty$, and empty matrix norms are
handled without taking a maximum of an empty spectrum.

No Lean source, catalogue status, or repository proof was modified or
compiled. Source inspection identifies reuse; compilation of this TeX
document checks only typesetting. The public snapshot is newer than
parts of the working checkout, so a public declaration is not necessarily
available in the local checkout without updating it.

\subsection{Coverage map}
The two rows marked Partial are the precise sharp boundaries listed
above. All other rows have complete derivations in the indicated sections.
\small
\begin{longtable}{@{}p{0.40\textwidth}p{0.55\textwidth}@{}}
\toprule
Catalogue identifier & Route and location\\
\midrule
\endfirsthead
\toprule
Catalogue identifier & Route and location\\
\midrule
\endhead

\code{two-interval-minimax} & \S\ref{polyroot:two-interval}: Quadratic folding and an existing residual polynomial.\\[2pt]
\code{chebyshev-coefficients-def} & \S\ref{polyseries:coefficients}: A finite orthogonal projection.\\[2pt]
\code{bernstein-ellipse-def} & \S\ref{polyroot:ellipse}: A quadratic equation and a focal-distance identity.\\[2pt]
\code{minimax-error-def} & \S\ref{polybest:minimax}: Minimize on a compact box of interpolation values.\\[2pt]
\code{lebesgue-constant-def} & \S\ref{polybest:lebesgue}: Norm of a finite interpolation operator.\\[2pt]
\code{chebyshev-series-convergence} & \S\ref{polyseries:convergence}: Finite Bessel inequalities and polynomial density.\\[2pt]
\code{chebyshev-aliasing} & \S\ref{polyseries:aliasing}: A folding map on coefficient sequences.\\[2pt]
\code{chebyshev-coeff-bv-decay} & \S\ref{polyseries:bv}: Step approximations and coefficient recursion.\\[2pt]
\code{chebyshev-approx-differentiable} & \S\ref{polyseries:differentiable}: An exact falling-factorial telescope.\\[2pt]
\code{equioscillation} & \S\ref{polybest:equioscillation}: Compact sign sets and a polynomial descent direction.\\[2pt]
\code{de-la-vallee-poussin} & \S\ref{polybest:dlvp}: A finite Lagrange annihilator.\\[2pt]
\code{lebesgue-constant-chebyshev} & \S\ref{polybest:lobatto}: A finite cardinal kernel and paired cotangents.\\[2pt]
\code{bernstein-walsh-ellipse} & \S\ref{polyroot:analytic-converse}: One polynomial on the unit disk.\\[2pt]
\code{exp-approx-lower-bound} & \S\ref{polyroot:lower}: Endpoint separation and Markov.\\[2pt]
\code{slit-plane-analytic-approx} & \S\ref{polyroot:slit}: Reuse the proved open-ellipse approximation theorem.\\[2pt]
\code{sign-polynomial-approx} & \S\ref{polyroot:slit}: Approximate the inverse square root and fold.\\[2pt]
\code{sign-approx-lower-bound} & \S\ref{polyroot:lower}: A mean value and Bernstein.\\[2pt]
\code{inverse-approx-lower-bound} & \S\ref{polyroot:lower}: Pass to the residual polynomial.\\[2pt]
\code{clenshaw-recurrence} & \S\ref{polyseries:clenshaw}: The transpose of a finite triangular recurrence.\\[2pt]
\code{chebyshev-projection-lebesgue-constant} & \S\ref{polybest:projection}: Norm of a finite Dirichlet-kernel operator.\\[2pt]
\code{lebesgue-constant-lower-bound} & \S\ref{polyretry:lebesgue}: Complete: cotangent sums and finite-matrix peak comparison; equispaced case in \S\ref{polybest:lower}.\\[2pt]
\code{bernstein-converse} & \S\ref{polyroot:analytic-converse}: A geometric telescoping series on smaller ellipses.\\[2pt]
\code{hermite-integral-formula} & \S\ref{polyquad:hermite}: Finite interpolation identity and Cauchy formula.\\[2pt]
\code{chebyshev-barycentric-weights} & \S\ref{polyquad:barycentric}: Differentiate the monic node polynomial.\\[2pt]
\code{gauss-quadrature-general-measure} & \S\ref{polyquad:gauss}: Compressed multiplication in a finite Gram space.\\[2pt]
\code{chebyshev-integral-formula} & \S\ref{polyquad:integrals}: A polynomial antiderivative.\\[2pt]
\code{gauss-clenshaw-curtis-quadrature-error} & \S\ref{polyquad:errors}: Corrected rates from parity, positivity and finite alias blocks.\\[2pt]
\code{colleague-matrix} & \S\ref{polyquad:colleague}: Multiplication modulo a polynomial.\\[2pt]
\code{chebyshev-series-abs-sign} & \S\ref{polyseries:abs-sign}: Scalar cosine integrals and summation by parts.\\[2pt]
\code{chebyshev-series-exp-bessel} & \S\ref{polyseries:bessel}: An absolutely convergent scalar double series.\\[2pt]
\code{fejer-chebyshev-means} & \S\ref{polyseries:fejer}: Positive contractions and polynomial density.\\[2pt]
\code{best-approximant-interpolates} & \S\ref{polybest:bestinterpolates}: Intermediate values between alternating extrema.\\[2pt]
\code{rational-equioscillation} & \S\ref{polybest:rational}: One polynomial linear map and rank--nullity.\\[2pt]
\code{abs-best-approx-asymptotics} & \S\ref{polyretry:bernstein}, \S\ref{retryroot:sharp-rate}, \S\ref{thirdstahl:main}: Partial: polynomial limit, rational exponent, and lower prefactor $4$ complete; limit $8$ remains.\\[2pt]
\code{exp-rational-halfline} & \S\ref{polyretry:exp}, \S\ref{polythird:spectral}, \S\ref{polythird:finite}: Partial: sharp reciprocal rate and exact finite reductions complete; diagonal Halphen equivalent remains.\\[2pt]

\bottomrule
\end{longtable}
\normalsize

\section{A single finite-dimensional operator transfer}
\label{polyroot:transfer}
Let $A=A^*$ act on a finite-dimensional real or complex inner-product
space. Choose an orthonormal eigenbasis, so
$A=U\diag(\lambda_1,\ldots,\lambda_d)U^*$.
For a scalar function $f$ defined on these eigenvalues (real-valued when
the vector space is real), define
$f(A)=U\diag(f(\lambda_i))U^*$. This definition is independent of the
choice of basis inside a repeated eigenspace. Polynomial evaluation
agrees with it, by the finite multiplication rule for diagonal matrices.
For any polynomial $p$ and any set $S$ containing the spectrum,
\begin{equation}\label{polyroot:transfer-bound}
 \|f(A)-p(A)\|=\max_i|f(\lambda_i)-p(\lambda_i)|
 \le \sup_{x\in S}|f(x)-p(x)|.
\end{equation}
The equality is stated for $d>0$; for $d=0$ the operator and its norm
are zero. Indeed, the squared norm of a diagonal operator applied to
$v$ is $\sum_i|f(\lambda_i)-p(\lambda_i)|^2|v_i|^2$; a coordinate
vector attaining the maximum gives the reverse inequality.
Consequently the same error multiplied by $\|v\|$ bounds a vector action,
and $d$ times the error bounds a trace error. For real $f,p$, a scalar
bound $|f-p|\le\varepsilon$ also gives the two Loewner inequalities
$-\varepsilon I\preceq f(A)-p(A)\preceq\varepsilon I$.

This is the useful division of work: construct or estimate one scalar
polynomial, then apply the diagonal identity once. It does not require
matrix concentration, a spectral measure, or an integral functional
calculus. It applies to rational functions as well whenever their
denominators are nonzero on the spectrum. No analogous norm-one transfer
is asserted for arbitrary nonnormal matrices.

\paragraph{Inspected reuse.}
Mathlib already provides \code{Matrix.IsHermitian.cfc_eq},
\code{Matrix.IsHermitian.cfc}, and \code{norm_cfc_le}. The last theorem
can be applied to the scalar error function once its spectrum-side bound
is known; the diagonal proof above explains the complete mathematics.

\section{Folding two intervals into one positive interval}
\label{polyroot:two-interval}
\paragraph{Target.}\code{two-interval-minimax}.
Assume $a<b<0<c<d$ and $d-c=b-a$. Put
\[
 s=b+c=a+d,\quad h(x)=x^2-sx,\quad
 \alpha=-bc>0,\quad\beta=-ad>\alpha.
\]
The endpoint identities are
$h(b)=h(c)=\alpha$ and $h(a)=h(d)=\beta$.
For $x<y\le b$,
$h(y)-h(x)=(y-x)(x+y-s)<0$, since $x+y-s<2b-(b+c)<0$.
For $c\le x<y$, the same difference is positive. By continuity,
each of $[a,b]$ and $[c,d]$ maps onto $[\alpha,\beta]$.
In particular $\beta>\alpha$ and $h(0)=0$.

For $k\ge0$ set
\[
 R_k(t)=
 \frac{T_k((\beta+\alpha-2t)/(\beta-\alpha))}
 {T_k((\beta+\alpha)/(\beta-\alpha))},
 \qquad p(x)=R_k(h(x)).
\]
The denominator is positive, $\deg p\le2k$, and $p(0)=1$.
Write $\eta=(\sqrt\beta-\sqrt\alpha)/(\sqrt\beta+\sqrt\alpha)\in(0,1)$.
The usual Chebyshev recurrence gives
\[
 T_k\!\left(\frac{\beta+\alpha}{\beta-\alpha}\right)
 =\frac{\eta^{-k}+\eta^k}{2}.
\]
Together with $|T_k|\le1$ on $[-1,1]$, this proves
\begin{equation}\label{polyroot:two-interval-bound}
 \sup_{[a,b]\cup[c,d]}|p|
 \le\frac{2\eta^k}{1+\eta^{2k}}\le2\eta^k.
\end{equation}
For $k=0$ the actual polynomial is $1$ and its error bound is $1$.
In the symmetric case $a=-d,b=-c$, $h(x)=x^2$ and
$\eta=(d-c)/(d+c)$. For a Hermitian $A$ with spectrum in the two
intervals, the entire construction is $p(A)=R_k(A^2-sA)$; the operator
$A^2-sA$ has spectrum in one positive interval.
The already proved \code{NLAlib.isLeast_chebyshev_minimax} supplies the
one-interval minimax contract, so only the quadratic endpoint algebra is new.

\section{Bernstein ellipses from a quadratic equation}
\label{polyroot:ellipse}
\paragraph{Target.}\code{bernstein-ellipse-def}.
For $\rho>1$ define
\[
 E_\rho=\{z\in\C:|z-1|+|z+1|<\rho+\rho^{-1}\},
 \qquad J(w)=\tfrac12(w+w^{-1})\quad(w\ne0).
\]
The key identity is purely algebraic:
\begin{equation}\label{polyroot:joukowski-distance}
 |J(w)-1|+|J(w)+1|
 =\frac{|w-1|^2+|w+1|^2}{2|w|}
 =|w|+|w|^{-1}.
\end{equation}
For any $z\in\C$, the quadratic $w^2-2zw+1=0$ has two roots with
product $1$. Both are nonzero, one has modulus at least $1$, and its
Joukowski image is $z$. The function $r\mapsto r+r^{-1}$ is strictly
increasing for $r\ge1$ (with strict increase between distinct values).
Thus \eqref{polyroot:joukowski-distance} proves
\[
 E_\rho=J(\{w:\rho^{-1}<|w|<\rho\})
 =J(\{w:1\le|w|<\rho\}).
\]
The closed inequality similarly describes
$J(\{w:1\le|w|\le\rho\})$. This is the closure of $E_\rho$:
points with $|w|=\rho$ are limits of $J(rw/\rho)$ as $r\uparrow\rho$;
all other points are already in $E_\rho$. The distance definition proves
openness directly, and also monotonicity in $\rho$ and inclusion of
$[-1,1]$, whose focal-distance sum is $2$.

Let $A_\rho=(\rho+\rho^{-1})/2$. For $z=J(w)$ with
$1\le|w|\le\rho$,
\[
 |z|\le A_\rho,\qquad \operatorname{Re}z\ge-A_\rho.
\]
Induction in the three-term recurrence, starting with $T_0=1,T_1=x$,
gives $T_n(J(w))=(w^n+w^{-n})/2$. Therefore
\[
 |T_n(z)|\le(\rho^n+\rho^{-n})/2
 \qquad(z\in\overline E_\rho).
\]
For $a<b$ and $\varphi(z)=((b-a)z+b+a)/2$, the triangle inequality gives
\[
 \operatorname{Re}\varphi(z)\ge\frac{b+a-(b-a)A_\rho}{2},\qquad
 |\varphi(z)|\le\frac{|b+a|+(b-a)A_\rho}{2}.
\]
These are all the catalogue's geometric bounds. The quadratic equation
and \eqref{polyroot:joukowski-distance} avoid a separate parametrized
ellipse geometry library. The objects \code{NLAlib.openBernsteinEllipse},
\code{NLAlib.filledBernsteinEllipse}, and \code{NLAlib.joukowski}
already exist; this candidate is an interface-completion task, not a
reason to introduce competing definitions.

\section{Polynomial growth and the analytic converse}
\label{polyroot:analytic-converse}
\paragraph{Targets.}\code{bernstein-walsh-ellipse}, \code{bernstein-converse}.
Let $p(z)=\sum_{j=0}^n c_jz^j$ be complex, with $|p(x)|\le M$ on
$[-1,1]$. In particular $M\ge0$. The expression
\[
 g(v)=\sum_{j=0}^n c_j2^{-j}v^{n-j}(1+v^2)^j
\]
is an ordinary polynomial; for $v\ne0$ it equals $v^np(J(v))$.
On $|v|=1$, $J(v)=\operatorname{Re}v\in[-1,1]$, hence $|g(v)|\le M$.
The maximum-modulus theorem on the closed unit disk gives $|g(v)|\le M$
inside it. Substituting $v=1/w$ proves
\begin{equation}\label{polyroot:walsh}
 |p(J(w))|\le M|w|^n\quad(|w|\ge1),\qquad
 |p(z)|\le M\rho^n\quad(z\in\overline E_\rho).
\end{equation}
This uses one standard disk theorem, with no singularity at $v=0$
because $g$ was explicitly expanded as a polynomial.

Now suppose $\deg q_n\le n$ and
$\|f-q_n\|_{[-1,1]}\le C\rho^{-n}$, where $\rho>1$ and $C>0$.
Put $d_j=q_j-q_{j-1}$ for $j\ge1$. Then
\[
 \deg d_j\le j,\qquad
 \|d_j\|_{[-1,1]}\le C(1+\rho)\rho^{-j}.
\]
For $1<r<\rho$, \eqref{polyroot:walsh} yields
\[
 \|d_j\|_{\overline E_r}\le C(1+\rho)(r/\rho)^j.
\]
Thus the telescoping series $q_0+\sum_{j\ge1}d_j$ converges uniformly
on every $\overline E_r$ to a function $F$. The definitions on different
ellipses agree because they are limits of the same $q_n$.
Every compact subset of $E_\rho$ lies in some $E_r$ with $r<\rho$:
maximize its continuous focal-distance sum and choose $r$ above the
corresponding maximum radius. Hence convergence is locally uniform, and
the standard locally-uniform limit theorem for holomorphic functions
makes $F$ holomorphic on $E_\rho$. On $[-1,1]$ its limit is $f$.
The proof supplies the explicit interior error
\begin{equation}\label{polyroot:analytic-tail}
 \|F-q_N\|_{\overline E_r}
 \le\frac{C(1+\rho)(r/\rho)^{N+1}}{1-r/\rho}.
\end{equation}
For the catalogue's limsup variant, fix any $1<r<s<\rho$.
The assumption $\limsup\|f-q_n\|^{1/n}\le\rho^{-1}$ implies
$\|f-q_n\|\le C_s s^{-n}$ for all $n$, after enlarging $C_s$ to
handle finitely many initial indices. Apply the preceding argument
with $s$, then exhaust $E_\rho$ by such interiors.
The only analytic ingredients are maximum modulus and preservation of
holomorphy under locally uniform limits; the estimates are finite
polynomial algebra followed by a geometric series.
The inspected Mathlib interfaces are
\code{Complex.norm_le_of_forall_mem_frontier_norm_le} in
\path{Mathlib/Analysis/Complex/AbsMax.lean} and
\code{TendstoLocallyUniformlyOn.differentiableOn} in
\path{Mathlib/Analysis/Complex/LocallyUniformLimit.lean}.

\section{Right-half-plane functions and a polynomial sign filter}
\label{polyroot:slit}
\paragraph{Targets.}\code{slit-plane-analytic-approx}, \code{sign-polynomial-approx}.
Let $0<a<b$, and set
\[
 \rho_0=\frac{\sqrt b+\sqrt a}{\sqrt b-\sqrt a},\qquad
 m_r=\frac{b+a-(b-a)(r+r^{-1})/2}{2},\qquad
 R_r=\frac{b+a+(b-a)(r+r^{-1})/2}{2}.
\]
A direct calculation gives
$\rho_0+\rho_0^{-1}=2(b+a)/(b-a)$.
Therefore $1<r<\rho_0$ implies $m_r>0$.
By the ellipse geometry above, the affine image of $\overline E_r$
has $\operatorname{Re}z\ge m_r$ and $|z|\le R_r$.
If $F$ is holomorphic in the right half-plane, real on $[a,b]$, and
bounded by $M$ on this enclosure, the already proved open-ellipse
approximation theorem applied to $F\circ\varphi$ gives a real $p$ with
\begin{equation}\label{polyroot:half-plane-approx}
 \deg p\le n,\qquad
 \sup_{[a,b]}|F-p|\le\frac{2M r^{-n}}{r-1}.
\end{equation}
No boundary regularity assumption has been added: the enclosure lies
strictly inside the right half-plane.

For real $\alpha$, choose the principal branch
$z^\alpha=\exp(\alpha\operatorname{Log}z)$. Since
$|z^\alpha|=|z|^\alpha$ and $m_r\le|z|\le R_r$, one may take
$M=\max(m_r^\alpha,R_r^\alpha)$.
Similarly $|\operatorname{Arg}z|<\pi/2$ gives
$M=\max(|\log m_r|,|\log R_r|)+\pi/2$ for $F=\operatorname{Log}$.
These are scalar branch facts; operator evaluation is then
\eqref{polyroot:transfer-bound}.

To approximate sign on $\gamma\le|x|\le1$, take $a=\gamma^2,b=1$
and $F(t)=t^{-1/2}$. Equation~\eqref{polyroot:half-plane-approx}
gives $q$ of degree at most $n$; set $p(x)=xq(x^2)$. Then
\[
 \deg p\le2n+1,\qquad
 |\operatorname{sgn}x-p(x)|
 =|x|\,|(x^2)^{-1/2}-q(x^2)|
 \le\frac{2m_r^{-1/2}r^{-n}}{r-1}.
\]
For $0<\gamma<1$, $r=1+\gamma$ lies strictly below
$(1+\gamma)/(1-\gamma)$, and exact simplification gives
$m_r=\gamma^2(3+\gamma)/4$. Hence
\begin{equation}\label{polyroot:sign-bound}
 \sup_{\gamma\le|x|\le1}|\operatorname{sgn}x-p(x)|
 \le\frac{4}{\sqrt3\,\gamma^2}(1+\gamma)^{-n}.
\end{equation}
For $0<\varepsilon<1$, it suffices to choose
\[
 n=\left\lceil
 \frac{\log(4/(\sqrt3\gamma^2\varepsilon))}{\log(1+\gamma)}
 \right\rceil.
\]
For a Hermitian contraction with no eigenvalues in $(-\gamma,\gamma)$,
the same polynomial approximates $\operatorname{sgn}(A)$ in operator
norm. The polynomial $(I+p(A))/2$ approximates the positive spectral
projector with half the error.

\paragraph{Inspected reuse.}
Use \code{NLAlib.exists_degree_le_abs_sub_eval_le_of_open_bernsteinEllipse}
or its \code{_interval} version from
\path{NLAlib/Polynomial/FunctionApproximation/AnalyticApproximation.lean}.
The latter takes the holomorphic extension in the normalized ellipse
variable, so pass $F\circ\varphi$, not the untranslated $F$.

\section{Three degree lower bounds from existing scalar inequalities}
\label{polyroot:lower}
\paragraph{Target.}\code{exp-approx-lower-bound}.
Suppose $b\ge1$, $0\le\delta\le1/8$, $\deg p\le q$, and
$|p(x)-e^{-x}|\le\delta$ on $[0,b]$. Then
\[
 p(0)-p(1)\ge1-e^{-1}-2\delta>0,
 \qquad \|p\|_{[0,b]}\le1+\delta.
\]
The mean-value theorem gives $\xi\in(0,1)$ with
$|p'(\xi)|\ge1-e^{-1}-2\delta$. The rescaled Markov inequality gives
$|p'(\xi)|\le 2q^2(1+\delta)/b$, so
\begin{equation}\label{polyroot:exp-lower}
 q^2\ge\frac{(1-e^{-1}-2\delta)b}{2(1+\delta)}\ge b/6.
\end{equation}
For the final numerical inequality, $e>1+1+1/2+1/6=8/3$ implies
$e^{-1}<3/8$; then the numerator is at least $3/8$ and the denominator
at most $9/4$. This retains the catalogue's exact constants.

\paragraph{Target.}\code{sign-approx-lower-bound}.
Assume $0<\gamma<1$, $0\le\varepsilon<1$, $\deg p\le q$,
$\|p\|_{[-1,1]}\le M$, and $p$ approximates sign within $\varepsilon$
on the two outer intervals. The mean-value theorem between $-\gamma$
and $\gamma$ gives $\xi\in(-\gamma,\gamma)$ with
$p'(\xi)\ge(1-\varepsilon)/\gamma$.
Bernstein's weighted derivative inequality then gives
\begin{equation}\label{polyroot:sign-lower}
 qM\ge |p'(\xi)|\sqrt{1-\xi^2}
 \ge\frac{(1-\varepsilon)\sqrt{1-\gamma^2}}{\gamma}.
\end{equation}
The global bound $M$ is a genuine hypothesis: an outer-interval error
bound by itself does not bound the polynomial inside the gap.

\paragraph{Target.}\code{inverse-approx-lower-bound}.
Let $0<a<b$, $\deg p<q$, and $|p(x)-1/x|\le\varepsilon$ on $[a,b]$.
The residual polynomial $r(x)=1-xp(x)$ satisfies $r(0)=1$,
$\deg r\le q$, and $\|r\|_{[a,b]}\le b\varepsilon$.
The already proved constrained Chebyshev minimax theorem therefore gives
\begin{equation}\label{polyroot:inverse-lower}
 \varepsilon\ge
 \frac{1}{bT_q((b+a)/(b-a))}
 =\frac{2\eta^q}{b(1+\eta^{2q})}
 \ge\frac{\eta^q}{b},\qquad
 \eta=\frac{\sqrt b-\sqrt a}{\sqrt b+\sqrt a}.
\end{equation}
This also handles $q=0$: the condition $\deg p<0$ forces $p=0$,
and $r=1$. Use polynomial degree with $\deg0=-\infty$ to state this
case faithfully, rather than natural degree.

\paragraph{Inspected reuse.}
The exact public declarations are
\code{NLAlib.abs_eval_derivative_le_two_mul_sq_div_mul_of_forall_abs_eval_le},
\code{NLAlib.abs_eval_derivative_le_div_sqrt_of_forall_abs_eval_le},
and \code{NLAlib.isLeast_chebyshev_minimax}. Thus the three new proofs
are applications of existing scalar bounds plus endpoint evaluation.



% BEGIN series.tex
\section{Chebyshev series through finite projections and positive operators}
\label{polyseries:section}

This section covers \code{chebyshev-coefficients-def},
\code{chebyshev-series-convergence}, \code{chebyshev-aliasing},
\code{chebyshev-coeff-bv-decay}, \code{chebyshev-approx-differentiable},
\code{clenshaw-recurrence}, \code{chebyshev-series-abs-sign},
\code{chebyshev-series-exp-bessel}, and \code{fejer-chebyshev-means}.
The useful reductions are finite orthogonal projections, summation by parts,
and positive contractions. No completeness theorem for a Fourier basis,
Rademacher theorem, or Stieltjes integral is required.

\subsection{Coefficient functionals and finite projections}
\label{polyseries:coefficients}
For $k\in\mathbb N$, write $T_k(\cos\theta)=\cos(k\theta)$, and define
\begin{equation}\label{polyseries:coeffdef}
 a_0(f)=\frac1\pi\int_0^\pi f(\cos\theta)\,d\theta,
 \qquad
 a_k(f)=\frac2\pi\int_0^\pi f(\cos\theta)\cos(k\theta)\,d\theta
 \quad(k\ge1),\qquad
 P_nf=\sum_{k=0}^n a_k(f)T_k.
\end{equation}
These formulas also make sense for bounded measurable $f$; continuity is
needed for the uniform approximation assertions below. The $k=0$
normalization is not doubled.

The substitution $x=\cos\theta$ gives
\[
 a_k(f)=\frac{2-\mathbf1_{k=0}}\pi
 \int_{-1}^1\frac{f(x)T_k(x)}{\sqrt{1-x^2}}\,dx.
\]
The endpoint singularities are integrable and their point values do not
affect the integral. Product-to-sum and the integral of $\cos(m\theta)$ give
\[
 \frac1\pi\int_0^\pi\cos(j\theta)\cos(k\theta)\,d\theta
 =\begin{cases}1&j=k=0,\\1/2&j=k>0,\\0&j\ne k.\end{cases}
\]
Consequently $a_k$ is linear, $a_k(T_j)=\mathbf1_{k=j}$, and
\begin{equation}\label{polyseries:coeffbound}
 |a_0(f)|\le\|f\|_\infty,\qquad |a_k(f)|\le2\|f\|_\infty\quad(k\ge1).
\end{equation}
The degrees and nonzero leading coefficients of $T_0,\ldots,T_n$ show by
finite triangular elimination that they form a basis of the polynomials
of degree at most $n$. Thus $P_np=p$ when $\deg p\le n$, and $a_k(p)=0$
for $k>\deg p$. Equivalently, $P_n$ is the orthogonal projection onto this
finite polynomial space for the probability inner product
\[
 \langle f,h\rangle_\mu=\frac1\pi\int_0^\pi
       f(\cos\theta)h(\cos\theta)\,d\theta.
\]
The norm bounds in \eqref{polyseries:coeffbound} concern the individual
coefficient functionals; they do not say that $\|P_n\|\le2$.

\subsection{Absolute convergence from finite Bessel inequalities}
\label{polyseries:convergence}
Suppose first that $f$ is continuous and $\sum_{k\ge0}|a_k(f)|<\infty$.
Since $|T_k(x)|\le1$ on $[-1,1]$, the series has a continuous uniform sum
$h$. Integrating its uniformly convergent partial sums and using
orthogonality gives $a_k(h)=a_k(f)$ for every $k$. Here is an elementary
way to identify $h$, rather than invoking Fourier inversion.

\begin{lemma}[Uniqueness from polynomial density]\label{polyseries:uniqueness}
If $v\in C[-1,1]$ and $a_k(v)=0$ for all $k$, then $v=0$.
\end{lemma}
\begin{proof}
Every polynomial $p$ satisfies $\langle v,p\rangle_\mu=0$. For
$\varepsilon>0$, polynomial density supplies $p$ with
$\|v-p\|_\infty<\varepsilon$. Hence
$\int v^2\,d\mu=\int v(v-p)\,d\mu\le\varepsilon\|v\|_\infty$.
Letting $\varepsilon$ tend to zero gives zero integral. If $v(x_0)\ne0$,
continuity gives a relative interval on which $v^2$ is bounded below by a
positive number; its $\mu$-measure is positive, including when $x_0$ is
an endpoint. This contradicts the zero integral.
\end{proof}
Applying the lemma to $v=f-h$ proves
\begin{equation}\label{polyseries:tail}
 f=\sum_{k=0}^\infty a_k(f)T_k\quad\hbox{uniformly},\qquad
 \|f-P_nf\|_\infty\le\sum_{k>n}|a_k(f)|.
\end{equation}
Any absolutely and uniformly convergent Chebyshev expansion has these
same coefficients, by termwise integration. This is the asserted uniqueness.

There is also a short proof of the catalogue's full Lipschitz assertion.
Let $f$ be $L$-Lipschitz and $g(\theta)=f(\cos\theta)$ on $\R$.
Then $g$ is even, $2\pi$-periodic, and $L$-Lipschitz. Normalize integrals
over a period by $\langle u,v\rangle=(2\pi)^{-1}\int_{-\pi}^\pi uv$.
For $h\ne0$, put $d_h(\theta)=(g(\theta+h)-g(\theta))/h$.
We have $|d_h|\le L$. Translation of a periodic integral gives, for $k\ge1$,
\[
 \langle d_h,\cos(k\cdot)\rangle
    =\frac{a_k(f)}{2h}(\cos(kh)-1),\qquad
 \langle d_h,\sin(k\cdot)\rangle
    =-\frac{a_k(f)}{2h}\sin(kh).
\]
For each finite $N$, apply the finite orthogonal-projection inequality to
$\sqrt2\cos(k\theta),\sqrt2\sin(k\theta)$, $1\le k\le N$:
\begin{equation}\label{polyseries:finitebessel}
 \sum_{k=1}^N a_k(f)^2\frac{1-\cos(kh)}{h^2}\le L^2.
\end{equation}
This inequality itself follows by expanding the nonnegative integral of
the squared difference between $d_h$ and its finite orthogonal projection.
Taking $h\to0$ in this \emph{finite} sum yields
\begin{equation}\label{polyseries:lipschitzenergy}
 \sum_{k=1}^N k^2a_k(f)^2\le2L^2\qquad\hbox{for every }N.
\end{equation}
Finite Cauchy--Schwarz and $\sum_{k\ge1}k^{-2}\le2$ give
$\sum_{k=1}^N|a_k(f)|\le2L$. The increasing partial sums therefore
converge, proving absolute summability and \eqref{polyseries:tail}.
One also obtains the explicit, though nonoptimal, bound
\[
 \sum_{k\ge1}|a_k(f)|\le2L,\qquad
 \|f-P_nf\|_\infty\le\frac{\sqrt2 L}{\sqrt n}\quad(n\ge1),
\]
because $\sum_{k>n}k^{-2}\le1/n$. All limiting steps involve scalar
monotone sums or a finite-dimensional projection; no derivative of $g$
has been asserted to exist.

\subsection{Aliasing as a folding operator on coefficient sequences}
\label{polyseries:aliasing}
Let $n\ge1$, $x_j=\cos(j\pi/n)$ for $0\le j\le n$, and let $r_n(k)$
be the fold into $\{0,\ldots,n\}$: if $s$ is the remainder of $k$ modulo
$2n$, set $r_n(k)=\min(s,2n-s)$. Directly,
\begin{equation}\label{polyseries:alias}
 T_k(x_j)=\cos(kj\pi/n)=T_{r_n(k)}(x_j).
\end{equation}
In particular $T_m,T_{2\ell n+m},T_{2\ell n-m}$ have the same values
for $0\le m\le n$, $\ell\ge1$. Since the nodes are distinct, polynomial
interpolation and degree at most $n$ imply $I_nT_k=T_{r_n(k)}$.

If $a=(a_k)_{k\ge0}$ is absolutely summable, define the finite vector
$c_j=\sum_{r_n(k)=j}a_k$. Absolute convergence justifies this partition,
and $\sum_{j=0}^n|c_j|\le\sum_{k\ge0}|a_k|$. Thus folding is a
contraction from $\ell^1(\mathbb N)$ to $\ell^1(\{0,\ldots,n\})$.
For $a_k=a_k(f)$ and $f$ continuous, \eqref{polyseries:tail} and
\eqref{polyseries:alias} show that
$p_n=\sum_{j=0}^nc_jT_j$ interpolates $f$. Explicitly,
\begin{align*}
 c_0&=a_0+\sum_{\ell\ge1}a_{2\ell n},&
 c_n&=\sum_{\ell\ge0}a_{(2\ell+1)n},\\
 c_j&=a_j+\sum_{\ell\ge1}(a_{2\ell n+j}+a_{2\ell n-j})
 &&(1\le j<n).
\end{align*}
The endpoint formulas avoid double counting. For $n=1$ there are no
interior coefficients. Moreover,
\begin{equation}\label{polyseries:aliaserror}
 f-p_n=\sum_{k>n}a_k(T_k-T_{r_n(k)}),\qquad
 \|f-p_n\|_\infty\le2\sum_{k>n}|a_k|.
\end{equation}
This proof requires only a finite partition of an absolutely convergent
series; continuity of the interpolation operator need not be developed.

At the $n+1$ roots $\cos((2j+1)\pi/(2n+2))$ of $T_{n+1}$, the
corresponding fold has signs: reduce modulo $2(n+1)$, use
$T_{2(n+1)\pm m}=-T_m$ on the grid, and use $T_{n+1}=0$ there.
Successive reductions give $I_nT_k$ equal to $T_r$, $-T_r$, or zero,
with $r\le n$. Hence $\|I_nT_k\|_\infty\le1$ and the same factor-two
tail bound follows. This variant includes $n=0$: the single root is zero.

\subsection{Bounded variation without a Stieltjes integral}
\label{polyseries:bv}
Use the following derivative interpretation, which also covers $|x|$:
there are functions $f_0=f,f_1,\ldots,f_\nu$ such that $f_j$ is absolutely
continuous for $j<\nu$, $f_j'=f_{j+1}$ almost everywhere, and $f_\nu$ has
a bounded-variation representative with variation $V$ on $[-1,1]$.
Classical derivatives satisfy this formulation whenever the catalogue's
hypotheses hold. Values assigned to derivatives at exceptional points
do not change coefficients. For $\nu=0$ only the coefficient estimate is
claimed; a BV function need not be continuous.

\begin{lemma}[An elementary oscillatory integral bound]
\label{polyseries:bvbase}
If $h$ has variation at most $V$ on $[0,\pi]$, then, for every integer $k\ge1$,
\[
 \left|\int_0^\pi h(\theta)\cos(k\theta)\,d\theta\right|\le\frac Vk.
\]
\end{lemma}
\begin{proof}
Partition $[0,\pi]$ into $N$ equal intervals with endpoints $t_j=j\pi/N$.
Let $h_N=h(t_j)$ on $[t_j,t_{j+1})$. On each cell the error is bounded
by the variation on that cell. Additivity of variation over adjacent
intervals therefore gives
$\int_0^\pi|h-h_N|\le\pi V/N$.
Writing $s_j=\sin(kt_j)$, direct integration and finite summation by parts give
\[
 \int_0^\pi h_N\cos(k\theta)\,d\theta
  =\frac1k\sum_{j=0}^{N-1}h(t_j)(s_{j+1}-s_j)
  =-\frac1k\sum_{j=1}^{N-1}(h(t_j)-h(t_{j-1}))s_j.
\]
The boundary terms vanish because $s_0=s_N=0$. The last sum has absolute
value at most $V/k$. The integral error tends to zero, proving the claim.
BV functions are bounded and measurable; alternatively these same step
approximations supply the needed integrability.
\end{proof}
Composition with the monotone map $\cos:[0,\pi]\to[-1,1]$ does not
increase variation. The lemma immediately yields
\begin{equation}\label{polyseries:bvzero}
 |a_k(g)|\le\frac{2\operatorname{Var}(g)}{\pi k}\quad(k\ge1).
\end{equation}

For an absolutely continuous $u$ with bounded derivative representative
$v$, integration by parts gives, for $k\ge2$,
\begin{equation}\label{polyseries:coeffrecursion}
 a_k(u)=\frac{a_{k-1}(v)-a_{k+1}(v)}{2k}.
\end{equation}
To check the endpoint issue explicitly, put
$F_k(x)=\sin(k\arccos x)$. This is absolutely continuous, vanishes at
both endpoints, and
$F_k'(x)=-kT_k(x)/\sqrt{1-x^2}$ in the interior. Its derivative is
integrable; the derivative identity integrated on trimmed intervals and
then at the endpoints proves absolute continuity. Thus
\[
 a_k(u)=\frac2{\pi k}\int_{-1}^1v(x)F_k(x)\,dx
 =\frac2{\pi k}\int_0^\pi
 v(\cos\theta)\sin\theta\sin(k\theta)\,d\theta.
\]
Product-to-sum proves \eqref{polyseries:coeffrecursion}. For $k=1$ its
correct exceptional form is $a_1(u)=a_0(v)-a_2(v)/2$; it is not needed
in the following induction.

\begin{theorem}[The full BV coefficient bound]\label{polyseries:bvtheorem}
For $k\ge\nu+1$,
\begin{equation}\label{polyseries:bvbound}
 |a_k(f)|\le\frac{2V}{\pi\,k(k-1)\cdots(k-\nu)}.
\end{equation}
\end{theorem}
\begin{proof}
The case $\nu=0$ is \eqref{polyseries:bvzero}. Suppose $\nu\ge1$.
Then $k\ge2$, and the induction hypothesis applied to $f_1$ at $k-1$
and $k+1$ bounds each coefficient in \eqref{polyseries:coeffrecursion}
by
$2V/[\pi(k-1)(k-2)\cdots(k-\nu)]$: the bound at $k+1$ has larger
positive factors in its denominator. Divide their sum by $2k$.
Every function used as a derivative representative is bounded: it is
continuous when another AC level remains, and BV at the last level.
Hence all the integrations by parts just used are legitimate.
\end{proof}
This finite recursion proves the same constants as the source while
avoiding its contour notation and its final BV integration-by-parts step.
For a continuously differentiable last derivative one may take
$V=\int_{-1}^1|f^{(\nu+1)}(x)|\,dx$, but this stronger regularity is
unnecessary.

\subsection{Algebraic approximation rates by an exact telescope}
\label{polyseries:differentiable}
Assume the preceding hypotheses with $\nu\ge1$ and $n\ge\nu$.
The positive summands admit the identity
\[
 \frac1{k(k-1)\cdots(k-\nu)}
 =\frac1\nu\left(
 \frac1{(k-1)\cdots(k-\nu)}-
 \frac1{k\cdots(k-\nu+1)}\right).
\]
Summing and letting the final denominator tend to infinity gives
\begin{equation}\label{polyseries:sharptail}
 \sum_{k=n+1}^\infty|a_k(f)|
 \le\frac{2V}{\pi\nu\,n(n-1)\cdots(n-\nu+1)}.
\end{equation}
There are $\nu$ factors in the denominator. This proves summability,
so \eqref{polyseries:tail} and \eqref{polyseries:aliaserror} apply and give
\[
 \|f-P_nf\|_\infty\le
 \frac{2V}{\pi\nu\,n(n-1)\cdots(n-\nu+1)},\qquad
 \|f-I_nf\|_\infty\le
 \frac{4V}{\pi\nu\,n(n-1)\cdots(n-\nu+1)}.
\]
For the catalogue's $n>\nu$, every denominator factor is at least
$n-\nu>0$, giving exactly $2V/[\pi\nu(n-\nu)^\nu]$ and
$4V/[\pi\nu(n-\nu)^\nu]$. The best degree-$n$ error is no larger
than the truncation error. If $V=0$, the high coefficients are zero;
the result remains valid without dividing by $V$.

\subsection{Clenshaw as the transpose of a finite recurrence}
\label{polyseries:clenshaw}
Take scalars $\alpha_k,\beta_k$ and vectors $y_k$ satisfying
$y_{k+1}=\alpha_ky_k-\beta_ky_{k-1}$ for $k\ge1$.
For $n\ge0$, define $u_{n+1}=u_{n+2}=0$ and, backwards,
\[
 u_k=c_k+\alpha_ku_{k+1}-\beta_{k+1}u_{k+2}
 \qquad(0\le k\le n).
\]
Then
\begin{equation}\label{polyseries:clenshawidentity}
 \sum_{k=0}^nc_ky_k=u_0y_0+u_1(y_1-\alpha_0y_0).
\end{equation}
Indeed substitute
$c_k=u_k-\alpha_ku_{k+1}+\beta_{k+1}u_{k+2}$ into the finite sum.
The coefficient of $u_j$ for $2\le j\le n$ is
$y_j-\alpha_{j-1}y_{j-1}+\beta_{j-1}y_{j-2}=0$.
The $u_0,u_1$ terms are exactly those displayed, and the remaining
terms involve the prescribed zeros. For $n=0$, $u_1=0$ and the formula
is simply $c_0y_0$.

The catalogue's general recurrence is $\alpha_k=t-a_k$,
$\beta_k=b_k$. For Chebyshev take $\alpha_k=2t$, $\beta_k=1$,
$y_0=1$, $y_1=t$, giving
\[
 \sum_{k=0}^nc_kT_k(t)=u_0-tu_1
 =c_0+tu_1-u_2.
\]
For an arbitrary linear operator $A$ on a finite-dimensional real or
complex vector space and a vector $b$, compute vectors
\[
 v_{n+1}=v_{n+2}=0,\qquad
 v_k=c_kb+2Av_{k+1}-v_{k+2}\quad(k=n,\ldots,1).
\]
Then the answer is $c_0b+Av_1-v_2$. The polynomial identity proves it
equals $\sum_{k=0}^nc_kT_k(A)b$. When $n\ge1$, initialize $v_n=c_nb$;
there are $n-1$ products for $v_{n-1},\ldots,v_1$ and one for the final
answer, hence $n$ matrix--vector products. For $n=0$ return $c_0b$
without a product. No symmetry, spectral restriction, or invertibility
is required. The interpretation is a finite triangular system and its
transpose, rather than an infinite recurrence limit.

\subsection{The absolute value and sign expansions}
\label{polyseries:abs-sign}
For $f(x)=|x|$, symmetry makes the odd coefficients zero. Splitting the
coefficient integral at $\pi/2$ and applying product-to-sum gives
\[
 a_0=\frac2\pi,\qquad
 a_{2j}=\frac4\pi\frac{(-1)^{j+1}}{4j^2-1}\quad(j\ge1).
\]
For example the calculation is
$a_{2j}=(4/\pi)\int_0^{\pi/2}\cos\theta\cos(2j\theta)\,d\theta$;
the two endpoint sine terms simplify to the displayed fraction.
Absolute summability and \eqref{polyseries:tail} prove the uniform series.
Since
$1/(4j^2-1)=\tfrac12((2j-1)^{-1}-(2j+1)^{-1})$,
\begin{equation}\label{polyseries:abstail}
 \left\||x|-\frac2\pi-\frac4\pi\sum_{j=1}^m
      \frac{(-1)^{j+1}T_{2j}(x)}{4j^2-1}\right\|_\infty
 \le\frac2{\pi(2m+1)}\quad(m\ge0).
\end{equation}
In fact equality holds at $x=0$: $T_{2j}(0)=(-1)^j$ makes every
omitted term have the same negative sign.

For $\operatorname{sign}(0)=0$, direct integration gives zero even
coefficients and
$a_{2j+1}=4(-1)^j/[\pi(2j+1)]$. These coefficients are not absolutely
summable, so a separate argument is needed to prove the pointwise
series. Put $\phi=\arcsin x$, so that
$(-1)^jT_{2j+1}(x)=\sin((2j+1)\phi)$. For $\phi\ne0$, finite geometric
sums give
\[
 \left|\sum_{j=p}^q\sin((2j+1)\phi)\right|
 \le\frac1{|\sin\phi|}.
\]
Finite summation by parts with the decreasing weights $(2j+1)^{-1}$
proves convergence and the tail bound
\begin{equation}\label{polyseries:signtail}
 \left|\sum_{j=m+1}^\infty
        \frac{\sin((2j+1)\phi)}{2j+1}\right|
 \le\frac1{(2m+3)|\sin\phi|}.
\end{equation}
For completeness the sum's value can be identified without a Fourier
convergence theorem. For $0\le r<1$, termwise differentiation of a
power series and a geometric sum show that
\[
 H(r)=\sum_{j\ge0}\frac{r^{2j+1}\sin((2j+1)\phi)}{2j+1},\qquad
 H'(r)=\frac{(1+r^2)\sin\phi}
 {(1-r^2)^2+4r^2\sin^2\phi}.
\]
Both $H(r)$ and
$\tfrac12\arctan(2r\sin\phi/(1-r^2))$ vanish at zero and have this
derivative. Their equality implies $H(r)\to(\pi/4)\operatorname{sign}(\phi)$
as $r\uparrow1$. Summation by parts gives the bound
\eqref{polyseries:signtail} also for the tails with weights
$r^{2j+1}/(2j+1)$, uniformly in $0\le r<1$. First choose a large
finite cutoff, then take $r\uparrow1$, and finally let the cutoff
increase: the ordinary sum equals this limit. Thus
\[
 \operatorname{sign}(x)=\frac4\pi\sum_{j\ge0}
            \frac{(-1)^j}{2j+1}T_{2j+1}(x).
\]
At $x=0$ every summand is zero, agreeing with our convention. On
$\delta\le|x|\le1$, $\delta>0$, the error after $j=m$ is at most
$4/[\pi(2m+3)\delta]$. There is no uniform convergence across zero.

\subsection{Exponential coefficients from a double power series}
\label{polyseries:bessel}
For real $a$ and $k\ge0$ define
\[
 I_k(a)=\frac1\pi\int_0^\pi e^{a\cos\theta}\cos(k\theta)\,d\theta.
\]
A scalar exponential power series, rather than a special-function API,
gives every required property. Put $b=a/2$ and $z=e^{i\theta}$. Then
\[
 e^{a\cos\theta}=e^{bz}e^{bz^{-1}}
 =\sum_{p,q\ge0}\frac{b^{p+q}}{p!q!}z^{p-q}.
\]
The double series is absolutely and uniformly convergent on $|z|=1$,
because its absolute majorant sums to $e^{2|b|}=e^{|a|}$.
Partition the pairs $(p,q)$ by $p-q$ and pair opposite differences.
This gives
\begin{equation}\label{polyseries:besselseries}
 I_k(a)=\sum_{j=0}^\infty\frac{(a/2)^{2j+k}}{j!(j+k)!},\qquad
 e^{ax}=I_0(a)+2\sum_{k=1}^\infty I_k(a)T_k(x)
 \quad(-1\le x\le1).
\end{equation}
The integral definition agrees with the displayed power series by
termwise integration and cosine orthogonality; this also verifies the
different normalization at $k=0$. Uniform convergence follows from the
same absolute majorant, independently of any preexisting Bessel identity.
For $a\ge0$, each term is nonnegative and $(j+k)!\ge k!$, whence
\begin{equation}\label{polyseries:besselbound}
 0\le I_k(a)\le\frac{(a/2)^k}{k!}e^{a^2/4}.
\end{equation}
For negative $a$, the power series gives $I_k(a)=(-1)^kI_k(|a|)$.
Evaluation at $x=1$ with $a\ge0$ also gives the useful exact majorant
$I_0(a)+2\sum_{k\ge1}I_k(a)=e^a$.
The degree-$n$ truncation therefore satisfies
\[
 \|e^{a\cdot}-P_ne^{a\cdot}\|_\infty
 \le2\sum_{k>n}|I_k(a)|
 \le2e^{a^2/4+|a|/2}\frac{(|a|/2)^{n+1}}{(n+1)!}.
\]
For the last estimate use $(n+1+j)!\ge(n+1)!j!$ and sum an
exponential series. The case $a=0$ has $I_0=1$ and $I_k=0$ for $k>0$,
so all statements include it without a limiting argument.

\subsection{Fej\'er means as positive contractions}
\label{polyseries:fejer}
Define
\[
 S_nf=\frac1{n+1}\sum_{m=0}^nP_mf
 =\sum_{k=0}^n\left(1-\frac{k}{n+1}\right)a_k(f)T_k.
\]
Only a finite sum identity is needed for its kernel. Set
\begin{equation}\label{polyseries:fejerkernel}
 K_n(t)=\frac1{n+1}\left|\sum_{j=0}^ne^{ijt}\right|^2
 =1+2\sum_{k=1}^n\left(1-\frac{k}{n+1}\right)\cos(kt).
\end{equation}
The identity follows by counting the pairs $(j,\ell)$ with $j-\ell=k$.
Consequently $K_n\ge0$ and $(2\pi)^{-1}\int_{-\pi}^\pi K_n=1$.
For $g(\theta)=f(\cos\theta)$, a change of periodic integration variable
and cosine addition give
\begin{equation}\label{polyseries:fejeroperator}
 S_nf(\cos\theta)=\frac1{2\pi}\int_{-\pi}^\pi
                K_n(t)g(\theta-t)\,dt.
\end{equation}
Every interchange here is finite. Positivity, preservation of constants,
and $\|S_nf\|_\infty\le\|f\|_\infty$ follow immediately. In particular
the operator norm is exactly one, because $S_n1=1$.

Uniform convergence now follows from a general elementary principle:
uniformly bounded operators that converge on a dense linear subspace
converge everywhere. Explicitly, if $p=\sum_{k=0}^da_k(p)T_k$ and $n\ge d$,
\[
 \|S_np-p\|_\infty
 \le\frac1{n+1}\sum_{k=1}^d k|a_k(p)|\longrightarrow0.
\]
For continuous $f$, choose a polynomial $p$ with
$\|f-p\|_\infty<\varepsilon$. The contraction property gives
\[
 \|S_nf-f\|_\infty
 \le2\varepsilon+\|S_np-p\|_\infty.
\]
Taking $n\to\infty$ and then $\varepsilon\downarrow0$ proves the full
Fej\'er assertion. No kernel concentration estimates or Fourier inversion
are needed. This proof is especially suitable when polynomial density is
already available.

\subsection{Finite-dimensional operator consequences and inspected reuse}
\label{polyseries:reuse}
For a real symmetric or complex Hermitian $A$ with spectrum in $[-1,1]$,
orthogonal or unitary diagonalization transfers each uniform scalar bound:
\[
 \|f(A)-p(A)\|_2=\max_{\lambda\in\operatorname{spec}(A)}
                       |f(\lambda)-p(\lambda)|
 \le\|f-p\|_{\infty,[-1,1]}.
\]
In the empty dimension the left side is zero; use the inequality without
an empty maximum. Thus the finite coefficient tails directly bound
matrix-function and matrix--vector errors. Fej\'er positivity gives
$S_nf(A)\succeq0$ whenever $f\ge0$ on the interval. The sign tail gives
an explicit operator bound for a spectral gap $|\lambda|\ge\delta>0$.
These are applications of an existing spectral decomposition, not new
infinite-dimensional functional calculus obligations.

The following reuse was checked by source/declaration inspection; it is
not a claim that the present derivations have been formalized.
\begin{itemize}
\item In Mathlib's
\code{Analysis/SpecialFunctions/Trigonometric/Chebyshev/Orthogonality.lean},
\code{Polynomial.Chebyshev.integral_measureT_eq_integral_cos},
\code{integrable_measureT},
\code{integral_eval_T_real_mul_eval_T_real_measureT_of_ne},
\code{integral_eval_T_real_mul_self_measureT_zero}, and
\code{integral_T_real_mul_self_measureT_of_ne_zero}
already provide the weighted integral and Gram matrix. The final four
names are in the same \code{Polynomial.Chebyshev} namespace.
\item \code{Polynomial.Chebyshev.chebyshevTsequence} and
\code{Polynomial.Sequence.span_degreeLT} provide the finite polynomial
basis argument. Mathlib's
\code{polynomialFunctions.topologicalClosure} supplies polynomial density
on compact real sets; the interval-specific density route can be used as
well. This is enough for Lemma~\ref{polyseries:uniqueness} and the positive
operator density argument.
\item \code{Orthonormal.sum_inner_products_le} is in
\code{Analysis/InnerProductSpace/Orthonormal.lean}. For
\eqref{polyseries:finitebessel}, expanding an integrated square may be
shorter than building an $L^2$ orthonormal family. The required trigonometric
Gram matrix is finite in either route.
\item
\code{AbsolutelyContinuousOnInterval.integral_mul_deriv_eq_deriv_mul}
is in Mathlib's
\code{MeasureTheory/Integral/IntervalIntegral/AbsolutelyContinuousFun.lean}.
The variation API, including \code{eVariationOn.Icc_add_Icc}, supports the
step approximation in Lemma~\ref{polyseries:bvbase}; no missing Stieltjes
integration API need be filled.
\item At the pinned public snapshot, NLAlib already has
\code{NLAlib.chebyshevAliasIndex},
\code{NLAlib.chebyshevAliasIndex_le},
\code{NLAlib.eval_T_node_eq_eval_T_alias},
\code{NLAlib.chebyshev_node_injOn},
\code{NLAlib.interpolate_eval_T_node_eq_T_alias}, and
\code{NLAlib.abs_eval_interpolate_eval_T_node_le_one}.
The last two already supply the interpolation identity and its norm bound,
as verified in the pinned source file. Reuse these for the finite fold and
extend only the summable-coefficient argument. The analytic interpolation
module already contains an analytic special case of the same sum/interpolation
bridge; its analytic assumptions must not be silently transferred to a
general continuous function.
\item The pinned declarations also contain
\code{NLAlib.realChebyshevCircleFunction},
\code{NLAlib.realChebyshevCircleFunction_coe}, and
\code{NLAlib.continuous_realChebyshevCircleFunction}; the pinned source
also proves \code{NLAlib.even_realChebyshevCircleFunction}. These avoid defining
a second real circle lift for \eqref{polyseries:fejeroperator}.
\code{NLAlib.sum_choose_mul_chebyshev_eq_two_mul_X_pow} gives an alternative
entirely real finite-polynomial entry point to the exponential expansion:
expand each exponential monomial in that existing Chebyshev identity,
then regroup the absolutely summable coefficients.
\end{itemize}
The source labels for the target statements are ATAP Theorems 3.1,
4.1--4.2, 7.1--7.2 and Exercises 3.6, 3.15, 6.6; the author's
\href{https://people.maths.ox.ac.uk/trefethen/ATAP/}{official book page}
and \href{https://people.maths.ox.ac.uk/trefethen/trefethen_sample.pdf}{sample
chapters} provide the source context. The arguments above give the
individual reductions explicitly instead of importing those convergence
theorems as black boxes.

% END series.tex


% BEGIN best.tex
\section{Best approximation as finite-dimensional optimization}
\label{polybest:main}

Throughout this section $I=[a,b]$ with $a<b$, $C(I)$ means real continuous
functions with the supremum norm, and $\mathcal P_n$ consists of real
polynomials of degree at most $n$, restricted to $I$.  The zero polynomial
belongs to every $\mathcal P_n$.  No construction below needs a measure on
the underlying matrix spectrum: after proving a scalar bound, the spectral
theorem transfers it to a self-adjoint operator with spectrum in $I$.

\subsection{Attainment using an interpolation-coordinate box}
\label{polybest:minimax}
\textbf{Target: \code{minimax-error-def}. Complete derivation.}
Define $E_n(f)=\inf_{p\in\mathcal P_n}\norm{f-p}_\infty$ for $f\in C(I)$.
This equals the infimum of all real $E$ for which some $p\in\mathcal P_n$
satisfies $|f(x)-p(x)|\le E$ for every $x\in I$.

\begin{theorem}
The infimum is attained.  Moreover
\[
 0\le E_{n+1}(f)\le E_n(f)\le\norm f_\infty,
 \qquad E_n(f)\longrightarrow0,
\]
and, for $p\in\mathcal P_n$ and $c\in\R$,
\[
 E_n(f+p)=E_n(f),\qquad E_n(cf)=|c|E_n(f),\qquad
 E_n(f+g)\le E_n(f)+E_n(g).
\]
If $\varphi(t)=((b-a)t+a+b)/2$, then
$E_n(f;[a,b])=E_n(f\circ\varphi;[-1,1])$.
\end{theorem}
\begin{proof}
Fix $n+1$ distinct points $t_0,\ldots,t_n$ in $I$, and their Lagrange
polynomials $\ell_j$.  Put $M=\norm f_\infty$.  For $y\in\R^{n+1}$ set
$P_y=\sum_jy_j\ell_j$.  The map
$y\mapsto\norm{f-P_y}_\infty$ is continuous, since its change between
$y$ and $z$ is at most
$\sum_j|y_j-z_j|\norm{\ell_j}_\infty$.
It therefore attains a minimum on the compact box $[-2M,2M]^{n+1}$.
Every polynomial with error at most $M$ lies in this box of sample values,
because $|p(t_j)|\le2M$; every polynomial outside the box has error larger
than $M$.  The zero polynomial has error $M$ and lies in the box.  Hence
the minimum on the box is the global minimum.  This also covers $M=0$.

Nestedness of the polynomial spaces and the competitor zero give the first
inequalities.  Weierstrass approximation gives, for every $\varepsilon>0$,
one polynomial $p$ with $\norm{f-p}_\infty<\varepsilon$; all $n\ge\deg p$
then have $E_n(f)<\varepsilon$.  Translation by $p$ is a bijection of
$\mathcal P_n$, as is multiplication by $c$ when $c\ne0$; $c=0$ is
immediate.  Add two minimizing polynomials for the subadditivity assertion.
Finally composition with $\varphi$ and with its affine inverse are inverse
degree-preserving maps, and preserve the relevant supremum norms.
\end{proof}

\textbf{Verified reuse.} Mathlib has
\code{Lagrange.eq_interpolate}, \code{Lagrange.degree_interpolate_lt},
\code{IsCompact.exists_isMinOn}, and
\code{exists_polynomial_near_of_continuousOn}.  Thus the only new existence
argument is a compact finite box; no nearest-point theorem in a Banach
space is required.  The indexed declaration
\code{NLAlib.isLeast_chebyshev_minimax} already uses the same
``least admissible error'' statement style, for the constrained Chebyshev
problem.

\subsection{Interpolation is an explicit bounded projection}
\label{polybest:lebesgue}
\textbf{Target: \code{lebesgue-constant-def}. Complete derivation.}
For distinct $v_0,\ldots,v_n\in I$, let
\[
 (Lf)(x)=\sum_{j=0}^n f(v_j)\ell_j(x),\qquad
 \lambda(x)=\sum_{j=0}^n|\ell_j(x)|,\qquad
 \Lambda=\max_{x\in I}\lambda(x).
\]
The maximum exists by continuity.  The identity $\sum_j\ell_j=1$ gives
$\Lambda\ge1$ and
\[
 |Lf(x)|\le\lambda(x)\max_j|f(v_j)|\le\Lambda\norm f_\infty.
\]
This bound is exact: choose $x_*$ maximizing $\lambda$, prescribe
$f(v_j)=\operatorname{sgn}(\ell_j(x_*))$, assigning either sign if the
coefficient is zero, and extend these finitely many values by linear
interpolation in the order of the nodes, constantly outside the extreme
nodes.  Then $f$ is continuous, $\norm f_\infty=1$, and
$Lf(x_*)=\Lambda$.  Thus $\norm L=\Lambda$.  The same argument with any
fixed $x$ proves sharpness of the pointwise estimate.

Here is the operator lemma that should be shared by interpolation and
Chebyshev truncation.  If $L:C(I)\to\mathcal P_n$ is linear, fixes
$\mathcal P_n$, and $\norm{Lg}_\infty\le B\norm g_\infty$, then
\begin{equation}\label{polybest:nearbest}
 \norm{f-Lf}_\infty\le(1+B)E_n(f).
\end{equation}
Indeed, for any $p\in\mathcal P_n$,
$f-Lf=(f-p)-L(f-p)$; use a minimizing $p$ from the previous subsection.
Taking $B=\Lambda$ proves the target.  This proof also works with an
infimum and arbitrarily near minimizing polynomials.

\textbf{Verified reuse.} The needed finite identities are already
\code{Lagrange.sum_basis}, \code{Lagrange.eval_interpolate_at_node}, and
\code{Lagrange.eq_interpolate}.  The converse norm inequality needs only
the explicit piecewise affine extension above, not an extension theorem.

\subsection{Alternation certificates without counting intermediate roots}
\label{polybest:dlvp}
\textbf{Target: \code{de-la-vallee-poussin}. Complete derivation.}
For $x_0<\cdots<x_{N+1}$ define
\[
 w_i=\frac1{\prod_{j\ne i}(x_i-x_j)},\qquad
 a_i=\frac{|w_i|}{\sum_j|w_j|}>0.
\]
Taking the coefficient of $X^{N+1}$ in the Lagrange interpolation formula
gives the finite linear-algebra identity
\begin{equation}\label{polybest:annihilator}
 \sum_{i=0}^{N+1}w_i h(x_i)=0\quad(h\in\mathcal P_N),
 \qquad \operatorname{sgn}w_i=(-1)^{N+1-i}.
\end{equation}
Consequently $\sum_i a_i\sigma(-1)^i h(x_i)=0$ for either
$\sigma\in\{1,-1\}$.  This is a signed positive quadrature formula
annihilating the approximation space.

Suppose $p\in\mathcal P_N$ and
$f(x_i)-p(x_i)=\sigma(-1)^i e_i$ with $e_i>0$.
For every $q\in\mathcal P_N$,
\begin{equation}\label{polybest:weightedcertificate}
 \max_i|f(x_i)-q(x_i)|
 \ \ge\ \sum_i a_i\sigma(-1)^i(f(x_i)-q(x_i))
 \ =\ \sum_i a_i e_i
 \ \ge\ \min_i e_i.
\end{equation}
This proves the catalogue bound, and also a sometimes stronger weighted
bound.  No continuity of $f$ is needed for this discrete assertion.
For continuous $f$ it implies $E_N(f)\ge\min_i e_i$.

\textbf{Verified reuse.} Mathlib's \code{Lagrange.coeff_eq_sum} and
\code{Lagrange.leadingCoeff_basis} give
\eqref{polybest:annihilator} directly.  This replaces the alternative
intermediate-value/root-count proof by one finite sum.

\subsection{The converse: separate the compact extremal sign sets}
\label{polybest:equioscillation}
\textbf{Target: \code{equioscillation}. Complete derivation.}

\begin{lemma}[A polynomial matching finitely many extremal sign blocks]
\label{polybest:signblocks}
Let $e\in C(I)$, $E=\norm e_\infty>0$, and
$K_\pm=\{x\in I:e(x)=\pm E\}$.  If their union contains no increasing
alternating list of length $N+2$, there exists $h\in\mathcal P_N$ such
that $e(x)h(x)>0$ on $K_+\cup K_-$.
\end{lemma}
\begin{proof}
The two sets are compact and disjoint, so, when both are nonempty, their
distance is some positive number.  Start at the smallest point of their
union.  After a selected point, select the smallest later point of the
opposite sign, if one exists.  Compactness gives each such smallest point;
positive separation makes the selection stop after finitely many steps.
The hypothesis makes the number of selected points at most $N+1$.

Between successive selected points $u<v$, all extremal points in $[u,v)$
have the sign at $u$.  The supremum of those points is strictly below $v$,
by separation from the opposite sign.  Choose a separator in that open
gap.  These separators partition the extremal set into nonempty,
successive constant-sign blocks whose signs alternate.  A product of
$X-z$ over the separators has exactly these sign changes.  Multiply it
by $1$ or $-1$ so its first sign agrees with $e$.  Its degree, the number
of separators, is at most $N$.  If there is only one sign block, use a
constant $1$ or $-1$.
\end{proof}

\begin{lemma}[Uniform improvement from a signed direction]
\label{polybest:improvement}
Let $E=\norm e_\infty>0$.  If $v_t\to v$ uniformly as $t\downarrow0$,
all these functions are continuous, and $e(x)v(x)>0$ whenever
$|e(x)|=E$, then $\norm{e-tv_t}_\infty<E$ for all sufficiently small
$t>0$.
\end{lemma}
\begin{proof}
On the compact extremal set $|v|$ has a positive minimum $c$.
Choose an open neighborhood $U$ of that set on which $e,v$ have the
same sign, $|v|>c/2$, and $|e|>E/2$.  Uniform convergence ensures the
same sign and $|v_t|>c/4$ on a possibly smaller such neighborhood,
and a common finite bound $B$ for $|v_t|$.  More explicitly, first
shrink $U$ so that $|v|\ge3c/4$ there, then require
$\norm{v_t-v}_\infty<c/4$.  On $I\setminus U$, compactness gives
$|e|\le E-\eta$ for some $\eta>0$; if this set is empty omit this
condition.  Take $tB<\min(E/2,\eta/2)$.  On $U$ there is no sign
overshoot and $|e-tv_t|\le E-tc/2$; off $U$ the bound is at most
$E-\eta/2$.
\end{proof}

\begin{theorem}[Polynomial equioscillation and uniqueness]
For $p\in\mathcal P_n$ and $f\in C(I)$, $p$ is best if and only if
there are $a\le x_0<\cdots<x_{n+1}\le b$ and $\sigma\in\{1,-1\}$
such that
\[
 f(x_i)-p(x_i)=\sigma(-1)^i\norm{f-p}_\infty.
\]
The best polynomial is unique.
\end{theorem}
\begin{proof}
When the error is zero, every list of $n+2$ distinct points works.
For positive error, alternation implies optimality by
\eqref{polybest:weightedcertificate}.  Conversely, if a best $p$ did not
alternate, Lemma~\ref{polybest:signblocks} would supply a polynomial
$h\in\mathcal P_n$ with the correct extremal signs.  Apply
Lemma~\ref{polybest:improvement} with $e=f-p$ and $v_t=v=h$ to improve
$p$ to $p+th$, a contradiction.

For uniqueness, let $p$ be a best polynomial with a positive error and
an alternating list, and let $q$ be another best polynomial.  At these
points $\sigma(-1)^i(q-p)(x_i)\ge0$.  Their positive weighted sum is
zero by \eqref{polybest:annihilator}; hence all the summands vanish.
The polynomial $q-p$ of degree at most $n$ vanishes at $n+2$ distinct
points, so is zero.  If the best error is zero, both polynomials equal
$f$ on $I$ and the same conclusion follows.
\end{proof}

\textbf{Formalization split.} The substantial new topological lemma is
the compact sign-block construction, shared with rational approximation
below.  The remainder is a finite annihilator, compact extrema, and one
uniform perturbation estimate.  Mathlib's root-counting or
\code{Lagrange.eq_zero_of_degree_lt_of_eval_finset_eq_zero} handles the
last uniqueness step.

\subsection{Every best polynomial is an interpolant}
\label{polybest:bestinterpolates}
\textbf{Target: \code{best-approximant-interpolates}. Complete derivation.}
If $p$ is best on $[-1,1]$ and the error is positive, the $n+2$
alternating extrema from the preceding theorem enclose $n+1$ disjoint
open intervals.  Continuity gives one zero of $f-p$ in each interval.
These are $n+1$ distinct points in $(-1,1)$.  If the error is zero,
choose any $n+1$ distinct interior points.  Lagrange uniqueness now gives
that $p$ is the interpolant at those points.  Reuse
\code{intermediate_value_Icc} and \code{Lagrange.eq_interpolate}.

\section{Logarithmic projection bounds through finite trigonometric kernels}
\label{polybest:kernels}

\subsection{Chebyshev--Lobatto interpolation: the stated sharp coefficient}
\label{polybest:lobatto}
\textbf{Target: \code{lebesgue-constant-chebyshev}. Complete derivation,
including the catalogue constant.}
For $n\ge1$ and nodes $x_j=\cos(j\pi/n)$, we prove
\begin{equation}\label{polybest:lobatto-bound}
 \Lambda_n\le 1+\frac2\pi\log(n+1).
\end{equation}
The route is a finite trigonometric interpolation operator and two scalar
inequalities.  It does not invoke an unproved comparison theorem for
Chebyshev roots and extrema.

Put $h=\pi/(2n)$ and define the continuous $2\pi$-periodic cardinal kernel
\[
 K_n(t)=\frac1{2n}\left(1+2\sum_{k=1}^{n-1}\cos(kt)+\cos(nt)\right)
       =\frac{\sin(nt)\cot(t/2)}{2n}
       \quad(t\notin2\pi\mathbb Z).
\]
The equality is a finite geometric-sum identity.  At the $2n$ equally
spaced angles $t_j=j\pi/n$, its translates satisfy
$K_n(t_i-t_j)=\delta_{ij}$.  For symmetric sampled data
$g(t_j)=g(-t_j)$, the polynomial
$\sum_{j=0}^{2n-1}g(t_j)K_n(t-t_j)$ is even and hence a polynomial of
degree at most $n$ in $\cos t$.  It agrees with the data at every
Lobatto node.  Therefore it is exactly the Lobatto interpolant, and
\begin{equation}\label{polybest:trig-norm-upper}
 \Lambda_n\le\max_t\sum_{j=0}^{2n-1}|K_n(t-t_j)|.
\end{equation}

Translation and reflection symmetry reduce the last expression to
$t=\alpha\pi/n$, $0\le\alpha\le1$.  At the endpoints it equals one.
For $0<\alpha<1$ it equals
\begin{equation}\label{polybest:cot-pair-sum}
 \frac{\sin(\pi\alpha)}{2n}
 \sum_{j=0}^{n-1}\left\{
  \cot(h(j+\alpha))+\cot(h(j+1-\alpha))\right\}.
\end{equation}
All the cotangents here have arguments in $(0,\pi/2)$.

We claim that each brace, after multiplication by $\sin(\pi\alpha)$,
is at most $2\cot(h(j+1/2))$.  Set $s=|\alpha-1/2|\le1/2$,
$u=h(j+1/2)$ and $v=hs$.  For $s<1/2$, the identity
\[
 \cot(u+v)+\cot(u-v)
 =\frac{\sin(2u)}{\sin^2u-\sin^2v}
\]
reduces the claim to
\begin{equation}\label{polybest:paired-scalar}
 \sin^2(hs)\le2\sin^2(\pi s/2)\sin^2u.
\end{equation}
For $s=0$ this is immediate.  For $s>0$, the function
$s\mapsto\sin(hs)/\sin(\pi s/2)$ is increasing on $(0,1/2]$:
its logarithmic derivative is nonnegative because $z\cot z$ is
decreasing on $(0,\pi/2]$, and $h\le\pi/2$.
The latter assertion follows by differentiating and using
$\sin z\cos z\le z$.
Thus
\[
 \frac{\sin(hs)}{\sin(\pi s/2)}
 \le \frac{\sin(h/2)}{\sin(\pi/4)}
 =\sqrt2\sin(h/2)\le\sqrt2\sin u,
\]
which proves \eqref{polybest:paired-scalar}.  Consequently the full
trigonometric Lebesgue constant is attained at a half-grid point and is
\begin{equation}\label{polybest:cot-exact}
 C_n=\frac1n\sum_{j=0}^{n-1}\cot\left(\frac{(j+1/2)\pi}{2n}\right),
 \qquad \Lambda_n\le C_n.
\end{equation}
The endpoint value one is no larger than $C_n$, since the cardinal
functions sum to one.  This also covers $n=1$, where $C_1=1$.

It remains to bound this finite sum.  The function $\cot u$ is convex
on $(0,\pi/2]$.  Integrate the midpoint inequality on each interval
$[jh,(j+1)h]$, $j=1,\ldots,n-1$, to obtain
\[
 C_n\le\frac1n\cot(h/2)+\frac1{nh}\int_h^{\pi/2}\cot u\,du
 \le\frac4\pi+\frac2\pi\log\frac1{\sin h}.
\]
Here $\cot(h/2)\le2/h$.  Also
\[
 \sin h\ge h-h^3/6\ge\frac{\pi}{2(n+1)}.
\]
Indeed, the second inequality is $\pi^2(n+1)\le24n^2$, which follows
from $\pi^2<12$ and $n+1\le2n^2$.  The first follows, for example,
by differentiating twice the inequality
$\sin u-u+u^3/6\ge0$ and using $\sin u\le u$.
Therefore
\[
 C_n\le\frac2\pi\log(n+1)+\frac2\pi\left(2+\log\frac2\pi\right)
       \le\frac2\pi\log(n+1)+1.
\]
For completeness, the last numerical inequality needs no decimal
oracle: $\log x\ge2(x-1)/(x+1)$ for $x\ge1$, by differentiation.
At $x=\pi/2$, this lower bound is at least $2-\pi/2$ because
$\pi^2+2\pi-16>0$, which follows already from $\pi>25/8$.
This proves \eqref{polybest:lobatto-bound}.  Combining it with
\eqref{polybest:nearbest} gives exactly
\[
 \norm{f-I_nf}_\infty\le
 \left(2+\frac2\pi\log(n+1)\right)E_n(f).
\]

\textbf{Reuse and scope.} The pinned NLAlib index includes
\code{NLAlib.chebyshev_node_injOn} and the polynomial interpolation APIs;
Mathlib has \code{Polynomial.Chebyshev.T_real_cos} and
\code{Real.pi_gt_d2}.  The additional scalar calculus facts have been
proved above, so there is no hidden Ehlich--Zeller hypothesis.  The
classical source is Ehlich and Zeller, \emph{Auswertung der Normen von
Interpolationsoperatoren}, Math. Ann. 164 (1966), 105--112,
\url{https://doi.org/10.1007/BF01429047};
the derivation here uses only the displayed finite kernel identities.

\subsection{Chebyshev truncation: exact norm and an elementary logarithmic bound}
\label{polybest:projection}
\textbf{Target: \code{chebyshev-projection-lebesgue-constant}.
Complete derivation.}
Let $S_nf=\sum_{k=0}^na_k(f)T_k$ with the usual cosine-integral
coefficients, and let
\[
 D_n(t)=1+2\sum_{k=1}^n\cos(kt)
       =\frac{\sin((n+1/2)t)}{\sin(t/2)},\qquad
 L_n=\frac1{2\pi}\int_{-\pi}^{\pi}|D_n(t)|\,dt.
\]
The trigonometric sum defines the values at multiples of $2\pi$.
For the even periodic function $g(t)=f(\cos t)$, a finite product-to-sum
calculation gives
\[
 (S_nf)(\cos\theta)=\frac1{2\pi}\int_{-\pi}^{\pi}
                  D_n(\theta-t)g(t)\,dt.
\]
Thus $\norm{S_n}\le L_n$.  To prove equality within the even subspace,
set $d_n(x)=1+2\sum_{k=1}^nT_k(x)$ and
$f_\varepsilon(x)=d_n(x)/\sqrt{d_n(x)^2+\varepsilon^2}$.
It is continuous and bounded by one.  At $\theta=0$ the preceding
integral differs from $L_n$ by at most $\varepsilon$, since
\[
 0\le|u|-\frac{u^2}{\sqrt{u^2+\varepsilon^2}}\le\varepsilon.
\]
Letting $\varepsilon\downarrow0$ proves $\norm{S_n}=L_n$ without a
discontinuous sign function.

Here is an explicit elementary estimate with the precise catalogue
constant.  For $0<u\le\pi/2<2$,
\[
 \sin u\ge u-u^3/6\ge\frac{u}{1+u^2/2},
 \qquad
 \frac1{\sin(t/2)}\le\frac2t+\frac t4\quad(0<t\le\pi).
\]
The middle inequality follows by multiplication: the surplus is
$u^3(4-u^2)/12\ge0$.  Write $M=n+1/2$ and $H_n=\sum_{j=1}^n1/j$.
Then
\begin{align*}
 L_n
 &\le\frac2\pi\int_0^{M\pi}\frac{|\sin u|}{u}\,du+\frac\pi8,\\
 \int_0^{M\pi}\frac{|\sin u|}{u}\,du
 &\le\pi+\sum_{j=1}^{n}\frac1{j\pi}
                       \int_{j\pi}^{(j+1)\pi}|\sin u|\,du
 =\pi+\frac2\pi H_n.
\end{align*}
The first interval uses $|\sin u|/u\le1$; in the last interval extending
the integral only increases it.  Since $H_n\le1+\log(n+1)$ for $n\ge1$,
\[
 L_n\le\frac4{\pi^2}\log(n+1)+2+\frac\pi8+\frac4{\pi^2}
       \le\frac4{\pi^2}\log(n+1)+3.
\]
The final inequality follows from $3<\pi<22/7$, since
$\pi/8+4/\pi^2<11/28+4/9<1$.
Cosine orthogonality proves that $S_n$ fixes $\mathcal P_n$.
Equation~\eqref{polybest:nearbest} now yields
\[
 \norm{f-S_nf}_\infty\le
 \left(4+\frac4{\pi^2}\log(n+1)\right)E_n(f).
\]
For $n=0$ the operator is the normalized cosine average and has norm one.

\textbf{Verified reuse.} The real even circle lift is already
\code{NLAlib.realChebyshevCircleFunction}, with
\code{NLAlib.continuous_realChebyshevCircleFunction} and
\code{NLAlib.even_realChebyshevCircleFunction}.  Chebyshev orthogonality
is in Mathlib's Chebyshev orthogonality file.  Only finite trigonometric
sums and ordinary interval integration are needed for this bound.

\subsection{Interpolation lower bounds: exact elementary part and the sharp boundary}
\label{polybest:lower}
\textbf{Target: \code{lebesgue-constant-lower-bound}. The equispaced
inequality is completely proved below. The exact universal additive
constant is now fully derived in Section~\ref{polyretry:lebesgue},
including the arbitrary-grid comparison. The weaker operator estimate
below remains useful as a shorter independent bound.}

The catalogue's sharp universal assertion is
\begin{equation}\label{polybest:brutman-import}
 \Lambda_n\ge\frac2\pi\left(\log(n+1)+\gamma+\log\frac4\pi\right)
 \quad\text{for every set of $n+1$ distinct nodes in $[-1,1]$}.
\end{equation}
This is the Erd\H{o}s--Brutman inequality, with Euler's constant $\gamma$.
It is a genuine sharp interpolation result; a proof of the projection
norm identity alone does not establish it.  A primary reference is
L.~Brutman, \emph{On the Lebesgue Function for Polynomial Interpolation},
SIAM J. Numer. Anal. 15 (1978), 694--704,
\url{https://doi.org/10.1137/0715046}.  A complete proof of
\eqref{polybest:brutman-import} is given in the second-pass sections
below. The following weaker universal logarithmic estimate is a shorter
self-contained operator alternative:
\begin{equation}\label{polybest:weaker-lower}
 \Lambda_n\ge\max\left(1,\frac2{\pi^2}\log(n+1)-\frac12\right).
\end{equation}

To prove it, transfer the interpolation operator to even continuous
functions on the circle and precompose with evenization
$Eh(t)=(h(t)+h(-t))/2$.  Denote the resulting operator by $Q$.
It has norm at most $\Lambda_n$, takes values in the span of
$1,\cos t,\ldots,\cos nt$, and fixes each of these cosine functions.
Let $R_sh(t)=h(t+s)$ and average its rotations:
\[
 A h=\frac1{2\pi}\int_{-\pi}^{\pi}R_{-s}QR_sh\,ds.
\]
This integral is elementary: $Q$ is a finite sum of evaluations times
trigonometric polynomials, so each value is an ordinary scalar integral;
the result is again a trigonometric polynomial of degree at most $n$.
The triangle inequality gives $\norm A\le\Lambda_n$.

For full detail of the multiplier calculation, on $e^{ikt}$ the
evenization gives $\cos(kt)$, whose image under $Q$ is a finite cosine
sum.  Averaging $e^{iks}\cos(j(t-s))$ kills every term unless $j=|k|$.
If $1\le|k|\le n$, reproduction makes that surviving coefficient
one and the average is $e^{ikt}/2$; if $|k|>n$ it is zero.  The constant
is fixed.  Hence $2A-M$ is the ordinary Fourier partial-sum operator,
where $Mh$ is the mean constant.  This calculation can equivalently be
made using sine and cosine pairs, without complexification.  Its extension
from finite trigonometric polynomials to all continuous functions follows
from their uniform density and boundedness of both operators.  Uniform
density follows from the real Stone--Weierstrass theorem applied to
$\cos t,\sin t$, which separate points on the circle.  Alternatively,
substitute the finite evaluation formula for $Q$ into the integral and
perform the same finite product-to-sum calculation directly.
The norm of the Fourier partial sum is $L_n$, with the lower bound
obtained by the continuous sign approximants used above.  Thus
$L_n\le2\Lambda_n+1$.

Since $\sin(t/2)\le t/2$,
\[
 L_n\ge\frac2\pi\int_0^{(n+1/2)\pi}\frac{|\sin u|}{u}\,du
 \ge\frac4{\pi^2}\sum_{j=1}^n\frac1j
 \ge\frac4{\pi^2}\log(n+1).
\]
The middle inequality integrates over the $n$ full intervals
$[(j-1)\pi,j\pi]$ and bounds $1/u$ below by $1/(j\pi)$.
Together with $\Lambda_n\ge1$, this proves
\eqref{polybest:weaker-lower}.  This weaker coefficient is explicitly
different from the sharp coefficient $2/\pi$ in
\eqref{polybest:brutman-import}.

For equispaced nodes, however, the exact requested bound has a short
finite proof.  Affine invariance reduces the nodes to $0,1,\ldots,n$.
At $t=1/2$,
\[
 \sum_{j=0}^n|\ell_j(1/2)|
 =\frac{|\prod_{k=0}^n(1/2-k)|}{n!}
       \sum_{j=0}^n\frac{\binom nj}{|j-1/2|}.
\]
For $n\ge1$,
\[
 \frac{|\prod_{k=0}^n(1/2-k)|}{n!}
 =\frac12\prod_{k=1}^n\left(1-\frac1{2k}\right)
 \ge\frac1{4n},
\]
because the first factor is $1/2$ and the factors for $k\ge2$ are
at least $(k-1)/k$.  Also $|j-1/2|<n$ for every $0\le j\le n$.
The binomial sum is $2^n$, so
\[
 \Lambda_n>\frac1{4n}\frac{2^n}{n}
           =\frac{2^{n-2}}{n^2}.
\]
This includes $n=1$ and retains the strict inequality.  It is the
\emph{sum} of cardinal magnitudes that gives exponential growth;
the single endpoint cardinal polynomial does not suffice.
The primary source for the equispaced phenomenon is Trefethen and
Weideman (1991),
\url{https://people.maths.ox.ac.uk/~trefethen/publication/PDF/1991_49.pdf}.

\section{Rational alternation through one polynomial linear map}
\label{polybest:rational}
\textbf{Target: \code{rational-equioscillation}. Complete derivation,
with the zero-function convention made explicit.}

Let $\mathcal R_{mn}$ be the real rational functions $p/q$ with
$\deg p\le m$, $\deg q\le n$, and no pole on $I=[-1,1]$.
Always cancel common factors and choose $q>0$ on $I$.
For a nonzero reduced fraction with exact degrees $(\mu,\nu)$ set
\[
 d=\min(m-\mu,n-\nu),\qquad N=m+n-d=\max(m+\nu,n+\mu).
\]
For the zero function use the reduced representation $0/1$ and
\emph{define $N=m$}, equivalently $d=n$ using $\deg0=-\infty$.
Using $\mu=0$ for the zero numerator would in general give the wrong
alternation count.  This case should be a separate branch in a formal
statement.

\begin{theorem}[Rational best approximation]
Every real $f\in C(I)$ has a unique best function in $\mathcal R_{mn}$.
A given $r\in\mathcal R_{mn}$ is best if and only if its error
$e=f-r$ has $N+2$ increasing extremal points with alternating signs:
\[
 e(x_i)=\sigma(-1)^i\norm e_\infty,\qquad 0\le i\le N+1.
\]
If, more generally, $e(x_i)$ has these alternating signs and
$|e(x_i)|\ge\varepsilon>0$, then every $s\in\mathcal R_{mn}$ has
uniform error at least $\varepsilon$.
\end{theorem}

\paragraph{Existence: compact coefficients followed by cancellation.}
Choose a minimizing sequence $p_j/q_j$, orienting $q_j>0$ on $I$,
and normalize the Euclidean norm of its coefficient vector to one.
The corresponding numerators are multiplied by the same factor.
Pass to a subsequence for which $q_j$ converges coefficientwise to $q$,
with coefficient norm one.  Thus $q\not\equiv0$ and $q\ge0$ on $I$.
The errors of a minimizing sequence may be chosen at most
$\norm f_\infty+1$, so the rational functions are uniformly bounded
by $2\norm f_\infty+1$.  The normalized denominators are uniformly
bounded by $\sqrt{n+1}$ on $I$; hence the numerators are uniformly
bounded too.  Evaluating at $m+1$ fixed distinct points and applying the
inverse Vandermonde map makes their coefficient vectors bounded.
Pass to a further subsequence $p_j\to p$ coefficientwise.

If $E$ is the infimum error, passage to the limit in the pointwise
inequality gives
\begin{equation}\label{polybest:rational-limit}
 |f(x)q(x)-p(x)|\le E q(x),\qquad x\in I.
\end{equation}
In particular $|p(x)|\le(\norm f_\infty+E)|q(x)|$.
At a zero $z$ of $q$ in $I$, $p$ has at least the same multiplicity:
otherwise factor $p=(X-z)^u p_0$, $q=(X-z)^v q_0$ with $u<v$ and
$p_0(z)q_0(z)\ne0$; then $|p(x)/q(x)|$ is unbounded as $x\to z$
through interior points.  The same one-sided argument covers endpoints.
If $p=0$, take the candidate $0/1$ directly.  Otherwise cancel the gcd
of $p,q$.  The reduced denominator has no zero on $I$, and its degrees
have not increased.  At points with $q(x)\ne0$, inequality
\eqref{polybest:rational-limit} says that this fraction has error at
most $E$.  Such points are dense since a nonzero polynomial has only
finitely many zeros; continuity extends the inequality to the canceled
points.  This gives a minimizer.  The same extension proves the claim
when $p=0$.

\paragraph{The finite linear map producing every descent direction.}
For nonzero reduced $r=p/q$, consider
\begin{equation}\label{polybest:rational-linear-map}
 \mathcal L:\mathcal P_m\times\mathcal P_n\longrightarrow\mathcal P_N,
 \qquad \mathcal L(P,Q)=Pq-pQ.
\end{equation}
Its kernel consists exactly of $(ph,qh)$ with $h\in\mathcal P_d$.
Indeed $Pq=pQ$ and coprimality imply $p\mid P$ and $q\mid Q$; after
substitution, cancellation makes the two quotients equal.  Conversely
these pairs lie in the kernel, and the degree restrictions give exactly
$\deg h\le d$.  Thus its rank is
$(m+1)+(n+1)-(d+1)=N+1$, the dimension of its codomain.  It is surjective.
For $r=0$, the same needed statement is simply that the numerator
directions $P\in\mathcal P_m$ give every element of $\mathcal P_m$.

Suppose a positive-error minimizer did not have the required $N+2$
alternating extrema.  Lemma~\ref{polybest:signblocks} gives
$h\in\mathcal P_N$ with the same sign as its error at every extremum.
Use surjectivity to choose $P,Q$ with $Pq-pQ=h$.  For all sufficiently
small $t>0$, compactness and $q>0$ give $q+tQ>0$ on $I$.
The type-$(m,n)$ perturbation satisfies
\[
 r_t-r=\frac{p+tP}{q+tQ}-\frac pq
       =t\frac{h}{q(q+tQ)}.
\]
The factor after $t$ converges uniformly to $h/q^2$, because $q$ is
bounded below away from zero.  It has the needed extremal signs.
Lemma~\ref{polybest:improvement} therefore contradicts optimality.
When $r=0$, just use $r_t=th$ with $h\in\mathcal P_m$.
This proves necessity.  At zero error any $N+2$ distinct points work.

\paragraph{Sufficiency and the rational de la Vall\'ee Poussin bound.}
Let $r=p/q$ have the alternating error values stated in the theorem,
and let $s=P/Q$ be any competitor with positive denominator $Q$.
The difference has numerator
\[
 H=Pq-pQ\in\mathcal P_N,
 \qquad s-r=\frac{H}{Qq}.
\]
For the zero baseline the numerator is $P$ and has degree at most $m=N$.
If $s$ had error strictly less than $\varepsilon$, then
$\sigma(-1)^i(s-r)(x_i)>0$ for every $i$.  Positivity of $Qq$ transfers
these signs to $H$.  The positive signed annihilator
\eqref{polybest:annihilator} cannot annihilate such a list; this is a
contradiction.  Thus every competitor has error at least $\varepsilon$.
Apply this with $\varepsilon=\norm{f-r}_\infty$ for positive error to
obtain sufficiency.  Zero error is automatically best.

\paragraph{Uniqueness without averaging rational functions pointwise.}
Let $p_1/q_1$ and $p_2/q_2$ be two minimizers, each with $q_i>0$ on
$I$, and common best error $E$.  The degree-preserving weighted midpoint
\[
 r=\frac{p_1+p_2}{q_1+q_2}
\]
belongs to $\mathcal R_{mn}$.  At each $x$, its error is the strictly
positive weighted average of the two errors, so is at most $E$ in
absolute value.  It is itself a minimizer.  Reduce this fraction and
let $N$ be its degree parameter above.  At its $N+2$ alternating extrema,
the weighted average can equal $+E$ or $-E$ only if both individual
errors have that same value.  Thus $r_1-r$ vanishes at $N+2$ distinct
points.  Its numerator relative to the reduced representation of $r$
has degree at most $N$, exactly as in the preceding paragraph, hence
is zero.  Similarly $r_2=r$.  If $E=0$, both minimizers already equal
$f$ on $I$ and uniqueness is immediate.

\textbf{Why this is a useful reroute.} All defect bookkeeping is isolated
in the single rank calculation \eqref{polybest:rational-linear-map}.
Existence uses finite coefficient compactness and cancellation of real
roots; necessity reuses the polynomial sign-block lemma; sufficiency
and uniqueness reuse a finite Vandermonde annihilator.  No rational
implicit-function theorem or nonlinear manifold formalism is needed.
Mathlib already supplies \code{RatFunc}, polynomial gcd/coprimality,
finite-dimensional rank--nullity, and Lagrange interpolation.  Those
interfaces still need assembling; this paragraph reports mathematical
reuse, not a checked Lean implementation.

% END best.tex


% BEGIN quadrature.tex
\section{Interpolation, compressed multiplication, and quadrature}
\label{polyquad:section}

\subsection{Hermite's interpolation formula from a finite rational identity}
\label{polyquad:hermite}
\noindent\textbf{Target: \code{hermite-integral-formula}.}
Let $N\ge1$, let $z_0,\ldots,z_{N-1}$ be distinct points of
$D(c,R)$, $R>0$, and let $f$ be holomorphic on a neighbourhood of
$\overline{D(c,R)}$. Put
\[
 \ell(z)=\prod_{j=0}^{N-1}(z-z_j),\qquad
 L_j(z)=\prod_{h\ne j}\frac{z-z_h}{z_j-z_h},\qquad
 p(z)=\sum_j f(z_j)L_j(z).
\]
For $x\in D(c,R)$ the two desired formulas are
\begin{align}
 p(x)&=\frac1{2\pi i}\oint_{|t-c|=R}
 f(t)\frac{\ell(t)-\ell(x)}{\ell(t)(t-x)}\,dt,
 \label{polyquad:hermite-p}\\
 f(x)-p(x)&=\frac{\ell(x)}{2\pi i}\oint_{|t-c|=R}
 \frac{f(t)}{\ell(t)(t-x)}\,dt.
 \label{polyquad:hermite-error}
\end{align}
No residue theorem or continuity of a moving contour is needed.

\begin{proof}
Fix $t$ on the circle. All $t-z_j$ and $t-x$ in question are nonzero.
The expression
\[
 K_t(x)=\frac{\ell(t)-\ell(x)}{\ell(t)(t-x)}
\]
is, as a function of $x$, a polynomial of degree at most $N-1$:
use the polynomial divided difference of $\ell$. At $x=z_j$ its
value is $1/(t-z_j)$. Polynomial interpolation uniqueness therefore gives
the exact identity
\begin{equation}
 \frac{\ell(t)-\ell(x)}{\ell(t)(t-x)}
   =\sum_j\frac{L_j(x)}{t-z_j}.
 \label{polyquad:finite-resolvent}
\end{equation}
Integrate this finite sum and apply the circle Cauchy formula to $f$
at each $z_j$. This proves \eqref{polyquad:hermite-p}. Subtract it
from the same Cauchy formula at $x$, using
$1/(t-x)-K_t(x)=\ell(x)/(\ell(t)(t-x))$.
This proves \eqref{polyquad:hermite-error}, including $x=z_j$:
the integrand is well-defined on the contour, and $\ell(x)=0$.
\end{proof}

The genuinely analytic input is only the ordinary circle Cauchy formula.
Everything before integration is finite interpolation. If $N=0$, take
$\ell=1$ and $p=0$; the first formula has zero integrand and the second
is exactly Cauchy's formula. The same proof works for a more general
contour whenever its Cauchy formula is supplied explicitly.

\paragraph{Inspected reuse.}
Use \code{Lagrange.eq_interpolate}, \code{Lagrange.nodal}, and
\code{Lagrange.eval_interpolate_at_node}. Mathlib already has
\code{DiffContOnCl.two_pi_i_inv_smul_circleIntegral_sub_inv_smul}
and the more general
\code{Complex.two_pi_I_inv_smul_circleIntegral_sub_inv_smul_of_differentiable_on_off_countable}.
The finite kernel identity is the useful new algebraic lemma; no
real-variable interpolation-remainder theorem is logically needed.

\subsection{Chebyshev endpoint-grid weights: differentiate the monic polynomial}
\label{polyquad:barycentric}
\noindent\textbf{Target: \code{chebyshev-barycentric-weights}.}
Let $n\ge1$, $x_j=\cos(j\pi/n)$ for $0\le j\le n$, and
$\delta_0=\delta_n=1/2$, $\delta_j=1$ otherwise. Then
\begin{align}
 \ell(x):=\prod_{j=0}^{n}(x-x_j)
 &=2^{-n}\bigl(T_{n+1}(x)-T_{n-1}(x)\bigr)
   =2^{1-n}(x^2-1)U_{n-1}(x),\label{polyquad:nodal}\\
 \frac1{\ell'(x_j)}&=\frac{2^{n-1}}{n}(-1)^j\delta_j.
 \label{polyquad:weights}
\end{align}
The catalogue's factor $2^{-n}$ in \eqref{polyquad:nodal} is correct:
$T_{n+1}$ has leading coefficient $2^n$. In particular, the $n=1$
case is $(T_2-T_0)/2=x^2-1$.

\begin{proof}
The Chebyshev recurrence gives
$T_{n+1}-T_{n-1}=2(x^2-1)U_{n-1}$. Its right-hand side vanishes at
the endpoints and, by
$U_{n-1}(\cos\theta)\sin\theta=\sin(n\theta)$, at all interior
grid points. Strict monotonicity of cosine on $[0,\pi]$ makes these
$n+1$ roots distinct. The displayed polynomial is monic of degree $n+1$,
so it is $\ell$.
For $0<j<n$, substitute $x=\cos\theta$ to obtain
\[
 \ell(\cos\theta)=-2^{1-n}\sin\theta\sin(n\theta).
\]
Differentiate at $\theta=j\pi/n$ and cancel the nonzero sine:
$\ell'(x_j)=2^{1-n}n(-1)^j$. At the endpoints differentiate the
factor $(x^2-1)U_{n-1}(x)$ directly. Since
$U_{n-1}(1)=n$ and $U_{n-1}(-1)=(-1)^{n-1}n$, this gives
$\ell'(1)=2^{2-n}n$ and
$\ell'(-1)=2^{2-n}n(-1)^n$. Taking reciprocals proves
\eqref{polyquad:weights}.
\end{proof}

For $x$ away from the nodes, the barycentric formula follows by canceling
the common factor $2^{n-1}/n$:
\[
 p(x)=\frac{\displaystyle\sum_{j=0}^n
          \frac{(-1)^j\delta_j f(x_j)}{x-x_j}}
             {\displaystyle\sum_{j=0}^n
          \frac{(-1)^j\delta_j}{x-x_j}}.
\]
Its denominator is nonzero: before canceling the common factor,
the denominator equals $1/\ell(x)$, by interpolation of the constant
polynomial. At a node use $p(x_j)=f(x_j)$, rather than evaluating a
quotient with a zero denominator. The contract excludes $n=0$ because
the stated cosine grid divides by $n$.

\paragraph{Inspected reuse.}
\code{Lagrange.nodalWeight_eq_eval_derivative_nodal},
\code{Lagrange.eval_interpolate_not_at_node'},
\code{Lagrange.sum_nodalWeight_mul_inv_sub_ne_zero},
\code{Polynomial.Chebyshev.U_real_cos}, and the existing Chebyshev
recurrences contain the needed ingredients. The proved NLAlib theorem
\code{NLAlib.nodal_chebyshevZero} concerns the \emph{roots} grid, so it
must not be mistaken for this endpoint-grid identity.

\subsection{General Gauss quadrature as a finite self-adjoint operator}
\label{polyquad:gauss}
\noindent\textbf{Target: \code{gauss-quadrature-general-measure}.}
Write $N=n+1\ge1$. Let $\mu$ be a finite positive Borel measure supported
on $[a,b]$, whose support contains at least $N$ distinct points. Define
$H=\mathcal P_{<N}$ and
\[
 \langle p,q\rangle_\mu=\int p(x)q(x)\,d\mu(x).
\]
This is an inner product: a nonzero polynomial of degree $<N$ cannot
vanish on the whole support. More explicitly, a support point where
it is nonzero has a neighbourhood of positive measure where its
absolute value is bounded below. All polynomial moments are finite,
since the support is compact.

Define $J:H\to H$ uniquely by
\begin{equation}
 \langle Jp,q\rangle_\mu=\int xp(x)q(x)\,d\mu(x)
 \qquad(p,q\in H).
 \label{polyquad:compression}
\end{equation}
This definition needs only an invertible finite Gram matrix. In the
monomial basis, that matrix is $(\int x^{i+j}\,d\mu)_{i,j<N}$; solve
one such linear system for each column of $J$. Equation
\eqref{polyquad:compression} makes $J$ self-adjoint and gives
\begin{equation}
 J^k1=x^k\quad(0\le k<N),\qquad
 a\|p\|_\mu^2\le\langle Jp,p\rangle_\mu\le b\|p\|_\mu^2.
 \label{polyquad:cyclic}
\end{equation}
The first identity follows by induction, since $Jp=xp$ whenever
$\deg p\le N-2$.

\begin{theorem}
The operator $J$ has $N$ distinct real eigenvalues
$\theta_1,\ldots,\theta_N\in[a,b]$. For unit orthogonal eigenvectors
$v_j$, put $w_j=\langle1,v_j\rangle_\mu^2$. Then
\begin{equation}
 w_j>0,\qquad \sum_jw_j=\mu([a,b]),\qquad
 \int p\,d\mu=\sum_{j=1}^Nw_jp(\theta_j)\quad(\deg p\le2N-1).
 \label{polyquad:gauss-exact}
\end{equation}
If the support has at least $N+1$ points, all $\theta_j\in(a,b)$.
\end{theorem}
\begin{proof}
Apply the finite-dimensional real spectral theorem to $J$. Its spectral
bounds follow from \eqref{polyquad:cyclic}. Any eigenvector $v$
orthogonal to $1$ is orthogonal to every $J^k1$, since
$\langle v,J^k1\rangle=\theta^k\langle v,1\rangle$. Those vectors
form a basis by \eqref{polyquad:cyclic}, a contradiction. Every
eigenspace is therefore one-dimensional: an eigenspace of dimension
at least two would contain a nonzero vector orthogonal to $1$.
The same observation proves each $w_j>0$. Parseval gives their sum.

For $0\le k\le2N-2$, choose $r,s<N$ with $r+s=k$ and compute
\[
 \int x^k\,d\mu=\langle x^r,x^s\rangle_\mu
 =\langle J^r1,J^s1\rangle_\mu
 =\langle1,J^k1\rangle_\mu=\sum_jw_j\theta_j^k.
\]
For the remaining moment $k=2N-1$, use
\[
 \int x^{2N-1}\,d\mu
 =\langle Jx^{N-1},x^{N-1}\rangle_\mu
 =\langle1,J^{2N-1}1\rangle_\mu.
\]
Linearity proves exactness. If an eigenvalue equals $a$ with eigenvector
$p$, then $\int(x-a)p(x)^2\,d\mu=0$. The nonnegative continuous
integrand vanishes on the support, putting the support in
$\{a\}\cup\{p=0\}$, a set with at most $N$ points. This contradicts
the stronger support hypothesis. The argument at $b$ is identical.
\end{proof}

To identify these nodes with orthogonal-polynomial roots, let
$\pi_N=x^N-h$ where $h\in H$ is the unique Gram projection of $x^N$.
Then $\pi_N$ is monic and orthogonal to $H$, and
$\pi_N(J)1=0$. Since $1$ is cyclic and polynomial operators commute,
$\pi_N(J)=0$. Both $\pi_N$ and $\det(xI-J)$ are monic of degree $N$;
their difference annihilates $1$ by Cayley--Hamilton and has degree
less than $N$. The linear independence in \eqref{polyquad:cyclic}
forces that difference to be zero. Thus
\[
 \pi_N(x)=\det(xI-J)=\prod_j(x-\theta_j).
\]
If $L_j$ are the cardinal polynomials of these nodes, exactness gives
$w_j=\int L_j\,d\mu=\int L_j^2\,d\mu>0$.

There is also a shorter proof once the distinct roots of $\pi_N$ are
available. Divide $p=\pi_Nq+r$ with $\deg q<N$, $\deg r<N$ for
$\deg p\le2N-1$. Orthogonality kills the first integral, and
interpolation integrates $r$ exactly. Apply this to $L_j^2$ for
positivity. For arbitrary distinct nodes, interpolation alone gives
exactness through degree $N-1$; positivity is not automatic.

\paragraph{Boundary contracts.}
With exactly $N$ support points, $\pi_N$ has zero squared norm and is
the support polynomial; the formula is still exact and its weights are
the positive atomic masses. Calling $\pi_N$ a \emph{normalized}
orthogonal polynomial would incorrectly exclude this case. Fewer than
$N$ support points require reducing to their actual number; a positive
$N$-node rule at distinct nodes is not the appropriate contract. The
zero measure uses the empty rule. The case $a=b$ is possible only in
the one-support-point situation.

\paragraph{Inspected reuse.}
The discrete-matrix specialization already has
\code{NLAlib.LanczosDecomp.dotProduct_aeval_mulVec_eq_sum},
\code{NLAlib.LanczosDecomp.integral_eval_charpoly_mul_eq_zero}, and
\code{NLAlib.integral_eval_spectralMeasure}.
The new interface is the finite moment Gram space, rather than an
infinite-dimensional multiplication operator. For the shorter proof,
use \code{Polynomial.modByMonic_add_div} and
\code{Lagrange.eq_interpolate}. The finite-dimensional spectral theorem
and \code{Matrix.aeval_self_charpoly} supply the operator proof.

\subsection{Chebyshev integrals from a polynomial antiderivative}
\label{polyquad:integrals}
\noindent\textbf{Target: \code{chebyshev-integral-formula}.}
Put $m_k=\int_{-1}^1T_k(x)\,dx$. Then
\begin{equation}
 m_k=\begin{cases}2/(1-k^2),&k\text{ even},\\0,&k\text{ odd}.
 \end{cases}
 \label{polyquad:moments}
\end{equation}
In particular $m_0=2$ and $m_1=0$; no expression $0/0$ is used.
For $k\ge2$, the polynomial
\[
 F_k=\frac{T_{k+1}}{2(k+1)}-\frac{T_{k-1}}{2(k-1)}
\]
has derivative $T_k$, since $T_j'=jU_{j-1}$ and
$U_k-U_{k-2}=2T_k$. Evaluating $F_k(1)-F_k(-1)$, using
$T_j(1)=1$, $T_j(-1)=(-1)^j$, proves \eqref{polyquad:moments}.
The cases $k=0,1$ are immediate. Thus, for every finite coefficient list,
\[
 \int_{-1}^1\sum_{k=0}^n c_kT_k(x)\,dx
    =\sum_{\substack{0\le k\le n\\ k\text{ even}}}
        \frac{2c_k}{1-k^2}.
\]
The same identity holds for an infinite series with
$\sum_k|c_k|<\infty$, by uniform convergence from $|T_k|\le1$.
This proof uses polynomial differentiation and the fundamental theorem
of calculus, avoiding trigonometric substitution.

\paragraph{Inspected reuse.}
The exact identities already exist as
\code{Polynomial.Chebyshev.T_derivative_eq_U},
\code{Polynomial.Chebyshev.two_mul_T_eq_U_sub_U},
\code{Polynomial.Chebyshev.T_eval_one}, and
\code{Polynomial.Chebyshev.T_eval_neg_one}.

\subsection{Quadrature errors: a scalar functional on Chebyshev coordinates}
\label{polyquad:errors}
\noindent\textbf{Target: \code{gauss-clenshaw-curtis-quadrature-error}.}
Here $I(f)=\int_{-1}^1 f(x)\,dx$ is specifically the \emph{Lebesgue}
integral. Let $C_n$ denote the $(n+1)$-point Clenshaw--Curtis rule,
and $G_n$ the $(n+1)$-point Gauss--Legendre rule. For a rule $Q$, write
$E_Q=I-Q$.

\paragraph{Audit of the requested bounds.}
The catalogue copied a printing error from the 2013 edition of ATAP.
The author's own
\href{https://people.maths.ox.ac.uk/trefethen/ATAP/}{errata page}
replaces the denominators in (19.10) and (19.12) by
$(2n+1-\nu)^\nu$. They are \emph{not}
$(n-2\nu-1)^{2\nu+1}$. We prove the corrected bounds, preserve their
constants, and give an explicit sufficient threshold for the eventual
Clenshaw--Curtis conclusion. All the estimates below follow from finite
aliasing and scalar series; a separate asymptotic quadrature theorem is
unnecessary.

\paragraph{Step 1: the rules are positive, symmetric, and exact.}
Gauss positivity and exactness through degree $2n+1$ follow from
Section~\ref{polyquad:gauss}. Reflection $p(x)\mapsto p(-x)$ is an
isometry of its polynomial space and conjugates $J$ to $-J$. Consequently
the nodes and their spectral weights are symmetric, so $G_n(T_k)=0$
for odd $k$.

For Clenshaw--Curtis put $\theta_j=j\pi/n$ and use the $\delta_j$
from Section~\ref{polyquad:barycentric}. Finite discrete cosine
orthogonality gives the cardinal polynomial
\[
 L_j(x)=\frac{2\delta_j}{n}
  \left(\frac12+\sum_{k=1}^{n-1}\cos(k\theta_j)T_k(x)
             +\frac12\cos(n\theta_j)T_n(x)\right).
\]
One can verify this identity just at the $n+1$ nodes: the needed cosine
sums follow from a finite geometric sum of $2n$-th roots of unity.
Integration with \eqref{polyquad:moments} gives
\begin{equation}
 w_j=\frac{2\delta_j}{n}\left[
  1-\sum_{k=1}^{\lfloor(n-1)/2\rfloor}
       \frac{2\cos(2k\theta_j)}{4k^2-1}
  -\boldsymbol{1}_{2\mid n}\frac{\cos(n\theta_j)}{n^2-1}
  \right] \qquad(n\ge2).
 \label{polyquad:cc-weights}
\end{equation}
For $n=2m+1$, the bracket is at least
$1-\sum_{k=1}^m2/(4k^2-1)=1/n>0$.
For $n=2m$, it is at least
$1-\sum_{k=1}^{m-1}2/(4k^2-1)-1/(n^2-1)=n/(n^2-1)>0$.
These are telescoping sums, since
$2/(4k^2-1)=1/(2k-1)-1/(2k+1)$.
For $n=1$, the two weights are $1,1$.
Thus the rule is positive. Interpolation gives exactness through degree
$n$, and reflection of the grid gives symmetry. In both cases the
weights sum to $2$, so
\begin{equation}
 |E_Q(f)|\le4\|f-p\|_\infty
 \quad\text{whenever $Q$ is exact on the polynomial $p$.}
 \label{polyquad:positive-error}
\end{equation}
Polynomial density now proves $C_n(f),G_n(f)\to I(f)$ for every
continuous $f$: fix a uniformly close polynomial, then take $n$ large
enough to integrate it exactly.

\paragraph{Step 2: one scalar bound produces both analytic constants.}
Define the folded residue $r_n(k)\in\{0,\ldots,n\}$ by reducing $k$
modulo $2n$ and reflecting the upper half. At the nodes,
$T_k(x_j)=T_{r_n(k)}(x_j)$, so
\begin{equation}
 C_n(T_k)=m_{r_n(k)}.
 \label{polyquad:alias-moment}
\end{equation}
Parity is preserved by $r_n$, hence both errors vanish on all odd modes.
For even $k>n$ with $n\ge2$, necessarily $k\ge4$, and
\[
 |E_{C_n}(T_k)|\le |m_k|+|m_{r_n(k)}|
                 \le\frac2{15}+2=\frac{32}{15}.
\]
For even $k>2n+1$ with $n\ge1$, Gauss positivity gives
$|G_n(T_k)|\le2$ and again $|m_k|\le2/15$. Thus
\begin{equation}
 |E_{C_n}(f)|\le\frac{32}{15}
        \sum_{\substack{k>n\\k\text{ even}}}|a_k|,
 \qquad
 |E_{G_n}(f)|\le\frac{32}{15}
        \sum_{\substack{k>2n+1\\k\text{ even}}}|a_k|,
 \label{polyquad:master-tail}
\end{equation}
whenever $f=\sum_{k\ge0}a_kT_k$ absolutely and uniformly.
For an analytic $f$ bounded by $M$ in the open Bernstein ellipse
$E_\rho$, $\rho>1$, the known coefficient estimate is
$|a_k|\le2M\rho^{-k}$ for $k\ge1$. Summing the even geometric tails
in \eqref{polyquad:master-tail} yields precisely
\begin{align}
 |E_{C_n}(f)|&\le\frac{64}{15}
       \frac{M\rho^{1-n}}{\rho^2-1}\quad(n\ge2),
       \label{polyquad:cc-analytic}\\
 |E_{G_n}(f)|&\le\frac{64}{15}
       \frac{M\rho^{-2n}}{\rho^2-1}\quad(n\ge1).
       \label{polyquad:gauss-analytic}
\end{align}
When $n$ is even, the first bound has the slightly stronger numerator
$M\rho^{-n}$. No boundary continuity on the ellipse is required when
using the already-proved open-ellipse coefficient estimate.

\paragraph{Step 3: the corrected variation bounds with their constants.}
Assume $\nu\ge1$, $f,\ldots,f^{(\nu-1)}$ are absolutely continuous,
and $f^{(\nu)}$ has total variation $V$. The coefficient-decay result
under precisely these hypotheses gives, for $k>\nu$,
\begin{equation}
 |a_k|\le\frac{2V}{\pi k(k-1)\cdots(k-\nu)}
       \le\frac{2V}{\pi(k-\nu)^{\nu+1}}.
 \label{polyquad:bv-input}
\end{equation}
This is the scalar bounded-variation input, rather than a new
assumption about quadrature. It also gives absolute uniform convergence
of the Chebyshev series after identifying its coefficients.
For the first even integer $m>s$, convexity of
$g(t)=(t-\nu)^{-\nu-1}$ on $t>\nu$ gives
\[
 \sum_{j=0}^\infty g(m+2j)
 \le\frac12\int_{m-1}^{\infty}g(t)\,dt
 =\frac1{2\nu(m-1-\nu)^\nu}.
\]
Indeed each midpoint value is at most the average on the interval of
length two centred there. Substitute this into
\eqref{polyquad:master-tail} to obtain
\begin{align}
 |E_{C_n}(f)|&\le\frac{32}{15}
       \frac{V}{\pi\nu(n-\nu)^\nu}\quad(n>\nu, n\ge2),
       \label{polyquad:cc-bv}\\
 |E_{G_n}(f)|&\le\frac{32}{15}
       \frac{V}{\pi\nu(2n+1-\nu)^\nu}\quad(2n+1>\nu, n\ge1).
       \label{polyquad:gauss-bv}
\end{align}
The Gauss condition is even weaker than retaining the old printed
restriction $n>2\nu+1$.

\paragraph{Step 4: the eventual Clenshaw--Curtis bound by finite blocks.}
We now prove the analogue of the corrected ATAP Theorem 19.5 with an
explicit threshold. Set $C_\nu=2^{\nu+2}V/\pi$. For
$n\ge\max(2,2\nu)$ and $k>n$, \eqref{polyquad:bv-input} implies
$|a_k|\le C_\nu k^{-\nu-1}$.
In any complete block of even integers of length $2n$, the folded
residues count $0$ once and every positive even residue at most twice.
Consequently, using a telescoping sum,
\begin{equation}
 \sum_{\text{one block}}|m_{r_n(k)}|
 \le2+2\sum_{j=1}^{\infty}\frac2{4j^2-1}=4.
 \label{polyquad:block-four}
\end{equation}
The initial half-block $n<k\le2n$ satisfies the same bound. For the
subsequent blocks $2\ell n<k\le2(\ell+1)n$, $\ell\ge1$, use
$k>2\ell n$. Since $\sum_{\ell\ge1}\ell^{-\nu-1}\le2$, this gives
\[
 \sum_{\substack{k>n\\k\text{ even}}}
       k^{-\nu-1}|m_{r_n(k)}|
 \le\frac4{n^{\nu+1}}+
       \frac4{(2n)^{\nu+1}}\sum_{\ell\ge1}\ell^{-\nu-1}
 \le\frac6{n^{\nu+1}}.
\]
Also $|m_k|\le(8/3)k^{-2}$ for $k\ge2$, so
\[
 \sum_{k>n}k^{-\nu-1}|m_k|
 \le\frac8{3(\nu+2)n^{\nu+2}}
 \le\frac1{n^{\nu+1}}.
\]
Together with \eqref{polyquad:alias-moment}, these inequalities prove
the fully explicit stronger-rate estimate
\begin{equation}
 |E_{C_n}(f)|\le\frac{7\,2^{\nu+2}V}{\pi n^{\nu+1}}
 \quad(n\ge\max(2,2\nu)).
 \label{polyquad:cc-extra-power}
\end{equation}
Since $2n+1-\nu\le2n$, comparison of constants gives, in particular,
\begin{equation}
 |E_{C_n}(f)|\le\frac{32}{15}
       \frac{V}{\pi\nu(2n+1-\nu)^\nu}
       \qquad\bigl(n\ge14\nu4^\nu\bigr).
 \label{polyquad:cc-eventual}
\end{equation}
Indeed $14>105/8$, the constant required after multiplying the two
positive denominators. This threshold is deliberately conservative;
its purpose is an elementary explicit proof of the original eventual
claim with its corrected constant. The case $V=0$ is included: the
coefficient tail vanishes and the rules are exact once their polynomial
degrees suffice.

\paragraph{Why the Lebesgue constant does not transfer to every measure.}
For an arbitrary positive measure of mass $m$, an $N$-node Gauss rule
does satisfy
\[
 |I_\mu(f)-Q_N(f)|\le2m\inf_{\deg p\le2N-1}\|f-p\|_\infty
 \le\frac{4mM\rho^{1-2N}}{\rho-1}.
\]
The last estimate is the direct analytic coefficient tail. The sharper
constant \eqref{polyquad:gauss-analytic} uses the specific Lebesgue
moments \eqref{polyquad:moments}, so mass scaling alone does not extend
it. For example, let
$\mu=(\delta_{-1}+\delta_1)/2+\delta_0$, of mass $2$.
Its two-node Gauss rule has nodes $\pm1/\sqrt2$ and weights $1,1$.
For $f=T_4$, the integral is $2$ and the rule gives $-2$, so the error
is $4$. On $E_\rho$, take $M=(\rho^4+\rho^{-4})/2$. The proposed
Lebesgue bound with $n=1$ tends to $32/15$ as $\rho\to\infty$, a
contradiction. If infinite support is required, replace $\mu$ by
$(1-\varepsilon)\mu+\varepsilon\,dx$ for sufficiently small positive
$\varepsilon$; the degree-two moment construction and its error vary
continuously and the strict violation persists.

\paragraph{Inspected reuse and proof boundary.}
The pinned declarations include \code{NLAlib.chebyshevAliasIndex},
\code{NLAlib.chebyshevAliasIndex_le},
\code{NLAlib.abs_eval_interpolate_eval_T_node_le_one},
\code{NLAlib.norm_chebyshevComplexCoeff_le_of_open_ellipse}, and
\code{NLAlib.hasSum_chebyshevAnalyticCoeff_mul_eval_T_of_open_ellipse}.
The scalar bounded-variation coefficient bound
\eqref{polyquad:bv-input} is the separate coefficient-decay target;
once supplied, the proof above adds only finite cosine sums, positive
finite sums, and scalar tail estimates. The reference-only asymptotic
proof of ATAP Theorem 19.5 has been replaced by
\eqref{polyquad:block-four}--\eqref{polyquad:cc-eventual}.

\subsection{Colleague matrices are multiplication modulo a polynomial}
\label{polyquad:colleague}
\noindent\textbf{Target: \code{colleague-matrix}.}
Let $p=\sum_{k=0}^n a_kT_k$ with $n\ge1$, $a_n\ne0$, and put
$q=p/(2^{n-1}a_n)$, a monic degree-$n$ polynomial. On
$H=\mathcal P_{<n}$ define the finite linear operator
\[
 M h=\text{the remainder of }xh\text{ upon division by }q.
\]
Polynomial division immediately gives
\begin{equation}
 h(M)1=h\bmod q\quad\text{for every polynomial }h.
 \label{polyquad:remainder-cyclic}
\end{equation}
For example, prove it first for monomials by induction, using uniqueness
of remainders, and then use linearity. Thus
$1,M1,\ldots,M^{n-1}1=1,x,\ldots,x^{n-1}$ is a basis.
Cayley--Hamilton and \eqref{polyquad:remainder-cyclic} say that the
characteristic polynomial of $M$ is divisible by $q$. Both are monic
of degree $n$, hence
\begin{equation}
 \det(xI-M)=q(x)=\frac{p(x)}{2^{n-1}a_n}.
 \label{polyquad:colleague-charpoly}
\end{equation}
This establishes multiplicities directly, even when $p$ has repeated
roots. There is no need to infer a characteristic polynomial from a
list of distinct eigenvectors.

In the basis $T_0,\ldots,T_{n-1}$, the coordinate matrix of $M$ is
the transpose of the following evaluation-recurrence matrix $C$.
For $n\ge2$, its first row has $C_{0,1}=1$ and all other entries zero;
for $1\le j\le n-2$, its $j$th row has $1/2$ in columns $j-1,j+1$;
and its last row is
\begin{equation}
 C_{n-1,k}=\frac12\boldsymbol{1}_{k=n-2}
                  -\frac{a_k}{2a_n}\qquad(0\le k<n).
 \label{polyquad:colleague-entries}
\end{equation}
Indeed $xT_0=T_1$, $xT_j=(T_{j-1}+T_{j+1})/2$ for $j\ge1$,
and modulo $p$ one replaces
$T_n$ by $-\sum_{k<n}(a_k/a_n)T_k$.
Transposition does not change the characteristic polynomial, so
\eqref{polyquad:colleague-charpoly} is the requested identity for $C$.
The leading coefficients of the Chebyshev polynomials make this a
basis by a triangular change from monomials.

\paragraph{The degree-one boundary must be stated separately.}
For $n=1$, $p=a_0+a_1x$ and the correct matrix is
$C=[-a_0/a_1]$. Applying the half-scaled last-row formula literally
would give $-a_0/(2a_1)$ and would be wrong. No root-finding matrix is
needed for a nonzero constant polynomial: its degree-zero counterpart
is the empty matrix with characteristic polynomial $1$. The zero
polynomial is excluded by $a_n\ne0$.

\paragraph{Inspected reuse.}
The algebra is \code{Polynomial.modByMonic_add_div},
\code{Polynomial.Chebyshev.T_add_two},
\code{Matrix.aeval_self_charpoly}, and invariance of characteristic
polynomials under change of basis and transposition. Mathlib also has
\code{charpoly_leftMulMatrix} for a power basis. Thus a companion-matrix
API is optional: the remainder operator gives the same theorem
directly and is reusable for Jacobi and other polynomial bases.

% END quadrature.tex

\section{Sharp asymptotic targets: exact boundaries and finite substitutes}
\label{polyroot:deep}
\paragraph{Targets.}\code{abs-best-approx-asymptotics},
\code{exp-rational-halfline}.
These entries ask for substantially more than existence of useful
approximants: exact limiting constants in best approximation. The
arguments below give elementary, explicit operator-ready estimates and
identify the additional theorems needed for the sharp catalogue limits.
The second-pass sections below now prove the polynomial limit and
the sharp reciprocal-polynomial rate, and improve the rational
absolute-value construction to the sharp root exponent.

\subsection{Polynomial approximation of absolute value: a finite dual certificate}
The Chebyshev series in Section~\ref{polyseries:abs-sign} gives, for $n\ge1$,
\begin{equation}\label{polyroot:abs-upper}
 E_n(|x|;[-1,1])\le \frac{2}{\pi n}.
\end{equation}
Here one uses its truncation at the largest even degree $2m\le n$,
whose error is $2/(\pi(2m+1))$.
There is also a finite trigonometric lower certificate of the correct
order. Put $K=\lfloor n/2\rfloor+1$ and define
\[
 F_K(t)=\frac1{K+1}\left|\sum_{j=0}^K e^{ijt}\right|^2
 =\sum_{j=-K}^K\left(1-\frac{|j|}{K+1}\right)e^{ijt}.
\]
It is nonnegative and has normalized integral $1$ over a period.
For a continuous $h$ on $[-1,1]$, set
\[
 L(h)=\frac1{2\pi}\int_{-\pi}^{\pi}
 h(\sin t)\cos(4Kt)F_K(2t)\,dt.
\]
Then $|L(h)|\le\|h\|_\infty$. The trigonometric polynomial
$p(\sin t)$ has frequencies at most $n$ if $\deg p\le n$.
Every frequency of $\cos(4Kt)F_K(2t)$ has absolute value at least
$2K>n$, so finite exponential orthogonality gives $L(p)=0$.
For $g(t)=|\sin t|$, direct integration gives the Fourier coefficients
\[
 \widehat g(2j)=-\frac{2}{\pi(4j^2-1)}\quad(j\ne0),
 \qquad \widehat g(0)=2/\pi,
\]
and zero odd coefficients. All frequencies used by $L$ are nonzero.
Writing $w_j=1-|j|/(K+1)$ and using $w_j=w_{-j}$ gives
\[
 |L(|x|)|
 =\sum_{j=-K}^K w_j\frac{2}{\pi(4(2K+j)^2-1)}
 \ge\frac{2(K+1)}{\pi(36K^2-1)}
 \ge\frac1{18\pi(n+1)}.
\]
The identity $\sum_jw_j=K+1$ is a finite arithmetic sum. Since
$L(p)=0$, every degree-$n$ polynomial has error at least this value:
\begin{equation}\label{polyroot:abs-lower}
 E_n(|x|;[-1,1])\ge\frac1{18\pi(n+1)}.
\end{equation}
Only finitely many Fourier coefficients and a finite Gram square are
used. These bounds alone do not establish existence of
$\lim_{n\to\infty}nE_n(|x|)=\beta$. Section~\ref{polyretry:bernstein}
now supplies that missing argument and an exact variational
characterization. The numerical investigation by
\href{https://doi.org/10.1007/BF01890040}{Varga--Carpenter (1985)}
is separate; no decimal approximation to $\beta$ is being treated as
an exact constant.

\subsection{A rational root-exponential bound from a finite product}
For an integer $m\ge2$, let
\[
 P_m(x)=\prod_{k=1}^{m^2}(x+e^{-k/m}),\qquad
 R_m(x)=x\frac{P_m(x)-P_m(-x)}{P_m(x)+P_m(-x)}.
\]
This is a real even rational function of type at most
$(m^2+1,m^2)$. Its denominator has no zero on $[-1,1]$:
for $x>0$, every factor satisfies
$|e^{-k/m}-x|<e^{-k/m}+x$, hence $|P_m(-x)|<P_m(x)$;
at zero the denominator is $2P_m(0)>0$, and it is even.

For $0\le x\le e^{-m}$ all factors of $P_m(-x)$ are nonnegative.
Thus $0\le R_m(x)\le x$ and $|x-R_m(x)|\le e^{-m}\le2^{-m}$.
For $e^{-m}\le x\le1$, put $t=-\log x\in[0,m]$.
Choose $j\in\{0,\ldots,m-1\}$ with $t\in[j,j+1]$.
For the $m$ indices $jm+1\le k\le(j+1)m$, $|t-k/m|\le1$, so
\[
 \frac{|e^{-k/m}-x|}{e^{-k/m}+x}
 \le\frac{e-1}{e+1}<\frac12.
\]
All other factor ratios are at most one. Hence
$|u|\le2^{-m}$ for $u=P_m(-x)/P_m(x)$, and
\[
 |x-R_m(x)|=\left|\frac{2xu}{1+u}\right|
 \le\frac{2\,2^{-m}}{1-2^{-m}}\le3\,2^{-m}.
\]
Evenness proves the same estimate on the negative half-interval:
\begin{equation}\label{polyroot:newman-simple}
 \||x|-R_m(x)\|_{[-1,1]}\le3\,2^{-m}.
\end{equation}
This is a complete elementary root-exponential construction with an
explicit degree and a nonvanishing denominator. It is weaker than the
catalogue's sharp $8e^{-\pi\sqrt n}$ asymptotic.

For the exact rational asymptotic, rational uniqueness implies that the
best approximant to $|x|$ is even. In reduced form an even rational
function with denominator nonzero at zero has both numerator and
denominator even: normalize the denominator at zero, use uniqueness of
coprime representation under $x\mapsto-x$, then compare coefficients.
Thus an even type-$(2m,2m)$ approximant is $s(x^2)$ with $s$ of type
$(m,m)$, and
\[
 E_{2m,2m}(|x|;[-1,1])=E_{m,m}(\sqrt t;[0,1]).
\]
The same error holds at odd type $(2m+1,2m+1)$, by evenness and degree.
\href{https://arxiv.org/abs/math/9301217}{Stahl's theorem}
asserts
$e^{2\pi\sqrt{\alpha m}}E_{m,m}(t^\alpha;[0,1])
\to4^{1+\alpha}|\sin\pi\alpha|$.
Taking $\alpha=1/2$ gives the catalogue limit $8$; the odd indices have
the same limit because $\sqrt{2m+1}-\sqrt{2m}\to0$.
This is a rigorous reduction to one precise deep theorem, not a new
finite-dimensional proof of Stahl's theorem.

\subsection{A geometric rational bound on the whole positive half-line}
For $n\ge1$ define
\[
 P_n(x)=\sum_{k=0}^n\frac{x^k}{k!},\qquad r_n(x)=P_n(x)^{-1}.
\]
The denominator is at least one on $[0,\infty)$, and
$0\le r_n(x)-e^{-x}$ because $P_n(x)\le e^x$.
For $0\le x\le n/4$, Taylor's remainder gives
\[
 r_n(x)-e^{-x}
 =\frac{e^x-P_n(x)}{e^xP_n(x)}
 \le\frac{x^{n+1}}{(n+1)!}
 \le(e/4)^{n+1}\le(4/5)^n.
\]
Here $N!\ge(N/e)^N$ follows by comparing
$\sum_{j=1}^N\log j$ with $\int_1^N\log t\,dt$; also $e<3$ gives
$(e/4)^{n+1}<(3/4)^{n+1}\le(4/5)^n$.
For $x\ge n/4$, the coefficient inequalities
$\binom nk/n^k\le1/k!$ imply
\[
 P_n(x)\ge(1+x/n)^n\ge(5/4)^n,
 \qquad 0\le r_n(x)-e^{-x}\le(4/5)^n.
\]
Consequently
\begin{equation}\label{polyroot:halfline-rational}
 \sup_{x\ge0}|e^{-x}-r_n(x)|\le(4/5)^n.
\end{equation}
For every finite-dimensional PSD $A$, independent of its largest
eigenvalue,
$\|e^{-A}-P_n(A)^{-1}\|\le(4/5)^n$.
Replacing $x$ by $-x$ gives the catalogue's negative half-line.
In contrast any nonconstant polynomial is unbounded there; the best
constant approximation to $e^x$ on $(-\infty,0]$ is $1/2$, with error
$1/2$. Thus polynomial best errors stay $1/2$ in every degree.

The rate $(4/5)^n$ proves geometric convergence but does not prove the
sharp type-$(0,n)$ root limit $1/3$, nor the type-$(n,n)$ equivalent
$2H^{n+1/2}$ with $H=1/9.28902549\ldots$.
The latter sharp constant is an analytic asymptotic theorem; see
\href{https://www.mathnet.ru/php/archive.phtml?wshow=paper&jrnid=sm&paperid=619&option_lang=eng}{Aptekarev (2002)}.
The type-$(0,n)$ limit does not follow from this Taylor construction,
but is now fully proved by the finite Gram argument in
Section~\ref{polyretry:exp}. Only the diagonal Halphen equivalent
remains an external mathematical input for this entry.


% BEGIN retry_lebesgue.tex
\section{Second attempt: the exact Lebesgue constant from a finite cotangent sum}
\label{polyretry:lebesgue}

\textbf{Status: complete derivation of the sharp universal lower bound.}
First, explicit cotangent samples prove the numerical bound for
\emph{every interval peak of the Chebyshev-root grid}.  Then
Sections~\ref{polyretry:finite-comparison}--\ref{polyretry:proper-inverse}
prove the Bernstein--Erd\H{o}s peak comparison needed for arbitrary grids.
The proof uses a positive sum-of-squares Wronskian, a finite matrix
with positive off-diagonal entries, compactness, and an explicit
lifting argument.  No specialized interpolation theorem remains
an imported hypothesis.  This closes the missing sharp part of
\code{lebesgue-constant-lower-bound}; the equispaced bound was
already proved in Section~\ref{polybest:lower}.

Write $N=n+1$ for the number of nodes and put
\[
 B_N=\frac2\pi\left(\log N+\gamma+\log\frac4\pi\right),
 \qquad H_m=\sum_{j=1}^{m}\frac1j,
\]
where $\gamma=\lim_m(H_m-\log m)$.
For the $N$ Chebyshev roots
$\xi_j=\cos((j+\tfrac12)\pi/N)$, $0\le j<N$, let
$\mathcal L_N(x)$ denote their Lebesgue function.  Its order of indexing
does not affect the function.

\subsection{Evaluate at extrema instead of finding the local maxima}
\label{polyretry:cotangent-samples}

For $N\ge2$ and $1\le k<N$ set $y_k=\cos(k\pi/N)$ and
$S_k=\mathcal L_N(y_k)$.  There is exactly one such point strictly
between each pair of consecutive roots.  The cardinal formula obtained
by differentiating $T_N$ gives
\begin{equation}\label{polyretry:sample}
 S_k=\frac1N\sum_{j=0}^{N-1}
 \frac{\sin((j+\tfrac12)\pi/N)}
 {|\cos(k\pi/N)-\cos((j+\tfrac12)\pi/N)|}.
\end{equation}
Indeed $|T_N(y_k)|=1$ and
$|T_N'(\xi_j)|=N/\sin((j+\tfrac12)\pi/N)$.

Here is a finite calculation locating the smallest \emph{sample},
which is enough for a lower bound on all interval maxima.  Put
$h=\pi/(2N)$ and
$C_r=\sum_{j=0}^{r-1}\cot((j+\tfrac12)h)$, $0\le r\le2N$.
The identity
\[
 \frac{\sin\theta}{\cos\phi-\cos\theta}
 =\frac12\left(\cot\frac{\theta-\phi}{2}
                  +\cot\frac{\theta+\phi}{2}\right)
\]
and $C_{N+r}=C_{N-r}$ give, by splitting the sum at $j=k$,
\begin{equation}\label{polyretry:sample-C}
 S_k=\frac{C_{N-k}+C_k-C_{2k}}N.
\end{equation}
Consequently
\[
 N(S_{k+1}-S_k)=2\cot((2k+1)h)
 -\cot((2k+\tfrac12)h)-\cot((2k+\tfrac32)h).
\]
For $1\le k<\lfloor N/2\rfloor$ this is nonpositive, since
$\cot''x=2\cot x\csc^2x>0$ on $(0,\pi/2)$.
Also $S_k=S_{N-k}$ by reflection.  Thus the smallest sample is
$S_m$, where $m=\lfloor N/2\rfloor$.
Using $\cot u-\tan u=2\cot(2u)$ in
\eqref{polyretry:sample-C} gives the explicit formula
\begin{equation}\label{polyretry:central-cot}
 S_m=\frac2N\sum_{j=0}^{m-1}
       \cot\frac{(j+\tfrac12)\pi}{N}
 +\begin{cases}
 0,&N=2m,\\[2pt]
 N^{-1}\tan\dfrac{\pi}{4N},&N=2m+1.
 \end{cases}
\end{equation}
In particular no theorem about the locations or monotonicity of the
\emph{true} local maxima is needed.

\subsection{The exact logarithmic constant by midpoint concavity}

Define $g(x)=\cot x-1/x$ on $(0,\pi/2]$, with $g(0)=0$.
This function is concave.  To check the only nontrivial sign,
\[
 g''(x)=\frac{2\cos x}{\sin^3x}-\frac2{x^3}\le0
 \quad\Longleftrightarrow\quad
 \sin^3x\ge x^3\cos x.
\]
For $0<x<\pi/2$, differentiate
$3\log(\sin x/x)-\log\cos x$.  The sign of its derivative is
the sign of
\[
 K(x)=x(2+\cos2x)-\tfrac32\sin2x.
\]
But $K(0)=0$ and
$K'(x)=4\sin x(\sin x-x\cos x)\ge0$, because
$(\sin x-x\cos x)'=x\sin x\ge0$.  Taking the limit at zero
proves the claim; continuity handles $\pi/2$.

The following explicit harmonic remainder will also be useful:
\begin{equation}\label{polyretry:harmonic-remainder}
 D_m:=2H_{2m}-H_m-\log m-\gamma-2\log2
 \ge\frac1{6(2m+1)^2}>0\qquad(m\ge1).
\end{equation}
Indeed $D_m\to0$ and
\[
 D_k-D_{k+1}
 =\log\frac{k+1}{k}-\frac2{2k+1}
 \ge\frac{2}{3(2k+1)^3}.
\]
The inequality follows without a logarithm power-series theorem:
for $0\le t<1$,
$\log((1+t)/(1-t))-2t$ has derivative
$2t^2/(1-t^2)\ge2t^2$, so its value is at least $2t^3/3$.
Telescope from $k=m$ to infinity, and bound the decreasing positive
sum below by its integral:
$\int_m^\infty2/[3(2x+1)^3]\,dx=1/[6(2m+1)^2]$.

Let $L=m\pi/N$.  Midpoint concavity on the $m$ equal subintervals
of $[0,L]$, together with
$\int_0^L g(x)\,dx=\log(\sin L/L)$, gives
\begin{align}
 \frac2N\sum_{j=0}^{m-1}\cot\frac{(j+\tfrac12)\pi}{N}
 &\ge\frac2\pi\left(2H_{2m}-H_m+
                  \log\frac{\sin L}{L}\right)\notag\\
 &=B_N+\frac2\pi\bigl(D_m+\log\sin L\bigr).
 \label{polyretry:midpoint-bound}
\end{align}
For even $N$, $L=\pi/2$, so this proves $S_m>B_N$ immediately.

For odd $N\ge3$, put $a=\pi/(2N)$; then $L=\pi/2-a$.
The extra tangent in \eqref{polyretry:central-cot} shows that it is
enough to prove
\[
 D_m\ge E(a):=-\log\cos a-a\tan(a/2).
\]
This small correction has a direct bound.  On $0\le a\le\pi/6$,
\[
 E'(a)=\tfrac12\sec^2(a/2)(\tan a-a),\qquad
 \tan a-a=\int_0^a\tan^2t\,dt\le\frac{a^3}{3\cos^2a}.
\]
Using $\cos^2a,\cos^2(a/2)\ge3/4$ and integrating gives
$E(a)\le2a^4/27<a^4/8$.  Since $\pi^4<192$ and $N\ge3$,
\[
 E(a)<\frac{\pi^4}{128N^4}<\frac1{6N^2}\le D_m.
\]
This completes the odd case as well.  We have proved, for every $N\ge2$,
\begin{equation}\label{polyretry:all-intervals}
 \min_{1\le k<N}S_k>B_N,
 \qquad
 \min_{\text{consecutive root intervals }J}
            \max_{x\in J}\mathcal L_N(x)>B_N.
\end{equation}
For $N=2$ the single sample is $1$; the formula and proof include this
case.  For $N=1$ the Lebesgue constant is $1>B_1$: for example,
$\gamma<1$, $\log(4/\pi)<1/3$, and $\pi>8/3$ suffice.

\subsection{Reduce arbitrary nodes to interval peak comparison}

Call a node set canonical if its first and last nodes are $-1,1$.
For such a set $X$ with $N\ge2$ nodes, write
$M_i(X)=\max_{[x_{i-1},x_i]}\sum_j|\ell_j^X(x)|$.
The comparison theorem proved in the following subsections is
\begin{equation}\label{polyretry:comparison}
 \min_i M_i(Y)\le\max_i M_i(X)
 \quad\text{for any two canonical sets $X,Y$ of the same size.}
\end{equation}
This is a theorem of Bernstein--Erd\H{o}s type, established by
Kilgore and de Boor--Pinkus.  It does not follow merely from the standard
alternation theorem for best polynomial approximation: here the nodes,
and hence the interpolation operator, vary.

Here is the exact reduction, with no loss of constants.  Affinely
stretch the extreme nodes of any arbitrary set $X\subset[-1,1]$ to
$-1,1$.  Its original Lebesgue constant is at least the norm on its
smaller node hull, which equals the norm of the stretched canonical
set.  For $Y$, stretch the Chebyshev roots by
$1/\cos(\pi/(2N))$.  This affine change preserves all the interval
peak values in \eqref{polyretry:all-intervals}.  Therefore
\eqref{polyretry:comparison} and the already proved finite cotangent
estimate yield the catalogue's exact bound $\Lambda_n\ge B_{n+1}$.
For $N=2$ there is only one canonical set, $\{-1,1\}$, so the
comparison is immediate; the case $N=1$ was handled separately.
The following arguments cover $N\ge3$.

For implementation planning, the comparison theorem has a useful
finite algebraic entry point.  Let $n=N-1\ge2$, let
$t_0=-1<t_1<\cdots<t_n=1$, and define the polynomial $F_i$,
$1\le i\le n$, by its nodal values
\[
 F_i(t_j)=
 \begin{cases}(-1)^{i-1-j},&j<i,\\(-1)^{j-i},&j\ge i.\end{cases}
\]
Then $F_i$ equals the Lebesgue function on $(t_{i-1},t_i)$.
Each $F_i$ has a simple real root in every other nodal gap.  Its
degree is $n-1$ or $n$; in the latter case the remaining root is
real, simple, and outside $[-1,1]$.  These statements follow from
the $n-1$ sign changes, the degree bound, and the fact that $F_i>0$
on its own gap.  Rolle's theorem then gives a unique nondegenerate
maximum $\tau_i$ in its own gap.  The positivity on that gap follows
directly from its cardinal representation; since $n\ge2$, at least
one cardinal function there is negative, so the Lebesgue function
is strictly greater than its endpoint value $1$.

Let $\omega(x)=\prod_{j=0}^n(x-t_j)$ and
$q_i(x)=F_i'(x)/(x-\tau_i)$, so $\deg q_i\le n-2$.
Differentiating the nodal equations and using interpolation
uniqueness gives
\[
 \frac{\partial F_i(x)}{\partial t_j}
       =-F_i'(t_j)\ell_j(x),
 \qquad
 \frac{\partial M_i}{\partial t_j}
       =-F_i'(t_j)\ell_j(\tau_i)
       =\frac{\omega(\tau_i)}{\omega'(t_j)}q_i(t_j)
       \quad(1\le j<n).
\]
Thus every maximal minor of the peak Jacobian is, up to explicit
nonzero diagonal factors, an evaluation determinant for $n-1$
polynomials in the $(n-1)$-dimensional space $\mathcal P_{n-2}$.
The following subsections prove the signs and nonsingularity of
these minors by a finite positive-matrix argument, establish
properness of $t\mapsto(M_{i+1}(t)-M_i(t))_{i=1}^{n-1}$ by
explicit cardinal-product estimates, and supply the global inverse
argument through finite chart covers.  A single radial path then
proves \eqref{polyretry:comparison}.  Thus each step beyond the
Jacobian identity is supplied rather than assumed.

\textbf{Verified local reuse.} Mathlib's
\code{Lagrange.eval_basis_self}, \code{Lagrange.eval_basis_of_ne},
\code{Lagrange.eval_basis_not_at_node},
\code{Lagrange.nodalWeight_eq_eval_derivative_nodal}, and
\code{Lagrange.eval_interpolate_at_node} provide the finite
interpolation identities.  The NLAlib declarations
\code{NLAlib.nodal_chebyshevZero} and
\code{NLAlib.abs_eval_nodal_chebyshevZero_le} already identify the
root grid's nodal polynomial.  Mathlib also supplies
\code{Real.tendsto_harmonic_sub_log} and
\code{Real.eulerMascheroniConstant_lt_two_thirds} in
\code{Mathlib.NumberTheory.Harmonic.EulerMascheroni}.
No statement formalizing the
Bernstein--Erd\H{o}s peak comparison was found in the searched
declaration index.  This is a mathematical handoff, not a claim
that all calculus or topology interfaces have been tested in Lean.

\textbf{Sources and scope.} The comparison theorem and its
interlacing/Jacobian route are in C.~de Boor and A.~Pinkus,
\emph{Proof of the conjectures of Bernstein and Erd\H{o}s concerning
the optimal nodes for polynomial interpolation}, J.~Approx.
Theory 24 (1978), 289--303,
\url{https://doi.org/10.1016/0021-9045(78)90014-X};
author-hosted copy:
\url{https://pinkus.net.technion.ac.il/files/2021/02/conj.pdf}.
The classical Chebyshev local-maximum analysis is due to
L.~Brutman, \url{https://doi.org/10.1137/0715046}.
The argument above instead evaluates explicit interior samples
and supplies its own cotangent and harmonic estimates; it does
not import Brutman's minimum-peak estimate.

% END retry_lebesgue.tex
% BEGIN retry_comparison.tex
\subsection{Closing the finite comparison step: a positive Wronskian certificate}
\label{polyretry:finite-comparison}

The remaining finite determinant calculation admits a useful
simplification: a sum of squares replaces differentiation of
interlacing polynomials, and a matrix with positive off-diagonal
entries replaces the sign-variation argument.  We keep
$n\ge2$, $-1=t_0<\cdots<t_n=1$, and the polynomials $F_i$ and
their unique nondegenerate interval maxima $\tau_i$ from above.
Thus $\tau_1<\cdots<\tau_n$,
$F_i(\tau_i)>1$, and $F_i''(\tau_i)<0$.

\paragraph{A finite rational identity.}
Fix $1\le i<j\le n$, and put
\[
 s=(-1)^{j-i},\qquad
 u=\frac{F_i+sF_j}{2},\qquad
 v=\frac{F_i-sF_j}{2}.
\]
The nodal values are particularly simple:
\[
 \begin{array}{c|cc}
 &k<i\text{ or }k\ge j&i\le k<j\\ \hline
 u(t_k)&F_i(t_k)&0\\
 v(t_k)&0&F_i(t_k).
 \end{array}
\]
The polynomial $v$ vanishes at
$t_0,\ldots,t_{i-1},t_j,\ldots,t_n$ and changes sign in each
of the $j-i-1$ gaps between $t_i,\ldots,t_{j-1}$.
These give exactly $n$ distinct real roots.  Since $\deg v\le n$
and $v(t_i)=1$, its degree is $n$ and all these roots are simple.
Call them $\beta_1<\cdots<\beta_n$; in particular $\beta_1=t_0$.

The polynomial $u$ has the $j-i$ roots
$t_i,\ldots,t_{j-1}$ and a sign-changing root in each gap
$(t_{r-1},t_r)$ with $r<i$ or $r>j$.  These are $n-1$
distinct roots, exactly one between every consecutive pair of
the $\beta_r$.  They are simple and there are no further roots
inside $[t_0,t_n]$.  Here are the multiplicity details: in each
outer gap the endpoint values have opposite signs, so the total
root multiplicity there is odd.  In $(t_{i-1},t_j)$ the endpoint
values of $u$ have sign ratio $(-1)^{j-i}$, so the total root
multiplicity there has parity $j-i$.  We already have $n-1$
distinct roots and degree at most $n$; putting an additional root
or multiplicity in any of these regions would violate its parity.
The other nodes have nonzero values.  Thus any remaining root,
if $\deg u=n$, lies outside $[t_0,t_n]$.

Polynomial interpolation, or division followed by interpolation,
now gives the exact finite partial-fraction identity
\begin{equation}\label{polyretry:positive-residues}
 \frac{u(x)}{v(x)}
   =c+\sum_{r=1}^n\frac{a_r}{x-\beta_r},
 \qquad a_r=\frac{u(\beta_r)}{v'(\beta_r)}>0.
\end{equation}
To verify the signs, $u$ has one simple root between successive
$\beta_r$, so the signs of $u(\beta_r)$ alternate, just as those
of $v'(\beta_r)$ do.  Hence all residues have the same sign.
At $t_i$, precisely $i$ roots of $v$ lie to the left; since
$v(t_i)=1$, its leading coefficient has sign $(-1)^{n-i}$.
Therefore $v'(t_0)$ has sign $(-1)^{i-1}$, which is also the
sign of $u(t_0)=F_i(t_0)$.  This proves strict positivity of
every $a_r$.  The constant $c$ is the ratio of leading
coefficients when $\deg u=n$, and zero otherwise.

Put $v_r(x)=v(x)/(x-\beta_r)$, a polynomial.  Differentiating
\eqref{polyretry:positive-residues} and clearing denominators
gives the polynomial identity
\begin{equation}\label{polyretry:wronskian-sos}
 s\bigl(F_i'(x)F_j(x)-F_i(x)F_j'(x)\bigr)
       =2\sum_{r=1}^n a_r v_r(x)^2>0
       \qquad(x\in\R).
\end{equation}
For the identity it suffices to check points where $v\ne0$,
then use polynomial equality.  Strict positivity holds at all
real points: away from the roots every $v_r$ is nonzero; at
$\beta_r$ the term $v_r(\beta_r)=v'(\beta_r)$ is nonzero.
This is a finite Gram certificate for the Wronskian sign.

\paragraph{Signs of the divided derivative polynomials.}
Recall $q_i(x)=F_i'(x)/(x-\tau_i)$, with $\deg q_i\le n-2$.
Evaluation of \eqref{polyretry:wronskian-sos} at $\tau_j$
and at $\tau_i$, using $F_j(\tau_j),F_i(\tau_i)>0$, gives
\[
 \operatorname{sgn}F_i'(\tau_j)=(-1)^{j-i},\qquad
 \operatorname{sgn}F_j'(\tau_i)=-(-1)^{j-i}
       \quad(i<j).
\]
Dividing by the sign of $\tau_j-\tau_i$ yields
\begin{equation}\label{polyretry:q-signs}
 (-1)^{i+j}q_i(\tau_j)>0\quad(i\ne j),
 \qquad q_i(\tau_i)=F_i''(\tau_i)<0.
\end{equation}
In particular, this calculation does not require a separate
theorem saying that derivatives preserve strict interlacing.

\paragraph{A finite matrix proves rank and a positive normal.}
\label{polyretry:peak-laplacian}
Let
\[
 b_r=\frac1{\prod_{k\ne r}|\tau_r-\tau_k|}>0,\qquad
 B_{ri}=(-1)^{r+i}q_i(\tau_r),\qquad A=\diag(b)B.
\]
The off-diagonal entries of $A$ are strictly positive, and its
diagonal entries are strictly negative.  For every polynomial
$q$ of degree at most $n-2$, comparison of leading coefficients
in its Lagrange formula at the $n$ points $\tau_r$ gives
\[
 \sum_{r=1}^n
 \frac{q(\tau_r)}{\prod_{k\ne r}(\tau_r-\tau_k)}=0.
\]
The denominator has sign $(-1)^{n-r}$.  Applying this identity
to each $q_i$ proves that every column of $A$ sums to zero.

Delete the same row and column $k$ from $A$.  Each remaining
column $i$ satisfies strict diagonal dominance:
\[
 |A_{ii}|=\sum_{r\ne i}A_{ri}
          >\sum_{r\ne i,k}A_{ri}.
\]
The resulting square matrix is invertible.  Indeed its
transpose is strictly row diagonally dominant; for a nonzero
kernel vector choose a coordinate of maximum absolute value,
and the triangle inequality contradicts the strict dominance
in that row.  It follows that any $n-1$ columns of $B$ are
linearly independent, and hence any $n-1$ of the polynomials
$q_1,\ldots,q_n$ are linearly independent.

The same matrix gives a positive nullvector explicitly.
Let $A^{(k)}$ denote its $k$th deleted principal submatrix.
Hold its negative diagonal fixed and multiply its off-diagonal
entries by $s\in[0,1]$.  Strict diagonal dominance persists,
so the determinant never vanishes.  At $s=0$ its sign is
$(-1)^{n-1}$.  Therefore
\begin{equation}\label{polyretry:laplacian-cofactor}
 \pi_k:=(-1)^{n-1}\det A^{(k)}>0\qquad(1\le k\le n).
\end{equation}
Since $A$ has rank $n-1$ and zero column sums, its left
nullspace is the span of $(1,\ldots,1)$.  The identities
$A\operatorname{adj}(A)=\operatorname{adj}(A)A=0$ show that
every row of its adjugate is constant.  The constant in row
$i$ is its diagonal entry $\det A^{(i)}$.  Thus every column
of the adjugate is $(-1)^{n-1}\pi$, and $A\pi=0$.
Undoing the row scalings gives
$\sum_i(-1)^i\pi_iq_i(\tau_r)=0$ for every $r$.
This polynomial has degree at most $n-2$ and $n$ distinct
zeros, so
\begin{equation}\label{polyretry:positive-q-relation}
 \sum_{i=1}^n(-1)^i\pi_iq_i=0.
\end{equation}

Evaluation at the $n-1$ distinct interior nodes
$t_1,\ldots,t_{n-1}$ is an isomorphism on
$\mathcal P_{n-2}$.  Consequently every maximal minor of
$(q_i(t_j))_{1\le i\le n,\,1\le j<n}$ is nonzero.
The explicit diagonal factors in the earlier peak-Jacobian
formula then prove the same fact for
\begin{equation}\label{polyretry:peak-jacobian}
 J_{ij}=\frac{\partial M_i}{\partial t_j}
       =\frac{\omega(\tau_i)}{\omega'(t_j)}q_i(t_j).
\end{equation}
In particular $\rank J=n-1$.  Since
$\operatorname{sgn}\omega(\tau_i)=(-1)^{n-i+1}$,
\eqref{polyretry:positive-q-relation} proves directly that
\[
 \widetilde w_i=\frac{\pi_i}{|\omega(\tau_i)|}>0,
 \qquad \widetilde w^TJ=0.
\]
Indeed
$\widetilde w_i\omega(\tau_i)=(-1)^{n+1}(-1)^i\pi_i$,
with a common sign independent of $i$.  Normalizing gives
the explicit positive normal
\begin{equation}\label{polyretry:positive-normal}
 w_i=\frac{\widetilde w_i}{\sum_r\widetilde w_r}>0,
 \qquad\sum_iw_i=1,\qquad w^TJ=0.
\end{equation}
This argument uses no extremization over node sets.

\paragraph{A quantitative boundary estimate.}
The following properness argument uses only products of
distances.  Let $\delta$ be the smallest node gap and let $y$
be the midpoint of a largest gap.  Since the total length is
$2$, every node has distance at least $1/n$ from $y$.
For $x$ in a gap of length $\delta$, every cardinal numerator
$\prod_{r\ne j}(x-t_r)$ includes at least one endpoint of
that small gap.  That factor has absolute value at most
$\delta$, and each of the other $n-1$ factors is at most $2$.
Cancellation of the cardinal denominators gives
\[
 |\ell_j(x)|\le C_n\delta\,|\ell_j(y)|,
 \qquad C_n=2^{n-1}n^n.
\]
If $M_{\mathrm{s}}$ and $M_{\mathrm{l}}$ are the peaks in
the chosen smallest and largest gaps, respectively, summing
and maximizing yields
\[
 1\le M_{\mathrm{s}}\le C_n\delta M_{\mathrm{l}}.
\]
Whenever $C_n\delta<1$, therefore,
\begin{equation}\label{polyretry:properness}
 \max_i M_i\ge\frac1{C_n\delta},\qquad
 \max_i M_i-\min_i M_i\ge\frac1{C_n\delta}-1.
\end{equation}
Thus $\max_i M_i$ tends to infinity at the boundary of the
open node simplex.  So does the norm of
$\Gamma(t)=(M_{i+1}(t)-M_i(t))_{i=1}^{n-1}$, because the
range of the $M_i$ is at most the sum of the absolute adjacent
differences.  In particular these functions have compact
sublevel sets, and $\Gamma$ is proper.

\paragraph{The local inverse hypotheses.}
The map $M$ is continuously differentiable: the cardinal
coefficients depend smoothly on distinct nodes, and the
nondegenerate maxima $\tau_i$ depend smoothly on the nodes
by the ordinary implicit function theorem.  The cofactor
formula above shows that the positive normalized normal
$w(t)$ is continuous as well.  Finally, $D\Gamma$ is
nonsingular.  If $D\Gamma h=0$, all
coordinates of $Jh$ have the same value $c$; multiplying by
$w^T$ gives $c=0$, and the rank of $J$ gives $h=0$.
The finite-dimensional inverse function theorem proves
that $\Gamma$ is a local diffeomorphism.  Together with
properness, these are exactly the hypotheses of the elementary
global argument in the next subsection.

\textbf{Scope of this improvement.} All finite polynomial,
matrix, properness, and normal-vector assertions in this
subsection have been proved.  The architecture comes from
the de Boor--Pinkus comparison theorem cited above; the
positive-residue Wronskian and strictly diagonally dominant
matrix provide an explicit alternative to its derivative
interlacing and polynomial sign-variation calculations.

\textbf{Additional verified reuse.} The strict column dominance
step already exists as
\code{det_ne_zero_of_sum_col_lt_diag} in
\code{Mathlib/LinearAlgebra/Matrix/Gershgorin.lean}.
The adjugate identities are \code{Matrix.mul_adjugate} and
\code{Matrix.adjugate_mul} in
\code{Mathlib/LinearAlgebra/Matrix/Adjugate.lean}.
The polynomial root-count conclusion can use
\code{Polynomial.eq_zero_of_natDegree_lt_card_of_eval_eq_zero}.
These are checked source declarations, not a claim that the
assembled argument has already been implemented in Lean.

% END retry_comparison.tex
% BEGIN retry_lebesgue_global.tex
\subsection{A global peak comparison from one radial path}
\label{polyretry:global-comparison}

The global step can be isolated as the following finite-dimensional
calculus lemma.  It is useful because the desired comparison follows
from one radial path; the stronger global Pareto theorem for all peak
vectors is unnecessary.

\begin{lemma}[Positive normal and proper peak differences]
Let $U\subset\R^d$ be nonempty, open, and convex, and let
$M:U\to\R^{d+1}$ be continuously differentiable.  Suppose, for every
$u\in U$, that $DM(u)$ has rank $d$ and its left nullspace contains a
vector with all coordinates strictly positive.  Set
\[
 \Gamma(u)=(M_{i+1}(u)-M_i(u))_{i=1}^{d}.
\]
Suppose also that $\Gamma$ is proper: the inverse image in $U$ of
every compact subset of $\R^d$ is compact.  Then there is a unique
$u_*\in U$ at which all coordinates of $M$ are equal.  If their
common value is $\lambda_*$, then for every $u\in U$,
\[
 \min_i M_i(u)\le\lambda_*\le\max_i M_i(u).
\]
Consequently $\min_i M_i(v)\le\max_i M_i(u)$ for all $u,v\in U$.
\end{lemma}

\begin{proof}
Normalize the positive left nullvector so that
$\sum_i w_i(u)=1$.  It is unique, since the left nullspace has
dimension one, and is continuous: locally choose a nonzero maximal
minor and solve the linear equations; the positivity fixes its
normalization.  Equivalently, use the signed maximal minors and
divide their positive representative by its coordinate sum.

If $D\Gamma(u)v=0$, then $DM(u)v=c(1,\ldots,1)$ for a scalar $c$.
Multiplication by $w(u)^T$ gives $0=c$.  The rank assumption then
gives $v=0$.  Thus $D\Gamma(u)$ is invertible.  The local inverse
theorem makes $\Gamma$ a local diffeomorphism.  The elementary
proper-map lemma proved below makes it a global diffeomorphism
$U\to\R^d$.  In particular $u_*:=\Gamma^{-1}(0)$ exists and is unique.

Fix $u\in U$ and put $h_i=M_i(u)$.  Lift the radial segment by
\[
 u(s)=\Gamma^{-1}(s\Gamma(u)),\qquad 0\le s\le1.
\]
Every adjacent difference of the vector $M(u(s))-s h$ is zero.
Hence for a continuously differentiable scalar function $c$,
\[
 M_i(u(s))=c(s)+s h_i\quad\hbox{for every }i,
 \qquad c(1)=0,\qquad c(0)=\lambda_*.
\]
Differentiate this identity.  Since $w(u(s))^TDM(u(s))=0$,
\[
 0=\sum_iw_i(u(s))\frac{d}{ds}M_i(u(s))
   =c'(s)+\sum_iw_i(u(s))h_i.
\]
The fundamental theorem of calculus now gives the exact identity
\begin{equation}\label{polyretry:peak-average}
 \lambda_*=\int_0^1\sum_iw_i(u(s))h_i\,ds.
\end{equation}
The integrand is a convex combination of the original peak values
$h_i$, proving both inequalities.  In fact both inequalities are
strict unless all $h_i$ are equal, because every weight is positive.
\end{proof}

For the interpolation problem, $d=n-1$ and
$U=\{(t_1,\ldots,t_{n-1}):-1<t_1<\cdots<t_{n-1}<1\}$.
Thus the determinant and collision estimates are exactly the local
and properness hypotheses of this lemma.  The proof introduces no
additional interpolation-theoretic assumption.

\subsection{The proper-map lemma without importing covering-space theory}
\label{polyretry:proper-inverse}

\begin{lemma}
Let $U\subset\R^d$ be nonempty, open, and convex.  Every proper
local homeomorphism $f:U\to\R^d$ is a global homeomorphism.
If $f$ is a local $C^1$ diffeomorphism, its global inverse is $C^1$.
\end{lemma}

\begin{proof}
The image is open by local invertibility.  It is closed as well:
if $f(x_j)\to y$, the set consisting of this convergent sequence
and its limit is compact.  Properness gives a subsequence of $x_j$
converging to a point $x\in U$, and continuity gives $f(x)=y$.
Connectedness of $\R^d$ therefore proves surjectivity.

We next give all details behind lifting a contracted loop.
For any $y$, its fiber $f^{-1}(y)$ is compact by properness and
discrete by local injectivity, and hence finite.  Write it as
$\{x_1,\ldots,x_r\}$.  Choose pairwise disjoint inverse-chart
neighborhoods $V_j$ of $x_j$.  There is a neighborhood $W$ of $y$
such that every point of $f^{-1}(W)$ lies in $\bigcup_jV_j$.
Otherwise choose $y_k\to y$ and $z_k\notin\bigcup_jV_j$ with
$f(z_k)=y_k$.  Properness gives a convergent subsequence whose limit
belongs to the fiber, contradicting the openness of its chart.
Shrink $W$ into every chart image.  It follows that $f^{-1}(W)$
is a disjoint union of $r$ open pieces, each carried homeomorphically
onto $W$.  This conclusion has been proved directly from properness.

Suppose now $f(x)=f(z)=y$.  Convexity supplies the path
$a(s)=(1-s)x+sz$ in $U$.  Its image is a loop based at $y$.
Contract that loop using the explicitly defined map
\[
 H(s,t)=(1-t)f(a(s))+ty,\qquad (s,t)\in[0,1]^2.
\]
The bottom edge $H(s,0)$ has the prescribed lift $a(s)$.
Cover the compact square by inverse images under $H$ of the
neighborhoods $W$ constructed above.  The Lebesgue-number lemma
gives a finite rectangular grid fine enough that the image of
each closed cell lies in one such $W$.

Extend the lift over the cells, from left to right in each row and
then from bottom to top.  For a cell whose bottom edge is already
lifted, choose the inverse branch containing the lifted bottom-left
corner and compose that branch with $H$ throughout the cell.
It agrees with the existing bottom edge: both lifts lie in
$f^{-1}(W)$, the edge is connected, and its starting point determines
the unique open piece.  The same argument gives agreement on the
left edge whenever that edge has already been filled.  Thus finite
pasting gives a continuous lift $\widetilde H$ on the full square,
with $\widetilde H(s,0)=a(s)$.

Both vertical edges of $H$ and its top edge are constant equal to
$y$.  Their lifts are continuous maps into the finite fiber, and
are therefore constant on each edge.  The left lifted edge is
constantly $x$ and the right lifted edge constantly $z$.  The
constant top edge connects these values, so $x=z$.  This proves
injectivity.  A bijective local homeomorphism is a homeomorphism,
and the differentiability of its inverse follows locally from
the inverse charts.
\end{proof}

The proof uses only finite chart covers, compactness, connectedness
of intervals, and a finite grid on a square.  The existing declarations
\code{IsClosedMap.isCoveringMapOn_of_isLocalHomeomorphOn} in
\code{Mathlib/Topology/Covering/Basic.lean} and
\code{lebesgue_number_lemma_of_metric} support these steps if a
formalizer prefers established topological interfaces.  The argument
above supplies the global step mathematically rather than leaving
the global inverse theorem as an unexplained imported assumption.

% END retry_lebesgue_global.tex
% BEGIN retry_abs.tex
\section{Second attempt: Bernstein's polynomial limit from coefficient compactness}
\label{polyretry:bernstein}

\paragraph{New result and remaining boundary.}
The polynomial half of \code{abs-best-approx-asymptotics} can in fact be
completed.  The proof below establishes
\[
 \lim_{n\to\infty}nE_n(|x|;[-1,1])=\beta,
 \qquad \frac1{18\pi}\le\beta\le\frac2\pi.
\]
It does not certify the decimal digits traditionally printed for
$\beta$, and it does not prove Stahl's rational prefactor $8$.
The useful change is to subtract an explicit comparison polynomial
\emph{before} passing to a limit.  The resulting bounded polynomial
sequence can be compactified coefficient by coefficient.  No
Phragm\'en--Lindel\"of theorem, general normal-family theorem, or library
of entire functions of exponential type is required.

The classical connection with entire approximation is discussed in
\href{https://lubinsky.math.gatech.edu/Research%20papers/BrnstnDec05CA.pdf}
{Lubinsky, \emph{On the Bernstein Constants of Polynomial Approximation}}.
The detailed argument here uses the explicit Chebyshev comparison
sequence to provide all growth estimates needed for this special case.

\subsection{A coefficient compactness lemma with an elementary proof}
\label{polyretry:coefficient-compactness}

\begin{lemma}
Let $D_j$ be complex polynomials.  Suppose that for every $R>0$ there
are a finite $M_R$ and an index $j_R$ such that
\[
 \max_{|z|\le R}|D_j(z)|\le M_R\qquad(j\ge j_R).
\]
There is a subsequence converging uniformly on every compact subset of
$\C$ to an entire function.  Real and even coefficients, when present
in every $D_j$, are preserved.
\end{lemma}

\begin{proof}
Write $D_j(z)=\sum_{k\ge0}d_{jk}z^k$, extending its finite coefficient
sequence by zeros.  Finite exponential orthogonality on a circle gives,
for $R>0$,
\[
 d_{jk}R^k=\frac1{2\pi}\int_{-\pi}^{\pi}
 D_j(Re^{it})e^{-ikt}\,dt.
\]
Consequently $|d_{jk}|\le M_RR^{-k}$ eventually in $j$.
Every coefficient sequence is bounded, since deleting a finite prefix
does not affect boundedness.  Successive applications of
Bolzano--Weierstrass, followed by a diagonal subsequence, give
$d_{j_\ell k}\to d_k$ for every $k$.

Fix $0<r<R$.  The eventual coefficient bound passes to the limit and
gives both
\[
 \sum_{k>K}|d_k|r^k
 \le\frac{M_R(r/R)^{K+1}}{1-r/R},\qquad
 \sum_{k>K}|d_{j_\ell k}|r^k
 \le\frac{M_R(r/R)^{K+1}}{1-r/R}
\]
for sufficiently large $\ell$.  The finite initial sums converge
coefficientwise.  Thus $D_{j_\ell}\to D:=\sum_{k\ge0}d_kz^k$ uniformly
on $|z|\le r$.  Since $R$ is arbitrary, this power series has infinite
radius of convergence.  Its coefficients retain reality and parity.
\end{proof}

This lemma only uses compactness of bounded subsets of $\R^2$, finite
Fourier orthogonality, and geometric-series tails.  Equivalently, one
may apply sequential compactness to the countable product of closed
disks containing the coefficients; no compactness theorem for spaces
of holomorphic functions is needed.

\subsection{The comparison polynomials and their explicit local limit}
\label{polyretry:comparison-limit}

For an even integer $n=2m\ge2$, take the degree-$n$ Chebyshev truncation
\[
 S_n(x)=\frac2\pi+
 \frac4\pi\sum_{j=1}^{m}
 \frac{(-1)^{j+1}}{4j^2-1}T_{2j}(x),
 \qquad B_n(z)=nS_n(z/n).
\]
The already derived absolute-value series gives
\begin{equation}\label{polyretry:comparison-real-error}
 \|S_n-|x|\|_{[-1,1]}=\frac{2}{\pi(n+1)},
 \qquad
 \|B_n-|x|\|_{[-n,n]}=\frac{2n}{\pi(n+1)}<\frac2\pi.
\end{equation}
The same polynomial formulas are used over $\R$ and $\C$.

Only a \emph{local} inverse of complex sine near zero is useful in the
following calculation.  Let $a$ be holomorphic near zero, with
\[
 a(0)=0,\qquad \sin(a(u))=u,\qquad a'(0)=1.
\]
Its existence follows from the ordinary local inverse theorem, since
$\sin'(0)=1$.  One may instead use
$a(u)=-i\log(\sqrt{1-u^2}+iu)$ on sufficiently small neighborhoods of
$u=0$, with the square root near $1$ and logarithm near $1$.
No choice of a global inverse-trigonometric branch is involved.
Differentiability at zero implies
\begin{equation}\label{polyretry:local-inverse-scaling}
 a_n(z):=n a(z/n)\longrightarrow z
 \quad\hbox{uniformly on every fixed compact subset of $\C$}.
\end{equation}
Indeed, if $|a(u)-u|\le\varepsilon|u|$ for $|u|\le\delta$, then
$|a_n(z)-z|\le\varepsilon R$ for $|z|\le R$ and $n\ge R/\delta$.

The recurrence for $T_k$ gives
$T_{2j}(\sin a)=(-1)^j\cos(2ja)$.
Using $\sum_{j=1}^{m}(4j^2-1)^{-1}=m/(2m+1)$, rewrite $B_n$ as
\begin{equation}\label{polyretry:comparison-riemann}
 B_n(z)=\frac{2n}{\pi(n+1)}+
 \frac{4n}{\pi}\sum_{j=1}^{n/2}
 \frac{1-\cos(2j a(z/n))}{4j^2-1}.
\end{equation}
Define the continuous function at its removable value by
\[
 \Phi(t,z)=
 \begin{cases}(1-\cos(tz))/t^2,&t\ne0,\\z^2/2,&t=0.\end{cases}
\]
An identity that also proves this continuity is
\[
 \Phi(t,z)=z^2\int_0^1(1-s)\cos(stz)\,ds.
\]
Writing $t_j=2j/n$, the sum in
\eqref{polyretry:comparison-riemann} equals
\[
 \frac2\pi\,\frac2n\sum_{j=1}^{n/2}
 \Phi(t_j,a_n(z))\frac{4j^2}{4j^2-1}.
\]
The extra weights can be removed uniformly on a fixed compact set:
if $|\Phi(t,a_n(z))|\le M$ there, their total contribution is at most
\[
 \frac{4M}{\pi n}\sum_{j=1}^{n/2}\frac1{4j^2-1}
 \le\frac{2M}{\pi n}.
\]
Uniform continuity of $\Phi$ on the relevant compact set, together with
\eqref{polyretry:local-inverse-scaling}, then reduces the remaining sum
to an ordinary Riemann sum, uniformly in $z$ on compact sets.
Therefore
\begin{equation}\label{polyretry:comparison-entire}
 B_n(z)\longrightarrow
 B(z):=\frac2\pi+\frac2\pi\int_0^1\frac{1-\cos(tz)}{t^2}\,dt
 \quad\hbox{locally uniformly in $\C$}.
\end{equation}
The quotient in the integral always means its continuous extension at
$t=0$.  Termwise integration of the locally uniformly convergent cosine
series gives the explicit entire power series
\[
 B(z)=\frac2\pi+
 \frac2\pi\sum_{k=1}^{\infty}
 \frac{(-1)^{k+1}z^{2k}}{(2k)!(2k-1)}.
\]
Moreover $|\cos w|\le e^{|\operatorname{Im}w|}$ and the integral formula
for $\Phi$ imply the global estimate
\begin{equation}\label{polyretry:comparison-growth}
 |B(z)|\le\frac2\pi+
 \frac{|z|^2}{\pi}e^{|\operatorname{Im}z|}
 \le\frac{2+|z|^2}{\pi}e^{|\operatorname{Im}z|}.
\end{equation}

\subsection{Subtract first, then take a coefficient limit}
\label{polyretry:scaled-extremal}

Let $E_n=E_n(|x|;[-1,1])$.  Symmetrizing a real polynomial cannot
increase its approximation error to $|x|$.  Therefore
\[
 E_{2m+1}=E_{2m},
\]
and for even $n$ choose an even real minimizer $p_n$.  Define
\[
 b=\liminf_{\substack{n\to\infty\\n\text{ even}}}nE_n,
 \qquad P_n(z)=n p_n(z/n),\qquad D_n=P_n-B_n.
\]
The earlier finite bounds imply
\[
 \frac1{18\pi}\le b\le\frac2\pi.
\]
Along some even subsequence $n_j\to\infty$, $n_jE_{n_j}\to b$.
For every even $n$, the comparison bound and optimality give
\begin{equation}\label{polyretry:difference-interval}
 \deg D_n\le n,\qquad
 \|D_n\|_{[-n,n]}
 \le nE_n+\frac{2n}{\pi(n+1)}<\frac4\pi.
\end{equation}

The finite polynomial Bernstein--Walsh inequality
\eqref{polyroot:walsh} has a particularly useful scaling limit.
For a fixed $z$, take $n$ large enough that $a(z/n)$ is defined.
The two numbers
\[
 w_\pm=\exp\bigl(\pm i(\pi/2-a(z/n))\bigr)
\]
satisfy $J(w_\pm)=z/n$, and the larger modulus is
$\exp(|\operatorname{Im}a(z/n)|)$.  Apply
\eqref{polyroot:walsh} to the polynomial $u\mapsto D_n(nu)$:
\begin{equation}\label{polyretry:difference-complex}
 |D_n(z)|\le\frac4\pi
 \exp\bigl(n|\operatorname{Im}a(z/n)|\bigr).
\end{equation}
This includes real $z/n$ in the interval, where the modulus is $1$.
On $|z|\le R$, the right side is eventually at most
$(4/\pi)e^{R+1}$, by \eqref{polyretry:local-inverse-scaling}.

Apply the coefficient compactness lemma to $D_{n_j}$, taking a further
subsequence if necessary.  It gives an entire even real function $D$.
For each fixed $z$, pass to the limit in
\eqref{polyretry:difference-complex} and use
\eqref{polyretry:local-inverse-scaling}:
\[
 |D(z)|\le\frac4\pi e^{|\operatorname{Im}z|}.
\]
Consequently $F=B+D$ is entire, even, real on the real axis, and
\begin{equation}\label{polyretry:extremal-growth}
 |F(z)|\le\frac{6+|z|^2}{\pi}e^{|\operatorname{Im}z|}.
\end{equation}
For every fixed real $x$, eventually $|x|\le n_j$, and then
\[
 |P_{n_j}(x)-|x||\le n_jE_{n_j}.
\]
Since $P_{n_j}=B_{n_j}+D_{n_j}\to F$ locally uniformly, we obtain
\begin{equation}\label{polyretry:global-real-error}
 |F(x)-|x||\le b\qquad(x\in\R).
\end{equation}
The unbounded real line causes no compactness issue: this last
inequality is obtained for each fixed $x$ after the subsequence has
already been chosen on all complex compact sets.

\subsection{A small contraction turns the liminf into the full limit}
\label{polyretry:contraction}

Fix $0<\alpha<1$ and put
\[
 t=\frac{1-\alpha}{2},\qquad \rho=e^t>1,
 \qquad c=t(1-e^{-t})>0.
\]
Since $\alpha\le e^{\alpha-1}=e^{-2t}$ and
$\sinh t\le te^t$, we have
\begin{equation}\label{polyretry:strict-exponent}
 \alpha\sinh t-t\le -c<0.
\end{equation}
For any integer $N\ge1$, consider the entire function
\[
 G_N(w)=\frac{F(\alpha Nw)}{\alpha N}.
\]
For real $w\in[-1,1]$, \eqref{polyretry:global-real-error} gives
\[
 |G_N(w)-|w||\le\frac{b}{\alpha N}.
\]
On the closed Bernstein ellipse with parameter $\rho=e^t$,
\[
 |w|\le\cosh t,\qquad
 |\operatorname{Im}w|\le\sinh t.
\]
The latter follows directly from
$J(re^{i\theta})=\tfrac12(r+r^{-1})\cos\theta+
\tfrac i2(r-r^{-1})\sin\theta$, $1\le r\le\rho$.
Thus \eqref{polyretry:extremal-growth} bounds $G_N$ there by
\[
 M_N=\frac{6+(\alpha N\cosh t)^2}{\pi\alpha N}
       e^{\alpha N\sinh t}.
\]
The already available analytic Chebyshev truncation theorem supplies a
real polynomial $q_N$ of degree at most $N$, satisfying
\[
 \|G_N-q_N\|_{[-1,1]}
 \le\frac{2M_N\rho^{-N}}{\rho-1}
 \le\frac{2\bigl(6+(\alpha N\cosh t)^2\bigr)}
          {\pi\alpha N(e^t-1)}e^{-cN}.
\]
Therefore the following estimate holds for every $N\ge1$:
\begin{equation}\label{polyretry:liminf-to-limsup}
 NE_N\le\frac b\alpha+
 \frac{2\bigl(6+(\alpha N\cosh t)^2\bigr)}
      {\pi\alpha(e^t-1)}e^{-cN}.
\end{equation}
Its second term tends to zero.  Hence
$\limsup_{N\to\infty}NE_N\le b/\alpha$.
Letting $\alpha\uparrow1$ gives $\limsup NE_N\le b$.
Finally, $E_{2m+1}=E_{2m}$ gives
$(2m+1)E_{2m+1}-2mE_{2m}=E_{2m}\to0$.
Thus the full sequence has the same liminf $b$ as the even sequence.
It follows that $NE_N\to b$.  Set $\beta=b$.

This proves the polynomial asymptotic $E_n\sim\beta/n$, since
$\beta\ge1/(18\pi)>0$.  In dimension zero the matrix consequence is
vacuous; for a nonempty Hermitian matrix with spectrum in $[-1,1]$,
the usual scalar-to-operator transfer applies to every polynomial
$q_N$ above.  No dimension-dependent constant is introduced.

\subsection{An exact variational characterization with an explicit growth bound}
Define $\mathcal F$ to consist of all entire functions $H$ that are
even, real on $\R$, and satisfy
\[
 |H(z)|\le\frac{6+|z|^2}{\pi}e^{|\operatorname{Im}z|}
 \qquad(z\in\C).
\]
The proof also yields
\begin{equation}\label{polyretry:bernstein-variational}
 \beta=\min_{H\in\mathcal F}\sup_{x\in\R}|H(x)-|x||.
\end{equation}
Indeed, the constructed $F$ belongs to $\mathcal F$ and has error at
most $\beta$.  For any other $H\in\mathcal F$ with finite error $A$,
the contraction proof, with $A$ replacing $b$, gives $\beta\le A$.
Functions with infinite error do not affect the assertion.  This
characterization is optional; the limit proof itself needs no
definition of exponential type or any theory of best entire
approximants.

\subsection{Reusable interfaces and the precise residual gap}
The following existing declarations are relevant; availability here
means their source was inspected, not that a Lean implementation of
this proof was checked.
\begin{itemize}
\item
\code{NLAlib.exists_degree_le_abs_sub_eval_le_of_open_bernsteinEllipse}
and
\code{NLAlib.abs_re_sub_eval_chebyshevAnalyticTrunc_le_of_open_ellipse}
provide the final approximation estimate directly.
\item
\code{Polynomial.fourierCoeff_toAddCircle_natCast}, in
\code{Mathlib/Analysis/Polynomial/Fourier.lean}, identifies polynomial
coefficients with finite-circle Fourier coefficients.  Apply it to
$D_j(RX)$ to obtain the coefficient bounds used above.
\item
\code{IsCompact.tendsto_subseq} and
\code{CompactSpace.tendsto_subseq}, in
\code{Mathlib/Topology/Sequences.lean}, supply sequential compactness.
The countable product of coefficient disks is an alternative to an
explicit diagonal-subsequence construction.
\item
\code{HasStrictFDerivAt.localInverse} and
\code{HasStrictFDerivAt.to_localInverse}, in
\code{Mathlib/Analysis/Calculus/InverseFunctionTheorem/FDeriv.lean},
support the local sine inverse.  Its use here is only local near zero;
a global complex arcsine API is unnecessary.
\item
\code{Complex.hasSum_cos}, in
\code{Mathlib/Analysis/SpecialFunctions/Trigonometric/Series.lean},
supports the explicit series for $B$.  Uniform convergence of the
parameter-dependent Riemann sums follows from uniform continuity on a
compact rectangle and the elementary estimate by the largest
oscillation in a partition cell.
\end{itemize}
The additional mathematical lemmas to package are the coefficient
compactness lemma, the comparison limit
\eqref{polyretry:comparison-entire}, and the scaled growth estimate
\eqref{polyretry:difference-complex}.  They have been proved above;
their Lean interfaces still require work.

The exact rational statement
$e^{\pi\sqrt n}E_{n,n}(|x|)\to8$ is a separate conclusion of the
catalogue entry.  The even-rational reduction to approximation of
$\sqrt t$ remains valid, but merely importing Stahl's square-root
theorem does not constitute another derivation.  The completed
Bernstein proof therefore removes the polynomial gap while leaving
that rational prefactor explicitly unresolved in this handoff.

% END retry_abs.tex
% BEGIN retry_rational_rate.tex
\section{A finite-product proof of the sharp rational exponent for absolute value}
\label{retryroot:sharp-rate}
This section proves a sharper intermediate conclusion than
\eqref{polyroot:newman-simple}. If
$\mathcal E_n=E_{n,n}(|x|;[-1,1])$, then
\begin{equation}\label{retryroot:root-limit}
 \lim_{n\to\infty}\frac{-\log\mathcal E_n}{\sqrt n}=\pi,
 \qquad
 \lim_{n\to\infty}\mathcal E_n^{1/\sqrt n}=e^{-\pi}.
\end{equation}
The argument uses finite products, a clipped scalar kernel, and elementary
quadrature estimates. It does not give the stronger multiplicative
constant $8$: that would require
$-\log\mathcal E_n=\pi\sqrt n-\log8+o(1)$, whereas the result below
controls the remainder only by $o(\sqrt n)$.

\subsection{One scalar kernel and a lower bound for every degree}
Set
\[
 \kappa(s)=\log\frac{1+e^{-|s|}}{1-e^{-|s|}}\quad(s\ne0),
 \qquad \kappa(0)=+\infty.
\]
It is even, nonnegative, decreasing on $(0,\infty)$, and
\begin{equation}\label{retryroot:kernel-mass}
 \int_{\R}\kappa(s)\,ds=\frac{\pi^2}{2}.
\end{equation}
Indeed on $s>0$ the nonnegative geometric-logarithmic expansion is
$\kappa(s)=2\sum_{j\ge0}e^{-(2j+1)s}/(2j+1)$.
Integrating its finite partial sums and passing monotonically to the
limit gives $2\sum_{j\ge0}(2j+1)^{-2}=\pi^2/4$ on the positive half-line.
Also $\int |s|\kappa(s)\,ds<\infty$, from the logarithmic singularity at
zero and exponential decay at infinity.

Let $p/q$ be the best type-$(n,n)$ approximant, in reduced form, with
$q>0$ on $[-1,1]$. Existence and uniqueness were proved in
Section~\ref{polybest:rational}. Its evenness makes both $p$ and $q$
even after normalizing $q(0)=1$. Write $\varepsilon=\mathcal E_n$.
The case $\varepsilon\ge1$ is harmless, and $\varepsilon=0$ is
impossible: equality with $x$ on $(0,1)$ would force the polynomial
identity $p=xq$, contradicting evenness. Thus assume $0<\varepsilon<1$.
Put $P=p+xq$, so $\deg P\le n+1$ and
\[
 \phi(x)=\frac{P(-x)}{P(x)}
 =\frac{p(x)/q(x)-x}{p(x)/q(x)+x}.
\]
For $\varepsilon<x<1$, the denominator is positive and
$|\phi(x)|\le\varepsilon/x$. Consequently
\[
 \int_\varepsilon^1\log\frac1{|\phi(x)|}\,\frac{dx}{x}
 \ge\int_\varepsilon^1\log\frac{x}{\varepsilon}\,\frac{dx}{x}
 =\frac12\log^2(1/\varepsilon).
\]
Factor $P$ over $\C$, counting multiplicity. A root $-a$ contributes
$\log|(x+a)/(x-a)|$. This contribution is nonpositive if $\Re a\le0$;
if $\Re a>0$, it is nonnegative and at most
$\log|(x+|a|)/(x-|a|)|=\kappa(\log x-\log|a|)$.
Its integral over any positive subinterval, with measure $dx/x$, is
therefore at most $\pi^2/2$. A root at zero contributes zero.
The finitely many logarithmic singularities are integrable; one may
delete small intervals around them and then take limits. Summing over
at most $n+1$ roots proves
\begin{equation}\label{retryroot:rational-lower}
 \boxed{\mathcal E_n\ge e^{-\pi\sqrt{n+1}}}\qquad(n\ge0).
\end{equation}
In particular $\limsup(-\log\mathcal E_n)/\sqrt n\le\pi$.

\subsection{A deterministic finite product with the matching exponent}
The following construction is deliberately explicit. Let $v\ge4$ be
an integer and put
\[
 m=v^2,\quad B=2\log v,\quad h=v^{-1},\quad K=\tfrac12\log v,
 \quad c=2/\pi^2,\quad J=\left\lceil\frac{m+B}{h}\right\rceil.
\]
On the grid $u_j=-B+jh$ ($0\le j<J$) assign the nonnegative integer
multiplicity
\[
 a_j=\lfloor h c(m-u_j)\rfloor.
\]
Every $u_j<m$. Form the polynomial
\begin{equation}\label{retryroot:taper-product}
 P_v(x)=
 \prod_{j=0}^{J-1}(x+e^{-u_j})^{a_j}
 \prod_{k=0}^{m}(x+e^{-k}),
 \qquad
 R_v(x)=x\frac{P_v(x)-P_v(-x)}{P_v(x)+P_v(-x)}.
\end{equation}
The second product has only $m+1$ factors. Its role is to keep a
denominator estimate uniform near the end of the interval.
All factors have positive parameters, so
$P_v(x)>|P_v(-x)|$ for $x>0$; the even denominator is nonzero also at
zero. Thus $R_v$ has no poles on the real interval and is even.
If $D_v=\deg P_v$, its rational type is at most $(D_v+1,D_v+1)$.

For a nonnegative function $f$ of bounded variation on a uniform grid,
the elementary cellwise estimate gives
\begin{equation}\label{retryroot:riemann-bv}
 \left|h\sum_{j=0}^{J-1}f(u_j)
       -\int_{-B}^{-B+Jh}f(u)\,du\right|
 \le h\operatorname{Var}(f).
\end{equation}
To prove it, bound the error on each cell by $h$ times the oscillation
there, then sum the oscillations, which are bounded by total variation.
Applying the decreasing-function version to $c(m-u)_+$ gives
\begin{equation}\label{retryroot:product-degree}
 D_v+1\le
 \frac{(m+B)^2}{\pi^2}+\frac{2h(m+B)}{\pi^2}+m+2.
\end{equation}
Indeed the sum of the floor multiplicities is no larger than the left
Riemann sum for that decreasing function; its integral is
$c(m+B)^2/2$, and the left-sum error is at most $hc(m+B)$.

To control the error, clip the singular kernel:
$\kappa_K=\min(\kappa,K)$, with value $K$ at zero.
Its variation on the real line is $2K$. Moreover,
\begin{equation}\label{retryroot:kernel-truncation}
 \int_\R(\kappa-\kappa_K)\le8e^{-K}\qquad(K\ge\log2).
\end{equation}
For $s>0$ use $\kappa(s)\le\log(1+2/s)$. The positive part of
$\log(1+2/s)-K$ is supported on $0<s<2/(e^K-1)$ and has integral
$-2\log(1-e^{-K})$. Doubling and using
$-\log(1-z)\le2z$ for $0\le z\le1/2$ proves the estimate.

For $t\ge0$ consider the finite sum
$S_v(t)=\sum_{j:a_j>0} a_j\kappa(t-u_j)$, allowing value $+\infty$.
Zero multiplicities are omitted, so no product $0\cdot\infty$ occurs.
Use $a_j\ge hc(m-u_j)-1$ and $\kappa\ge\kappa_K$.
For $\rho(u)=c(m-u)_+$ on $[-B,-B+Jh]$, the product
$\rho(u)\kappa_K(t-u)$ has variation at most $3c(m+B)K$:
use the product variation inequality, $\sup\rho=c(m+B)$,
$\operatorname{Var}\rho=c(m+B)$, and
$\operatorname{Var}\kappa_K\le2K$.
Applying \eqref{retryroot:riemann-bv} twice yields
\begin{align*}
 S_v(t)&\ge
 \int_{-B}^{m}\rho(u)\kappa_K(t-u)\,du
 -3hc(m+B)K-\frac{\pi^2}{2h}-2K.
\end{align*}
Here the second application bounds
$\sum_j\kappa_K(t-u_j)\le\pi^2/(2h)+2K$.

The full, untruncated convolution satisfies
\[
 \int_{-\infty}^{m}c(m-u)\kappa(t-u)\,du\ge m-t.
\]
Extend its density as $c(m-u)_+$ to the entire real line, bound it
below by the affine function $c(m-u)$, and use the mass identity
\eqref{retryroot:kernel-mass} and the zero first moment of the even
kernel. Both sides are integrable, including the signed affine term.
For the omitted part $u<-B$, use
$\kappa(t-u)\le4e^{u-t}$ (because $t-u\ge B\ge\log2$) to bound its
integral by $4ce^{-B}(m+B+1)$. On $[-B,m]$, clipping loses at most
$8c(m+B)e^{-K}$. Therefore
\begin{equation}\label{retryroot:potential-lower}
 S_v(t)\ge m-t-A_v,
\end{equation}
where the entirely explicit remainder is
\[
 A_v=4ce^{-B}(m+B+1)+8c(m+B)e^{-K}
      +3hc(m+B)K+\frac{\pi^2}{2h}+2K.
\]
In particular $A_v/m\to0$: the largest term is of order $v^{3/2}$,
whereas $m=v^2$.

For $x=e^{-t}$ with $0\le t\le m$, one of the additional integer
grid parameters $k\in\{0,\ldots,m\}$ satisfies $|t-k|\le1$.
The corresponding factor in $|P_v(-x)/P_v(x)|$ is at most
$\tanh(1/2)<1/2$; all other factors are at most one. Thus with
$\phi=P_v(-x)/P_v(x)$,
\[
 |x-R_v(x)|=\frac{2x|\phi|}{|1+\phi|}
 \le4x|\phi|\le4e^{-t-S_v(t)}
 \le4e^{-m+A_v}.
\]
For $0\le x\le e^{-m}$, every factor of $P_v(-x)$ is nonnegative,
so $0\le R_v(x)\le x$ and the error is at most $e^{-m}$.
Evenness now proves the uniform estimate
\begin{equation}\label{retryroot:rational-upper}
 \boxed{\||x|-R_v(x)\|_{[-1,1]}\le4e^{-v^2+A_v}}.
\end{equation}

For completeness this yields a full limit over all degrees, not just
a subsequence. Since $2\log v\le v$ for $v\ge4$, the degree bound
\eqref{retryroot:product-degree} is at most
$n_v=\lceil v^4/\pi^2+10v^3\rceil$.
The integers $n_v$ increase, $n_v/v^4\to1/\pi^2$, and
$n_{v+1}/n_v\to1$. For $n_v\le n<n_{v+1}$, monotonicity gives
$\mathcal E_n\le4e^{-v^2+A_v}$. For sufficiently large $v$,
$v^2-A_v-\log4>0$. Dividing this lower bound for the negative logarithm by
$\sqrt{n_{v+1}}$ proves
$\liminf(-\log\mathcal E_n)/\sqrt n\ge\pi$.
Together with \eqref{retryroot:rational-lower} this proves
\eqref{retryroot:root-limit}.

\paragraph{Operator transfer and formalization boundary.}
For a Hermitian contraction $A$, the same rational function satisfies
$\||A|-R_v(A)\|\le4e^{-v^2+A_v}$, with no dependence on the matrix
dimension. The proof needs finite products and scalar integrals, not
the existence of an equilibrium measure or an infinite-dimensional
Hankel operator. Complex polynomial factorization can use
\code{Polynomial.C_leadingCoeff_mul_prod_multiset_X_sub_C} and the
fundamental theorem of algebra. The kernel mass uses
\code{hasSum_zeta_two} from
\code{Mathlib/NumberTheory/ZetaValues.lean}; split that nonnegative
series into its even and odd terms. The bounded-variation quadrature
estimate has its proof above.
These are proof contracts, not claims that all connecting Lean lemmas
already exist. The unresolved part of Stahl's target is now specifically
the constant-order remainder $-\log8+o(1)$.

% END retry_rational_rate.tex
% BEGIN retry_exp.tex
\section{Second pass: the sharp reciprocal-polynomial rate from finite Gram algebra}
\label{polyretry:exp}
\paragraph{Status.}
The type-$(0,n)$ part of \code{exp-rational-halfline} is now proved,
including its exact root limit $1/3$.  The diagonal type-$(n,n)$
Halphen equivalent is not proved by this argument.  The distinction is
between two conclusions inside one catalogue entry.

For a nonnegative integer $n$, write
\[
 \varepsilon_n=\inf_{\substack{p\in\R[x],\ \deg p\le n\\
                   p(x)\ne0\ (x\ge0)}}
       \sup_{x\ge0}|e^{-x}-p(x)^{-1}|.
\]
Allowing an arbitrary constant numerator gives the same infimum:
absorb a nonzero numerator into $p$; the zero numerator has error $1$,
whereas the constant approximant $1/2$ has error $1/2$.
The change of variable $x\mapsto x/4$ leaves this infimum unchanged.

\begin{theorem}[Explicit finite-degree bounds]
For every $n\ge0$, with $a_n=(4n+4)\log3+\log4$,
\begin{equation}\label{polyretry:reciprocal-sharp}
 \frac{1}{6\sqrt{a_n}}\le 3^n\varepsilon_n\le\sqrt2.
\end{equation}
Consequently $\lim_{n\to\infty}\varepsilon_n^{1/n}=1/3$.
\end{theorem}

The source of the rate is Sch\"onhage's
\href{https://doi.org/10.1016/0021-9045(73)90042-7}{1973 paper}.
The proof below reconstructs the method with a finite Gram computation
and a nilpotent polynomial operator.  For comparison, the proof in
\href{https://people.math.ethz.ch/~hiptmair/Seminars/RAP_22/PEP87.pdf}{
Petrushev--Popov, Section 4.5, Theorems 4.10--4.11}
uses a Laguerre expansion.  No completeness theorem for that expansion
is needed here.

\subsection{One finite Gram computation gives the exact squared error}
Let $f(x)=e^{x/4}$ and, only for the functions used below, set
\[
 \langle u,v\rangle_w=\int_0^\infty u(x)v(x)e^{-x}\,dx.
\]
All relevant integrals are finite: they are finite sums of polynomial
times exponential integrals.  Define the polynomials
\[
 L_k(x)=\sum_{j=0}^k(-1)^j\binom{k}{j}\frac{x^j}{j!}
       =\frac{e^x}{k!}\frac{d^k}{dx^k}(e^{-x}x^k).
\]
The second identity follows from the finite Leibniz rule.
Integration by parts $k$ times gives, for every polynomial $r$,
\[
 \langle L_k,r\rangle_w
 =\frac{(-1)^k}{k!}\int_0^\infty x^ke^{-x}r^{(k)}(x)\,dx.
\]
Every boundary term vanishes: at zero, derivatives of $x^ke^{-x}$
of order less than $k$ vanish; at infinity, an exponential dominates
each fixed polynomial.  Therefore $L_k$ is orthogonal to every
polynomial of degree less than $k$.  Since its leading coefficient is
$(-1)^k/k!$, the same identity and
$\int_0^\infty x^ke^{-x}\,dx=k!$ give
\[
 \langle L_j,L_k\rangle_w=\delta_{jk}.
\]
Their nonzero, triangular leading coefficients show that
$L_0,\ldots,L_n$ are a basis of the finite-dimensional space
$\mathcal P_n$ of polynomials of degree at most $n$.

The same integration by parts with $r=f$ gives
\[
 c_k:=\langle f,L_k\rangle_w
 =\frac{(-1)^k}{4^kk!}
   \int_0^\infty x^ke^{-3x/4}\,dx
 =\frac{4(-1)^k}{3^{k+1}},\qquad \langle f,f\rangle_w=2.
\]
Here the moment identity
$\int_0^\infty x^ke^{-s x}\,dx=k!/s^{k+1}$, $s>0$,
follows by scaling the preceding factorial integral.
Put $g_n=\sum_{k=0}^nc_kL_k$.  If
$p=\sum_{k=0}^nb_kL_k$, expansion of one finite square gives
\begin{align}
 \int_0^\infty(f-p)^2e^{-x}\,dx
 &=2-\sum_{k=0}^nc_k^2+\sum_{k=0}^n(b_k-c_k)^2 \notag\\
 &=\frac{2}{3^{2n+2}}+\sum_{k=0}^n(b_k-c_k)^2.
 \label{polyretry:finite-gram}
\end{align}
Thus the exact minimum is $2/3^{2n+2}$ and is attained by $g_n$.
This is a finite orthogonal projection identity, not Parseval's
identity for an infinite system.

\subsection{A nilpotent resolvent turns the Gram bound into a weighted uniform bound}
Let $D$ be differentiation on $\mathcal P_n$.  Since $D^{n+1}=0$,
\[
 (I-D)^{-1}=I+D+\cdots+D^n.
\]
Define the polynomial
\[
 s_n=\frac34(I-D)^{-1}g_n.
\]
For any $g\in\mathcal P_n$, repeated integration by parts gives
\[
 ((I-D)^{-1}g)(x)=e^x\int_x^\infty g(t)e^{-t}\,dt.
\]
Also $f(x)=\frac34e^x\int_x^\infty f(t)e^{-t}\,dt$.
Cauchy--Schwarz and \eqref{polyretry:finite-gram} now imply, for $x\ge0$,
\begin{align*}
 |f(x)-s_n(x)|
 &\le\frac34e^x
   \left(\int_x^\infty e^{-t}\,dt\right)^{1/2}
   \left(\int_x^\infty(f(t)-g_n(t))^2e^{-t}\,dt\right)^{1/2}\\
 &\le\frac{\sqrt2}{4\,3^n}e^{x/2}.
\end{align*}
In particular, for every $a>0$ the weighted best error satisfies
\begin{equation}\label{polyretry:weighted-error}
 \delta_{n,a}:=\min_{q\in\mathcal P_n}
    \max_{0\le x\le a} e^{-x/2}|f(x)-q(x)|
 \le\frac{\sqrt2}{4\,3^n}.
\end{equation}

\subsection{Weighted minimax gives a denominator with controlled tails}
We need the following elementary fact about a minimizer $q=q_{n,a}$
in \eqref{polyretry:weighted-error}:
\begin{equation}\label{polyretry:minimax-tail}
 0<q(x)\le f(x),\qquad q'(x)\ge0\qquad(x\ge a).
\end{equation}
For $n=0$ the derivative is zero, which is allowed.

Here are all ingredients.  Existence of a minimizer follows from the
same finite sample-value compactness proof as for ordinary minimax:
on $[0,a]$ the weight $e^{-x/2}$ is continuous and bounded below by
$e^{-a/2}>0$, so any minimizing sequence has bounded values at $n+1$
fixed, distinct nodes.  The inverse Vandermonde map gives bounded
coefficients, and a convergent subsequence gives a minimizer.

The weighted residual $e^{-x/2}(f-q)$ has $n+2$ alternating extreme
points.  To see that the usual proof still applies, suppose there were
at most $n+1$ alternating sign blocks in its nonzero extremal set.
Place the roots of a polynomial $h$ in the gaps between these blocks
and choose its sign to agree with $f-q$ on that set.  It has degree at
most $n$.  Compactness and the strictly positive weight imply that
$q+\eta h$ has smaller weighted error for sufficiently small $\eta>0$,
a contradiction.  The extremal sets of opposite signs have positive
distance, so the blocks and separating gaps in this argument are
finite.  The error is nonzero because $f$ is not a polynomial of
degree at most $n$.

Between consecutive alternating points, $f-q$ has a zero.  Hence $q$
interpolates $f$ at $n+1$ distinct nodes
$0<t_0<\cdots<t_n<a$.  Its Newton form is
\[
 q(x)=\sum_{j=0}^n f[t_0,\ldots,t_j]
                 \prod_{\ell=0}^{j-1}(x-t_\ell).
\]
Every divided difference is positive.  Indeed, repeated Rolle's
theorem applied to $f$ minus its degree-$j$ interpolant gives
$f[t_0,\ldots,t_j]=f^{(j)}(\xi_j)/j!>0$ for $j\ge1$;
for $j=0$ it is $f(t_0)>0$.  When $x\ge a$, all factors in the
Newton formula are positive, and differentiating each product shows
$q'(x)\ge0$.

Finally, for $x>t_n$, apply repeated Rolle's theorem to
\[
 f(t)-q(t)-c\prod_{j=0}^n(t-t_j),\qquad
 c=\frac{f(x)-q(x)}{\prod_{j=0}^n(x-t_j)}.
\]
This function has $n+2$ distinct zeros, so
$c=f^{(n+1)}(\xi)/(n+1)!>0$ for some $\xi\in(t_0,x)$.
Thus $q(x)<f(x)$ for $x\ge a$.  This proves
\eqref{polyretry:minimax-tail}, with no assumption that the interpolation
nodes or minimizer are explicitly computable.

\subsection{Upper bound: positivity on one interval and monotonicity beyond it}
Choose $a=4n\log3+\log4>0$, and take $q=q_{n,a}$.
By \eqref{polyretry:weighted-error}, for $0\le x\le a$,
\[
 q(x)\ge f(x)-\delta_{n,a}e^{x/2}
 = f(x)(1-\delta_{n,a}e^{x/4})\ge\tfrac12f(x),
\]
because $e^{a/4}=\sqrt2\,3^n$.
Thus $q$ has no zero on $[0,a]$, and
\[
 |e^{-x/4}-q(x)^{-1}|
 =\frac{|f(x)-q(x)|}{f(x)q(x)}
 \le2\delta_{n,a}\le\frac{\sqrt2}{2\,3^n}.
\]
For $x\ge a$, \eqref{polyretry:minimax-tail} gives
$q(a)\le q(x)\le f(x)$, and therefore
\[
 0\le q(x)^{-1}-e^{-x/4}
 \le q(a)^{-1}\le2e^{-a/4}=\frac{\sqrt2}{3^n}.
\]
The polynomial $q$ is consequently positive on the entire half-line
and proves the upper bound in \eqref{polyretry:reciprocal-sharp}.

\subsection{Lower bound: a small reciprocal error would contradict the finite Gram minimum}
Put $a=a_n=(4n+4)\log3+\log4$ and
$b_n=1/(2\sqrt a\,3^{n+1})$.
Suppose an admissible reciprocal polynomial $1/p$ has uniform error
$\epsilon<b_n$ when approximating $e^{-x/4}$.
Since $b_n<1$, evaluation at zero and the absence of zeros imply
$p(x)>0$ on the whole half-line.  Moreover $a>2$ and
\[
 \epsilon e^{a/4}<\frac{\sqrt2}{2\sqrt a}<\frac12.
\]
For $0\le x\le a$, the reciprocal error bound therefore gives
\[
 p(x)\le\frac{f(x)}{1-\epsilon f(x)}\le2f(x),
 \qquad |p(x)-f(x)|\le\epsilon f(x)p(x)
                         \le2\epsilon e^{x/2}.
\]
Choose the weighted minimizer $q=q_{n,a}$ again.  Its optimality
implies $|q(x)-f(x)|\le2\epsilon e^{x/2}$ on $[0,a]$.
Beyond $a$, \eqref{polyretry:minimax-tail} gives
$0<q(x)\le f(x)$, so $|q(x)-f(x)|\le f(x)$.
The exact Gram lower bound \eqref{polyretry:finite-gram} yields
\begin{align*}
 \frac{2}{3^{2n+2}}
 &\le \int_0^\infty(q(x)-f(x))^2e^{-x}\,dx\\
 &\le 4a\epsilon^2+\int_a^\infty e^{-x/2}\,dx
  =4a\epsilon^2+\frac{1}{3^{2n+2}}.
\end{align*}
Thus $\epsilon\ge b_n$, a contradiction.  Every admissible reciprocal
has error at least $b_n$; no existence theorem for a best rational
approximant on the noncompact half-line was used.  Taking the infimum
proves the lower bound in \eqref{polyretry:reciprocal-sharp}.
For $n\to\infty$, both $(6\sqrt{a_n})^{1/n}$ and
$(\sqrt2)^{1/n}$ tend to $1$, proving the claimed root limit.
The bounds also hold at $n=0$; separately $\varepsilon_0=1/2$.

\subsection{Operator corollary and inspected interfaces}
For every $n\ge0$, the upper-bound construction gives a real polynomial
$q_n$, positive on $[0,\infty)$, such that for every finite-dimensional
PSD operator $A$,
\[
 \|e^{-A}-q_n(4A)^{-1}\|\le\sqrt2\,3^{-n}.
\]
The factor $4$ reverses the scalar normalization $e^{-x/4}$.
Spectral decomposition proves the estimate; it is independent of the
largest eigenvalue and remains valid in dimension zero.

The declaration index was searched before choosing the proof.
\code{NLAlib.exists_sub_eval_interpolate_eq} already supplies the
Rolle interpolation remainder.  Mathlib's
\code{Lagrange.eval_interpolate_at_node},
\code{Lagrange.eq_interpolate_of_eval_eq}, and
\code{Polynomial.iterate_derivative_eq_zero} support the polynomial
steps.  For the moments, the inspected declarations
\code{Real.integral_rpow_mul_exp_neg_mul_Ioi} and
\code{Real.Gamma_nat_eq_factorial} in
\code{Analysis/SpecialFunctions/Gamma/Basic.lean} give the exact
factorial integrals, while
\code{integral_exp_neg_Ioi} and \code{integral_exp_mul_Ioi} give the
exponential tails.  A dedicated Laguerre interface was not found in
the inspected files; it is unnecessary to assume one exists.
The remaining mathematical helper package is finite: orthogonality
of the displayed $L_k$, weighted polynomial minimax, positivity of
Newton coefficients, and the finite differentiation resolvent.
These are proposed formalization interfaces, not checked Lean code.

\subsection{A finite Loewner-matrix certificate for the diagonal problem}
Although it does not establish the Halphen equivalent, the following
bound turns any choice of finitely many sample points into a rigorous
lower certificate.  Let $n\ge1$, $N=n+1$, and choose nonnegative real
numbers $x_1,\ldots,x_N,y_1,\ldots,y_N$ with $x_i\ne y_j$ for every
$i,j$.  Set
\[
 L_{ij}=\frac{e^{-x_i}-e^{-y_j}}{x_i-y_j},\qquad
 C=\left(\sum_{i,j}|x_i-y_j|^{-2}\right)^{1/2}.
\]
Then, with $\sigma_{\min}$ denoting the smallest Euclidean singular
value and $\|\cdot\|_2$ the spectral norm,
\begin{equation}\label{polyretry:loewner-certificate}
 E_{n,n}(e^{-x};[0,\infty))
 \ge\frac{\sigma_{\min}(L)}{2C}
 \ge\frac{|\det L|}{2C\|L\|_2^n}.
\end{equation}
Here $C>0$ and $\|L\|_2>0$: the exponential is strictly decreasing,
so every displayed divided difference is nonzero.

To prove the bound, take any admissible $r=p/q$ of type $(n,n)$,
with uniform error $\epsilon$, and define
$L^r_{ij}=(r(x_i)-r(y_j))/(x_i-y_j)$.
The polynomial
\[
 B(X,Y)=\frac{p(X)q(Y)-p(Y)q(X)}{X-Y}
\]
has degree at most $n-1$ in each variable.  This follows termwise:
for $i>j$, the quotient of $X^iY^j-X^jY^i$ by $X-Y$ is
$X^jY^j\sum_{k=0}^{i-j-1}X^{i-j-1-k}Y^k$, whose two degrees
are at most $i-1$.  Consequently, if $B_0$ is its $n\times n$
coefficient matrix and $V_x,V_y$ are the $N\times n$ monomial
evaluation matrices,
\[
 L^r=\operatorname{diag}(q(x_i)^{-1})\,V_xB_0V_y^\top
           \operatorname{diag}(q(y_j)^{-1}),
 \qquad \operatorname{rank}L^r\le n.
\]
Choose a Euclidean unit vector $v$ in its kernel.  Entrywise,
$|L_{ij}-L^r_{ij}|\le2\epsilon/|x_i-y_j|$, and hence
\[
 \sigma_{\min}(L)\le\|Lv\|_2
 =\|(L-L^r)v\|_2\le\|L-L^r\|_F\le2\epsilon C.
\]
Take the infimum over $r$.  Finally,
$|\det L|\le\sigma_{\min}(L)\|L\|_2^n$ proves the determinant
version.  For $n=0$ the same proof uses $1\times1$ matrices and
$L^r=0$ for constant $r$, yielding
$E_{0,0}\ge|e^{-x_1}-e^{-y_1}|/2$; sending $x_1=0$ and
$y_1\to\infty$ recovers the exact constant-approximation error $1/2$.

This certificate uses polynomial division, matrix rank, a kernel vector,
and singular values.  It is valid for every sample choice, including
choices for which the lower bound is zero.  To obtain the sharp Halphen
equivalent from such certificates would additionally require a family
of sample sets with matching sharp lower asymptotics and matching
upper approximants; neither is established here.

\subsection{Why this does not yet prove the diagonal Halphen prefactor}
The proof closes the $1/3$ claim completely, but it does not show
$E_{n,n}(e^{-x};[0,\infty))\sim2H^{n+1/2}$.
The denominator-only class is a proper restriction of the diagonal
class.  The above projection minimizes an additive weighted
polynomial error, whereas a diagonal rational error contains a
variable numerator and division by a variable denominator.  No
identity identifies that uniform error with the smallest singular
value of the finite Gram matrix above.

This is a real obstruction to the proposed shortcut: Sch\"onhage's
\href{https://doi.org/10.1137/0719077}{1982 paper} explicitly distinguishes
the $1/9$ rate in a related Laguerre-space problem from the original
rational uniform-approximation problem.  A finite matrix encoding can
give certificates for a fixed degree, but its large-degree singular
value asymptotics and the exact transfer back to the minimax error
still need proof.  In particular, root asymptotics alone cannot supply
the multiplicative constant $2\sqrt H$.
No new proof of that strong asymptotic, or a simpler substitute implying
it, has been obtained in this second pass; Aptekarev's strong-asymptotic
theorem remains the outstanding mathematical input for this subclaim.

% END retry_exp.tex

% BEGIN THIRD PASS

% BEGIN third_stahl.tex
\section{Third pass: an endpoint weight raises the rational lower prefactor to four}
\label{thirdstahl:main}
This argument strengthens the elementary kernel proof in
Section~\ref{retryroot:sharp-rate}.

Write $\mathcal E_n=E_{n,n}(|x|;[-1,1])$ and set
\[
 d_n=2\lfloor n/2\rfloor+1,
 \qquad S_n=\sqrt{\pi^2d_n+3+(\log2)^2}.
\]
We prove the explicit bound
\begin{equation}\label{thirdstahl:four-finite}
 \boxed{\mathcal E_n\ge
 \frac{4e^{-S_n}}{1+4e^{-2S_n}}\quad(n\ge0),
 \qquad
 \liminf_{n\to\infty}e^{\pi\sqrt n}\mathcal E_n\ge4.}
\end{equation}
This improves the previous elementary lower prefactor from one to
four.  It does not establish the desired limiting prefactor eight.
All steps use finite polynomial algebra, alternation, and one scalar
integral identity; no potential-theoretic existence theorem is used.

\subsection{The optimal rational function has a finite positive-node representation}
\label{thirdstahl:positive-nodes}
Let $m=\lfloor n/2\rfloor$.  Evenness and uniqueness, already proved in
Section~\ref{polybest:rational}, identify the best approximant as
$r(x)=p(x^2)/q(x^2)$, where $p/q$ is the reduced best type-$(m,m)$
approximant to $\sqrt t$ on $[0,1]$ and $q>0$ there.
Its error is positive, and $p\ne0$: the constant $1/2$ has error
$1/2$, whereas the zero function has error one.

We first pay the possible degree defect.  Put
$\delta=\min(m-\deg p,m-\deg q)$.  Rational alternation supplies
$2m+2-\delta$ alternating extremal points.  Between consecutive points
there is an interpolation point.  After taking positive square roots,
the polynomial
\[
 B(x)=p(x^2)-xq(x^2)
\]
therefore has at least $2m+1-\delta$ distinct roots in $(0,1)$.
But $\deg B\le2(m-\delta)+1$, so root counting forces $\delta=0$.
It follows that $B$ has degree exactly $2m+1$ and has exactly that many
distinct roots, all in $(0,1)$.  Thus, after a positive scalar
normalization,
\begin{equation}\label{thirdstahl:positive-product}
 P(x)=p(x^2)+xq(x^2)=\prod_{j=1}^{2m+1}(x+a_j),
 \qquad 0<a_1<\cdots<a_{2m+1}<1,
\end{equation}
and
\begin{equation}\label{thirdstahl:product-rational}
 r(x)=x\frac{P(x)+P(-x)}{P(x)-P(-x)}.
\end{equation}
The apparent zero-over-zero at $x=0$ is removable, with
\[
 r(0)=\frac{P(0)}{P'(0)}
     =\left(\sum_{j=1}^{2m+1}a_j^{-1}\right)^{-1}.
\]
The positive scalar assertion follows from $q(0)>0$: a polynomial
whose roots are all negative has coefficients of one sign, and the
coefficient of $x$ in $P$ is $q(0)$.

Conversely, every product in \eqref{thirdstahl:positive-product}
gives an admissible even type-$(2m,2m)$ rational function by
\eqref{thirdstahl:product-rational}.  Its denominator after cancelling
$x$ has strictly positive coefficients as a polynomial in $x^2$.
This is an exact finite-node minimization problem, and the minimum is
attained by the nodes above.  In particular, this reduction has no
unproved assertion about where complex zeros or poles accumulate.

\subsection{An exact integral identity at the right endpoint}
Retain the kernel $\kappa(s)=\log\coth(|s|/2)$ and its already-proved
mass $\int_{\R}\kappa=\pi^2/2$ from
\eqref{retryroot:kernel-mass}.  For $0<a\le1$ we claim
\begin{equation}\label{thirdstahl:weighted-kernel}
 \int_0^1
   \log\left|\frac{x+a}{x-a}\right|
   \frac{dx}{x\sqrt{1-x^2}}=\frac{\pi^2}{2}.
\end{equation}
This identity has a direct substitution proof.  Write
$a=1/\cosh v$, $v\ge0$, and substitute $x=1/\cosh u$.
The differential becomes $-du$, and the hyperbolic addition formulas
give, away from the logarithmic singularity,
\[
 \log\left|\frac{x+a}{x-a}\right|
 =\log\left|\frac{\cosh u+\cosh v}{\cosh u-\cosh v}\right|
 =\kappa(u+v)+\kappa(u-v).
\]
Evenness of $\kappa$ now yields
\[
 \int_0^\infty\bigl(\kappa(u+v)+\kappa(u-v)\bigr)\,du
 =\int_v^\infty\kappa+\int_{-v}^\infty\kappa
 =2\int_0^\infty\kappa=\frac{\pi^2}{2}.
\]
All terms are nonnegative.  The singularities are integrable, including
$v=0$; the substitution may first be performed on truncated intervals
and then extended by monotone convergence.  Thus no differentiation
through a principal-value integral is needed.

For the optimal product, put
\[
 U(x)=\log\left|\frac{P(x)}{P(-x)}\right|
     =\log\frac{r(x)+x}{|r(x)-x|}\qquad(0<x<1).
\]
Values $+\infty$ at interpolation nodes cause only integrable
logarithmic singularities.  Finite summation of
\eqref{thirdstahl:weighted-kernel} gives the exact identity
\begin{equation}\label{thirdstahl:mass-identity}
 \int_0^1 U(x)\frac{dx}{x\sqrt{1-x^2}}
 =\frac{d_n\pi^2}{2}.
\end{equation}
Moreover $U\ge0$, since $P(x)>|P(-x)|$ for $x>0$.

\subsection{The finite bound and its asymptotic consequence}
Let $\varepsilon=\mathcal E_n\in(0,1/2]$.
On $[\varepsilon,1]$, the uniform error inequality gives
\[
 U(x)\ge\log\frac{2x-\varepsilon}{\varepsilon}.
\]
Consequently the left side of \eqref{thirdstahl:mass-identity} is
at least
\[
 F(\varepsilon)=\int_\varepsilon^1
   \log\frac{2x-\varepsilon}{\varepsilon}
   \frac{dx}{x\sqrt{1-x^2}}.
\]
Define $A\ge0$ by $\cosh A=1/\varepsilon$; explicitly,
$A=\log((1+\sqrt{1-\varepsilon^2})/\varepsilon)$.
Using $\log(1-z)\ge-2z$ for $0\le z\le1/2$ gives
\begin{align*}
 F(\varepsilon)&\ge
 G(\varepsilon)-\varepsilon
       \int_\varepsilon^1\frac{dx}{x^2\sqrt{1-x^2}}
 \ge G(\varepsilon)-1,\\
 G(\varepsilon)&=
 \int_\varepsilon^1\log\frac{2x}{\varepsilon}
       \frac{dx}{x\sqrt{1-x^2}}
 =A\log(2\cosh A)-\int_0^A\log\cosh u\,du.
\end{align*}
Here the subtracted integral, including its factor $\varepsilon$,
equals $\sqrt{1-\varepsilon^2}\le1$.
The elementary identity
$\log\cosh u=u-\log2+\log(1+e^{-2u})$ shows that
\[
 G(\varepsilon)=\frac{A^2}{2}+A\log2
 +A\log(1+e^{-2A})-\int_0^A\log(1+e^{-2u})\,du.
\]
The last integral lies between zero and
$\int_0^\infty e^{-2u}\,du=1/2$; the preceding extra term is
nonnegative.  Hence
\[
 F(\varepsilon)\ge
 \frac12(A+\log2)^2-\frac12(\log2)^2-\frac32.
\]
Combining this with \eqref{thirdstahl:mass-identity} yields
$A+\log2\le S_n$.  Since $\cosh$ is increasing on the nonnegative
half-line and $S_n-\log2>0$,
\[
 \varepsilon=\frac1{\cosh A}
 \ge\frac1{\cosh(S_n-\log2)}
 =\frac{4e^{-S_n}}{1+4e^{-2S_n}}.
\]
Finally $d_n-n\in\{0,1\}$, so $S_n-\pi\sqrt n\to0$.
This proves both assertions in \eqref{thirdstahl:four-finite},
including the odd-degree sequence.  Degree zero is covered by the
same proof; its exact error remains $1/2$.

\subsection{A precise scalar statement that still separates four from eight}
For the actual optimal approximant $r_n$, write
$L_n=\log(1/\mathcal E_n)$ and define its nonnegative logarithmic slack
\begin{align}
 D_n={}&\int_0^{\mathcal E_n}U_n(x)\frac{dx}{x\sqrt{1-x^2}}
 \nonumber\\
 &+\int_{\mathcal E_n}^1
 \left[U_n(x)-\log\frac{2x-\mathcal E_n}{\mathcal E_n}\right]
 \frac{dx}{x\sqrt{1-x^2}}.
 \label{thirdstahl:slack}
\end{align}
Every integrand is nonnegative and the integrals are finite.
The explicit calculation above also gives, as $\varepsilon\downarrow0$,
\[
 F(\varepsilon)=\frac12\log^2\frac4\varepsilon+O(1).
\]
For clarity, this uses only
$A=\log(2/\varepsilon)+O(\varepsilon^2)$,
$0\le A\log(1+e^{-2A})\le A e^{-2A}$, the bound $1/2$ on the
last integral, and $G(\varepsilon)-1\le F(\varepsilon)\le G(\varepsilon)$.
Thus the exact identity \eqref{thirdstahl:mass-identity} implies
\begin{equation}\label{thirdstahl:slack-equivalence}
 \pi^2n=(L_n+\log4)^2+2D_n+O(1).
\end{equation}
Since the sharp root exponent already gives
$L_n/(\pi\sqrt n)\to1$, elementary factorization of the difference
of squares in \eqref{thirdstahl:slack-equivalence} proves, for every
finite $c\ge0$,
\[
 \frac{D_n}{L_n}\longrightarrow c
 \quad\Longleftrightarrow\quad
 e^{\pi\sqrt n}\mathcal E_n\longrightarrow4e^c.
\]
In particular the remaining prefactor-eight target is exactly the
scalar assertion $D_n/L_n\to\log2$ for these finite optimal products.
This last assertion is \emph{not proved here}.  Alternation by itself
does not supply the average logarithmic slack between extremal
points, so it would be incorrect to replace that missing estimate by
an informal sinusoidal-ripple heuristic.

\paragraph{Formalization boundary.}
The new paid ingredients are degree-defect elimination by root count,
the finite positive-node representation, the hyperbolic substitution,
and the explicit lower bound \eqref{thirdstahl:four-finite}.
The existing rational alternation argument and scalar kernel-mass
proof are reused unchanged.  One may define $A$ by its logarithmic
formula, avoiding a separate inverse-hyperbolic-cosine API; Mathlib's
\code{Real.cosh_log} and hyperbolic differentiation and addition
lemmas provide the elementary identities.  The slack limit remains
a new mathematical obligation, not merely a missing Lean interface.
% END third_stahl.tex

% BEGIN third_bridges.tex
\section{Exact finite dual certificates for rational uniform approximation}
\label{polythird:dual}
The Loewner singular-value estimate gives a useful lower bound, but it
introduces norm inequalities before the sharp asymptotic question is
reached.  The following finite certificate preserves the minimax
constant exactly.  Its proof uses only polynomial evaluation and finite
sums; its completeness uses the rational alternation theorem already
proved in Section~\ref{polybest:rational}.

\subsection{A sufficient certificate with no denominator supplied}
Let $I=[a,b]$ with $a<b$, let $f\in C(I)$, and let $m,n\ge0$.
Write $E_{m,n}(f;I)$ for the best uniform error among real fractions
$p/q$ of degrees at most $(m,n)$ with no pole on $I$.  Cancel common
factors and orient the denominator so that $q>0$ on $I$.
Choose finitely many points $t_1,\ldots,t_K\in I$, signs
$s_i\in\{-1,1\}$, and nonnegative weights $a_i$, not all zero.
Suppose that, for a real number $\eta\ge0$,
\begin{align}
 \sum_{i=1}^K a_i s_i t_i^j&=0 &&(0\le j\le m),
       \label{polythird:dual-numerator}\\
 \sum_{i=1}^K a_i\bigl(s_i f(t_i)-\eta\bigr)t_i^j&=0
       &&(0\le j\le n).\label{polythird:dual-denominator}
\end{align}
Then
\begin{equation}\label{polythird:dual-lower}
 E_{m,n}(f;I)\ge\eta.
\end{equation}

Indeed, for every admissible $p/q$, the two moment identities give
\[
 \sum_i a_i s_i p(t_i)=0,\qquad
 \sum_i a_i s_i f(t_i)q(t_i)=\eta\sum_i a_iq(t_i).
\]
Subtract the first equality from the second and use $q(t_i)>0$:
\[
 \eta\sum_i a_iq(t_i)
 =\sum_i a_i s_i q(t_i)\bigl(f(t_i)-p(t_i)/q(t_i)\bigr)
 \le\|f-p/q\|_\infty\sum_i a_iq(t_i).
\]
The final sum is strictly positive, so cancellation proves the claim.
For this sufficient direction the sample points need not be distinct
or ordered, the signs need not alternate, and no full-degree
assumption is present.  The same argument applies directly on the
half-line whenever the denominator is positive there.

Let $V_j$ be the $K\times(j+1)$ matrix with entries
$(V_j)_{ik}=t_i^k$ for $0\le k\le j$, and set
$F=\diag(f(t_i))$ and $S=\diag(s_i)$.  The certificate is the finite
matrix system
\begin{equation}\label{polythird:dual-pencil}
 \begin{bmatrix}V_m^\top S\\ V_n^\top(FS-\eta I)\end{bmatrix}a=0,
 \qquad a\ge0,\qquad \sum_i a_i=1.
\end{equation}
The normalization loses nothing.  For fixed nodes, signs, and $\eta$,
this is linear feasibility in the vector $a$.  Checking this system
needs no reconstruction of a competitor's denominator.

\subsection{Completeness, including degree defects and the zero function}
Let $r_*=p/q$ be a best approximant in reduced form with $q>0$.
If $p\ne0$, write $\mu=\deg p$, $\nu=\deg q$, and
\[
 N=\max(m+\nu,n+\mu),\qquad K=N+2.
\]
If $p=0$, use $p/q=0/1$, set $N=m$, and again put $K=N+2$.
The already proved alternation theorem supplies distinct increasing
nodes $t_1<\cdots<t_K$ and alternating signs $s_i$ such that
\[
 f(t_i)-r_*(t_i)=s_i E,
 \qquad E=E_{m,n}(f;I).
\]
When $E=0$, any $K$ distinct nodes and alternating signs suffice.
Let
\[
 \omega(X)=\prod_{i=1}^K(X-t_i),\qquad
 b_i=\frac1{|\omega'(t_i)|},\qquad
 \widetilde a_i=b_iq(t_i)>0.
\]
For every polynomial $h$ of degree at most $K-2=N$,
\begin{equation}\label{polythird:dual-annihilator}
 \sum_i s_i b_i h(t_i)=0.
\end{equation}
To see this, write the Lagrange interpolation formula for $h$ and
compare the coefficient of $X^{K-1}$.  It gives
$\sum_i h(t_i)/\omega'(t_i)=0$; the alternating sign pattern of
$\omega'(t_i)$ differs from $s_i$ only by a common sign.

Apply \eqref{polythird:dual-annihilator} first to $h=qX^j$ for
$j\le m$.  Its degree is at most $m+\nu\le N$; for $p=0$ it is
at most $m=N$.  This proves \eqref{polythird:dual-numerator} with
$a_i=\widetilde a_i$.  At each node, the extremal error equality gives
\[
 \widetilde a_i\bigl(s_i f(t_i)-E\bigr)=s_i b_i p(t_i).
\]
Apply the same annihilator to $h=pX^j$ for $j\le n$.
Its degree is at most $n+\mu\le N$ when $p\ne0$, and it is the
zero polynomial when $p=0$.  This proves
\eqref{polythird:dual-denominator}.  Divide every weight by
$\sum_i\widetilde a_i$ to obtain \eqref{polythird:dual-pencil}.

Consequently the exact best error is the maximum of all numbers
$\eta\ge0$ admitting the finite certificate
\eqref{polythird:dual-pencil}, with $2\le K\le m+n+2$.
The maximum is attained; this follows from the explicit certificate
at a best approximant, without a compactness assertion about the
space of node configurations.  Restricting to distinct increasing
nodes and alternating signs still gives the same maximum.
All degree defects are already accounted for by $K=N+2$.

For example, for $|x|$ on $[-1,1]$ and type $(1,1)$, take
\[
 (t_1,t_2,t_3)=(-1,0,1),\quad
 (s_1,s_2,s_3)=(1,-1,1),\quad
 (a_1,a_2,a_3)=(1/4,1/2,1/4),\quad \eta=1/2.
\]
The four moment equalities are immediate.  They certify
$E_{1,1}\ge1/2$, while the constant $r=1/2$ gives equality.
This is a defective best approximant: forcing a full square system
with four distinct alternating extremal nodes would be incorrect.

\subsection{Rigorous sample enclosures and the remaining asymptotic task}
Exact function values are not essential for a certified finite bound.
Suppose real sample data $c_i$ satisfy $|f(t_i)-c_i|\le\delta$.
If \eqref{polythird:dual-numerator} holds and
\eqref{polythird:dual-denominator} holds with $c_i$ in place of
$f(t_i)$, then the identical proof, with an additional term bounded
by $\delta\sum_i a_iq(t_i)$, gives
\begin{equation}\label{polythird:dual-enclosures}
 E_{m,n}(f;I)\ge\eta-\delta.
\end{equation}
Thus rational nodes, rational sample enclosures, and exact finite
linear algebra can produce rigorous lower bounds for exponential
approximation.  An approximate null vector alone does not suffice:
the exact moment equalities are the step that removes all unknown
numerator and denominator coefficients.

This is an exact finite-dimensional reformulation, not a proof of
either remaining strong asymptotic.  To prove Stahl's prefactor by
this route one still needs explicit certificates with
$\eta_n=(8+o(1))e^{-\pi\sqrt n}$ and matching upper approximants.
For the half-line exponential one needs
$\eta_n=(2\sqrt H+o(1))H^n$ and matching upper approximants; any
sample-enclosure error must be $o(H^n)$.  Existence of the maximizing
nodes through an unknown minimizer is not an explicit asymptotic
construction.

The direct certificate proof uses finite sums and monomial evaluation.
Its completeness reuses the existing proof of rational alternation
and the leading-coefficient argument for Lagrange interpolation.
The inspected Mathlib declarations
\code{Lagrange.eq_interpolate} and
\code{Lagrange.leadingCoeff_basis} supply the interpolation identities.
The current NLAlib declaration index contains no public declaration
whose name includes \code{alternation} or \code{equioscillation};
the reuse claimed here is mathematical reuse within this document,
not a claim that this entire certificate has already been formalized.
% END third_bridges.tex

% BEGIN third_halphen.tex
\section{A symmetric finite operator for alternating rational error}
\label{polythird:spectral}
There is a more precise operator replacement for the Loewner lower
bound: orthogonal compression of a signed diagonal matrix.  It
preserves the alternating error exactly.  The symmetric eigenproblem
below is the one used in the rational Remez method; see
\href{https://people.maths.ox.ac.uk/trefethen/minimax_published.pdf}{
Filip--Nakatsukasa--Trefethen--Beckermann (2018), Section 4.1,
equations (4.6)--(4.9)}.  We give the finite proof and an explicit
residual certificate so that a future formalization need not import
an algorithm-convergence theorem.

\subsection{Barycentric cancellation produces the complementary subspace}
Fix $n\ge0$, $2n+2$ distinct increasing real nodes
$t_0<\cdots<t_{2n+1}$ in an interval $I$, and either sign
$\sigma\in\{-1,1\}$.  Put
\[
 s_i=\sigma(-1)^i,\quad S=\diag(s_i),\quad
 a_i=\frac1{\prod_{j\ne i}|t_i-t_j|},\quad W=\diag(a_i),
 \qquad V_{ij}=t_i^j\quad(0\le j\le n).
\]
For $f:I\to\R$, let $F=\diag(f(t_i))$.
Comparison of the leading coefficient in Lagrange interpolation gives
\[
 \sum_{i=0}^{2n+1}s_i a_i h(t_i)=0
       \qquad(\deg h\le2n).
\]
Indeed, if $\omega(X)=\prod_i(X-t_i)$, then
$1/\omega'(t_i)$ differs from $s_i a_i$ by one common sign.
Applying this identity to all products of two monomials of degree at
most $n$ gives
\begin{equation}\label{polythird:isotropic-vandermonde}
 V^\top WSV=0.
\end{equation}

The matrix $V$ has full column rank: a nonzero polynomial of degree
at most $n$ cannot vanish at $2n+2$ distinct points.  Choose any
matrix $M$ whose columns are an orthonormal basis of
$\operatorname{range}(W^{1/2}V)$.  Thus $M$ has size
$(2n+2)\times(n+1)$.  An explicit choice is
$M=W^{1/2}VB^{-1/2}$, where $B=V^\top WV$ is positive definite;
Gram--Schmidt is an equally valid construction and avoids matrix
square roots.  Equation~\eqref{polythird:isotropic-vandermonde}
implies
\begin{equation}\label{polythird:orthogonal-split}
 M^\top M=I,\qquad M^\top SM=0,\qquad
 MM^\top+SMM^\top S=I.
\end{equation}
For the last equality, the $2n+2$ columns of $[M\;SM]$ are
orthonormal and form a basis of $\R^{2n+2}$.
Define the real symmetric operator
\begin{equation}\label{polythird:signed-compression}
 T=M^\top SFM\in\R^{(n+1)\times(n+1)}.
\end{equation}

\subsection{Every positive sample vector gives a rigorous lower bound}
Let $b\in\R^{n+1}$ satisfy $g:=Mb>0$ coordinatewise.  There are
unique polynomials $q,p$ of degree at most $n$ such that
\[
 W^{1/2}(q(t_i))_i=Mb,
 \qquad
 W^{1/2}(p(t_i))_i=MM^\top FMb.
\]
The existence follows from the definition of $M$; uniqueness is
polynomial root counting.  In particular $q(t_i)>0$, although no
positivity between sample points is asserted or required below.
The complementary-subspace identity gives
\[
 W^{1/2}(f(t_i)q(t_i)-p(t_i))_i
 =(I-MM^\top)FMb=SMTb.
\]
Consequently the sample residual is exactly
\begin{equation}\label{polythird:sample-residual}
 f(t_i)-\frac{p(t_i)}{q(t_i)}
 =s_i\frac{(MTb)_i}{(Mb)_i}.
\end{equation}
Write
\[
 \ell(b)=\min_i\frac{(MTb)_i}{(Mb)_i}.
\]
Then the best error among all pole-free type-$(n,n)$ fractions on
$I$ obeys
\begin{equation}\label{polythird:cone-certificate}
 E_{n,n}(f;I)\ge\max\{0,\ell(b)\}.
\end{equation}
The assertion holds on a compact interval or on the half-line,
provided the nodes belong to the approximation domain.

For a proof when $\ell(b)>0$, suppose a competing fraction $u/v$
has uniform error less than $\ell(b)$.  Cancel its common factors
and change its overall sign so that $v>0$ on $I$.
Equation~\eqref{polythird:sample-residual} implies
\[
 s_i\bigl(u(t_i)q(t_i)-p(t_i)v(t_i)\bigr)>0
          \quad(0\le i\le2n+1).
\]
The polynomial $uq-pv$ has degree at most $2n$ and alternating
strict signs at $2n+2$ increasing nodes.  The intermediate value
theorem gives at least $2n+1$ distinct roots, a contradiction.
If $\ell(b)\le0$, the claimed lower bound is simply nonnegativity.

\subsection{An eigenpair, or a certified approximate eigenpair}
If $Tb=\lambda b$ and $Mb>0$, then
\eqref{polythird:sample-residual} becomes exact equioscillation:
\[
 f(t_i)-p(t_i)/q(t_i)=s_i\lambda,
 \qquad E_{n,n}(f;I)\ge\max\{0,\lambda\}.
\]
The relevant eigenvalue is selected by the positivity of $Mb$;
it is not automatically the largest, smallest, or smallest-magnitude
eigenvalue.  For a fixed $S$, two distinct eigenvalues cannot both
have such vectors: eigenvectors of a symmetric matrix for distinct
eigenvalues are orthogonal, whereas
$b_1^\top b_2=(Mb_1)^\top(Mb_2)>0$.
This does not assert simplicity of a repeated eigenvalue.

An exact eigenvector is unnecessary for a lower certificate.
For any real $\lambda$, set
\[
 e=Tb-\lambda b,\qquad
 A(b)=\max_i\frac{\|M_{i,:}\|_2}{(Mb)_i}.
\]
The finite Cauchy--Schwarz inequality gives
\[
 \frac{(MTb)_i}{(Mb)_i}
 =\lambda+\frac{(Me)_i}{(Mb)_i}
 \ge\lambda-A(b)\|e\|_2.
\]
Thus
\begin{equation}\label{polythird:spectral-residual-certificate}
 E_{n,n}(f;I)
 \ge\max\{0,\lambda-A(b)\|Tb-\lambda b\|_2\}.
\end{equation}
All quantities on the right are finite-dimensional.
The denominator positivity margin in $A(b)$ is essential; a small
eigen-residual by itself is not enough.

The same formula tolerates certified sample errors.  If
$|c_i-f(t_i)|\le\delta$, form $T_c=M^\top S\diag(c_i)M$.
Reconstruct $p$ using $c_i$ instead of $f(t_i)$ in the preceding
projection.  Its actual signed residual differs by at most $\delta$
at each node.  Therefore
\begin{equation}\label{polythird:spectral-sample-certificate}
 E_{n,n}(f;I)\ge
 \max\{0,\lambda-A(b)\|T_cb-\lambda b\|_2-\delta\}.
\end{equation}
This is useful with rational enclosures of exponential values.

Conversely, suppose an admissible fraction $p_*/q_*$, with $q_*>0$,
has global uniform error $E$ and equal alternating residual $s_iE$
at these $2n+2$ nodes.
Write $W^{1/2}(q_*(t_i))_i=Mb$.  Multiplication of
$W^{1/2}(f q_*-p_*)=ESMb$ by $M^\top S$ gives $Tb=Eb$.
Thus at such a full alternation the matrix certificate attains the
exact error, with no Frobenius-norm or determinant loss.  If the
best fraction is defective and has fewer extremal nodes, use the
rectangular certificate in Section~\ref{polythird:dual}; one must
not assume a full alternation that is absent.

\subsection{The precise Halphen task left by this reduction}
For the half-line problem one may use finite $x$-nodes directly.
To include the endpoint at infinity, put $t=x/(1+x)$ and
\[
 \widehat f(t)=\exp\!\left(-\frac{t}{1-t}\right)\quad(0\le t<1),
 \qquad\widehat f(1)=0.
\]
This continuous compactification preserves the uniform error and
the type-$(n,n)$ rational class, after cancellation of removable
endpoint factors.  More explicitly, any bounded pole-free $p/q$
on the half-line has $\deg p\le\deg q=k\le n$; multiplication
by $(1-t)^k$ after substitution gives a denominator positive on
$[0,1]$, including its positive leading coefficient at $t=1$.
The reverse substitution has no pole at finite $x\ge0$.

A genuinely sufficient new proof of the Halphen equivalent would
now consist of node sets, vectors $b_n$ with $M_nb_n>0$, and real
$\lambda_n$ satisfying
\[
 \lambda_n=(2\sqrt H+o(1))H^n,
 \qquad
 A(b_n)\|T_nb_n-\lambda_nb_n\|_2=o(H^n),
\]
together with admissible upper approximants of error
$(2\sqrt H+o(1))H^n$.  Sample enclosures must also have width
$o(H^n)$.  These conditions imply the result by the two-sided
squeeze and retain the original multiplicative constant.
No such asymptotic node and eigenvector construction has been
proved here.  The reduction identifies a finite operator target,
but does not replace that missing asymptotic estimate by a
numerical eigenvalue fit or by root convergence alone.

For formalization, the core objects are one polynomial evaluation
subspace, its orthonormal basis, the signed complementary subspace,
and a symmetric compression.  The inspected declaration
\code{NLAlib.exists_orthonormal_complement} provides the existing
orthogonal-subspace interface; the explicit complement $SM$ here
is even more specific.  The barycentric annihilator and rational
root-counting argument reuse the preceding derivations.  This is
a proposed mathematical interface, not a claim of a completed Lean
formalization or a completed Halphen asymptotic proof.
% END third_halphen.tex

% BEGIN third_finite.tex
\section{Removing the half-line and transcendental target by a quantitative finite reduction}
\label{polythird:finite}

This reduction preserves the strong asymptotic, not merely its root rate.
It does not evaluate that asymptotic.  Its advantage is that the finite
problems involve only polynomials on $[0,1]$, with an explicit error in
the reduction.  Write
\[
 F_n=E_{n,n}(e^{-x};[0,\infty)),\qquad
 A_{n,M}=E_{n,n}(t^M;[0,1])\quad(M\ge1).
\]
The elementary substitution and estimate below are Lemmas 2.1--2.2 of
Nakatsukasa--Trefethen,
\href{https://people.maths.ox.ac.uk/trefethen/ratxn.pdf}{Rational approximation of $x^n$}.
Their final asymptotic theorem uses Aptekarev's theorem; we use only the
elementary comparison, so that theorem is not being invoked circularly.

\begin{proposition}[Quantitative polynomial-target reduction]
\label{polythird:power-reduction}
For every $n\ge0$ and integer $M\ge1$,
\begin{equation}\label{polythird:power-error}
 |F_n-A_{n,M}|\le\frac1{eM}.
\end{equation}
Consequently, with $M_n=2^{n^2}$ and any fixed $0<H<1$,
\[
 \frac{F_n-A_{n,M_n}}{H^{n+1/2}}\longrightarrow0.
\]
In particular, the original Halphen equivalent holds if and only if
$A_{n,2^{n^2}}/(2H^{n+1/2})\to1$.
\end{proposition}
\begin{proof}
Put $g_M(x)=(1+x/M)^{-M}$.  The inequality
$\log(1+u)<u$ for $u>0$ gives $g_M(x)>e^{-x}$ for $x>0$.
Their difference is zero at $0$ and tends to zero at infinity, so its
positive maximum occurs at an interior point $a$.  Differentiation gives
$(1+a/M)^{-M-1}=e^{-a}$ there.  Thus
\[
 0\le g_M(x)-e^{-x}\le g_M(a)-e^{-a}
   =\frac{ae^{-a}}M\le\frac1{eM}.
\]
The last inequality follows by differentiating $xe^{-x}$.

The substitution $t=M/(M+x)$ is a bijection between $[0,\infty)$
and $(0,1]$ and sends $g_M$ to $t^M$.  It preserves type $(n,n)$:
after composition, multiplying numerator and denominator by $t^n$
(or, in the inverse direction, by $(M+x)^n$) produces polynomials
of degrees at most $n$.  A bounded rational function on $(0,1]$
has at most a removable singularity at zero: factoring its numerator
and denominator by their powers of $t$ proves this directly.  After
cancellation it therefore extends to an admissible rational function
on $[0,1]$.  Only bounded candidates matter for these minimax errors.
Continuity at the endpoint preserves the supremum.  Hence
$E_{n,n}(g_M;[0,\infty))=A_{n,M}$.

For any approximation class, its distance functional satisfies
$|E(f)-E(g)|\le\|f-g\|_\infty$, by the triangle inequality and
taking infima.  This proves \eqref{polythird:power-error}.
Finally $2^{-n^2}H^{-n-1/2}\to0$, since the negative quadratic term
dominates the linear term in its logarithm.
\end{proof}

The choice $2^{n^2}$ is deliberately conservative.  It proves an exact
reduction without assuming a numerical enclosure for $H$; it is not
a recommendation for numerical computation.  More generally, any
integer sequence with $M_nH^n\to\infty$ gives the same conclusion.

\subsection{A finite coefficient test for every strict polynomial inequality}
\label{polythird:bernstein-positive}
The next elementary lemma converts strict positivity into finite linear
inequalities.  It uses no sum-of-squares representation theorem.
For a real polynomial $h(t)=\sum_{k=0}^d a_kt^k$ and $N\ge d$, write
\[
 B_{j,N}(t)=\binom Nj t^j(1-t)^{N-j},\qquad
 b_{j,N}(h)=\sum_{k=0}^{\min(j,d)}a_k\frac{\binom jk}{\binom Nk}.
\]
Then
\begin{equation}\label{polythird:bernstein-expansion}
 h(t)=\sum_{j=0}^N b_{j,N}(h)B_{j,N}(t),\qquad
 \sum_{j=0}^N B_{j,N}(t)=1.
\end{equation}
Indeed,
$\binom Nj\binom jk/\binom Nk=\binom{N-k}{j-k}$;
factoring out $t^k$ and applying the binomial theorem proves the
identity for each monomial $t^k$, then for $h$ by linearity.

\begin{lemma}[Strict positivity has a finite Bernstein certificate]
\label{polythird:bernstein-certificate}
A real polynomial $h$ is strictly positive throughout $[0,1]$ if and
only if, at some degree $N\ge\deg h$, all $b_{j,N}(h)$ are positive.
The zero polynomial is excluded by either condition.  More explicitly,
for $N\ge\max(d,1)$,
\begin{equation}\label{polythird:bernstein-estimate}
 \left|b_{j,N}(h)-h(j/N)\right|
 \le\frac{C_h}{N-d+1},\qquad
 C_h=\frac12\sum_{k=0}^d|a_k|k(k-1).
\end{equation}
\end{lemma}
\begin{proof}
If all coefficients are positive, \eqref{polythird:bernstein-expansion}
expresses $h(t)$ as a convex combination of positive numbers.
For the converse, let $h\ge\eta>0$ on $[0,1]$, using compactness.
The coefficient estimate makes all coefficients positive for sufficiently
large $N$.

To prove the estimate, the terms $k=0,1$ agree exactly.  For $k\ge2$
and $j\ge k$, set $u=j/N$.  The product
\[
 \frac{\binom jk}{\binom Nk}
 =\prod_{i=0}^{k-1}\frac{j-i}{N-i}
\]
has factors in $[0,u]$.  The telescoping product identity gives
\[
 0\le u^k-\frac{\binom jk}{\binom Nk}
 \le\sum_{i=0}^{k-1}\left(u-\frac{j-i}{N-i}\right)
 \le\sum_{i=0}^{k-1}\frac{i}{N-k+1}
 =\frac{k(k-1)}{2(N-k+1)}.
\]
When $j<k$ the binomial ratio is zero, while
$(j/N)^k\le j/N\le(k-1)/N\le k(k-1)/(2(N-k+1))$.
Sum these inequalities with weights $|a_k|$, and use
$N-k+1\ge N-d+1$.
\end{proof}

\subsection{Exact finite linear feasibility formulations}
Fix $\epsilon>0$ and a polynomial target $f$ on $[0,1]$.  Choose a
nonnegative integer degree bound $d_f$ for $f$, allowing $d_f=0$ when
$f=0$.  For each $N\ge n+d_f$ (enlarging to $N\ge1$ if necessary), consider unknown
coefficients of real polynomials $p,q$ of degrees at most $n$.  Require
every degree-$N$ Bernstein coefficient of each of
\begin{equation}\label{polythird:lp-polynomials}
 q,\qquad \epsilon q+(p-fq),\qquad \epsilon q-(p-fq)
\end{equation}
to be at least $1$.  At fixed $\epsilon,N,f$ these are finitely many
\emph{linear} inequalities in the coefficients of $p,q$.

\begin{proposition}[An exact hierarchy for strict minimax bounds]
\label{polythird:lp-equivalence}
$E_{n,n}(f;[0,1])<\epsilon$ if and only if the displayed finite system
is feasible for some $N$.  For $f(x)=|x|$ on $[-1,1]$, the same
equivalence holds using, for each $\sigma\in\{-1,1\}$, the three
polynomials
\[
 q(\sigma t),\quad
 \epsilon q(\sigma t)+(p(\sigma t)-tq(\sigma t)),\quad
 \epsilon q(\sigma t)-(p(\sigma t)-tq(\sigma t))
\]
on $[0,1]$, with $N\ge n+1$.
\end{proposition}
\begin{proof}
A feasible system gives $q\ge1$ and
$\epsilon q\pm(p-fq)\ge1$.  If $Q=\max_{[0,1]}q$, then
$|p/q-f|\le\epsilon-1/Q<\epsilon$.
For the two-interval formulation use the maximum on $[-1,1]$.

Conversely, take an admissible rational function with uniform error
strictly below $\epsilon$.  Its real denominator has constant sign;
replace $(p,q)$ by $(-p,-q)$ if necessary.  The polynomials in
\eqref{polythird:lp-polynomials} are all strictly positive.
Lemma~\ref{polythird:bernstein-certificate} makes all their coefficients
positive for a common sufficiently large $N$; its estimate applies at
every sufficiently large degree.  Scale $(p,q)$ by one positive number
so that the smallest of the finitely many coefficients is at least $1$.
The same argument applies to the six polynomials for $|x|$.
\end{proof}

This gives complete finite certificates for every strict upper bound,
with only polynomial arithmetic, binomial coefficients, and linear
inequalities.  For rational $\epsilon$ and a rational-coefficient
target such as $t^M$, the system has rational coefficients; a real
feasible point can be scaled to make the inequalities strict and then
approximated by a rational point.  Thus the certificate can itself be
rational.  Failure at one chosen $N$ is \emph{not} a lower bound for the
whole hierarchy.  The signed certificates of
Section~\ref{polythird:dual} supply the complementary lower-bound route.

Combining Proposition~\ref{polythird:power-reduction} with this hierarchy
removes both the unbounded domain and the exponential from each finite
Halphen approximation problem, with a proved error negligible relative
to $H^n$.  What remains is a uniform construction and evaluation of
certificates as $n\to\infty$.  This last assertion is not supplied by
linear feasibility at each fixed degree.

\paragraph{Inspected formalization interfaces.}
Mathlib already defines \code{bernsteinPolynomial} in
\code{RingTheory/Polynomial/Bernstein.lean};
\code{bernstein_nonneg} and \code{bernstein.probability} in
\code{Analysis/SpecialFunctions/Bernstein.lean} give nonnegativity and
the partition of unity.  The explicit coefficient conversion and
estimate \eqref{polythird:bernstein-estimate} are the additional finite
lemmas proposed here.  Uniform approximation by Bernstein polynomials
alone is not the needed statement: these are exact coefficients of the
\emph{same} polynomial at successively larger degrees.
% END third_finite.tex

% END THIRD PASS

\section{Suggested mathematical interfaces for the next agent}
The highest-reuse units are small and have independent proof contracts:
\begin{enumerate}
\item A bounded linear projection reproducing $\mathcal P_n$ has error
at most $(1+\|L\|)E_n(f)$. Apply it to interpolation and truncation.
\item A finite Lagrange coefficient identity gives the signed annihilator
used in polynomial and rational alternation, including uniqueness.
\item Finite Bessel inequalities for difference quotients prove Lipschitz
coefficient summability. Step approximations plus coefficient recursion
give the BV decay bound, and one telescope supplies all algebraic tails.
\item Compressed multiplication in the finite polynomial Gram space
gives general Gauss nodes, positivity, and exactness. Multiplication
modulo a polynomial gives the colleague characteristic polynomial.
\item A positive, exact quadrature functional is controlled by uniform
approximation. Lebesgue-specific Chebyshev moments and finite alias
blocks provide the sharper constants and the explicit eventual threshold.
\item The Hermitian diagonal norm identity transfers every scalar
polynomial or rational bound to operators and vector actions.

\item Positive-residue rational identities give sum-of-squares
Wronskians. A finite interpolation annihilator then turns the
evaluation matrix into a column-sum-zero matrix; strict diagonal
dominance supplies its minors and its adjugate supplies a positive
normal vector. A proper difference map gives the global comparison.
\item Subtract an explicit Chebyshev comparison before taking
coefficient limits. Its scaled entire limit and a small contraction
convert a polynomial-error liminf into a full limit.
\item A finite Laguerre Gram projection and the resolvent
$(I-D)^{-1}=\sum_{j=0}^nD^j$ convert an exact weighted squared error
into the sharp reciprocal-polynomial root rate.
\item A clipped logarithmic kernel and finite Riemann sums construct
rational approximants with sharp root exponent. Loewner matrix rank
gives finite lower certificates for diagonal rational approximation.
\item A weighted scalar kernel and the exact positive-node representation
give the improved lower prefactor $4$; the unproved logarithmic-slack
limit is stated explicitly rather than hidden in an asymptotic step.
\item Signed Vandermonde moments give lossless rational lower certificates,
including degree defects.  Orthonormal compression turns full alternation
into a symmetric eigenproblem, with a proved residual-error bound.
\item A quantitative M\"obius substitution reduces the half-line
exponential to polynomial targets.  Exact Bernstein coefficients and a
finite coefficient estimate certify every strict rational upper bound.
\end{enumerate}
The mathematical proofs were cross-reviewed for constants, multiplicities,
support size, zero error, degree zero, and exceptional rational degrees.
The two remaining sharp rational boundaries in
Section~\ref{poly:scope} remain visible. The full Lebesgue bound,
Bernstein polynomial limit, and reciprocal root limit have mathematical
proofs in this document; no Lean implementation of them has been checked.
Neither remaining strong asymptotic is certified by a root-rate estimate
or by a numerical experiment.
\end{document}
